\chapter{Plasma physics for engineers}
\label{vol04:ch13}

\epigraph{The most important property of a plasma is its tendency to
remain electrically neutral and to oppose any local concentration of
electric charge.}{Francis F.\ Chen, \emph{Introduction to Plasma
Physics and Controlled Fusion}, 3rd ed.\ (Springer, 2016), Section~1.4
\cite[Section~1.4]{text:chen-plasma-physics}.}

\chapmeta{Half-life: durable. Archetypes: stability, transport, scaling. Exercise target: 27. Project track: bench observation plus analysis (hybrid). Hazard class: medium (plasma globe and discharge tube safe at low voltage; industrial plasmas require RF safety, high-voltage interlocks, ventilation). Reader path default: standard, with sections explicitly tagged. Named cases: none directly; ITER disruption pattern described as a class.}

\noindent
A \engterm{plasma} is an ionised gas in which a sufficient fraction of
the constituent atoms or molecules carries net charge that the
collective electromagnetic behaviour of the charged particles
dominates the collisional behaviour of the neutral particles. Plasma
is the fourth state of matter at high temperature; it is also the
working state of fluorescent lights, neon signs, integrated-circuit
etching, plasma cutting, and the central engineering problem of
controlled fusion. This chapter develops the working content of
plasma physics for the practising engineer: the operational
definitions (Debye length, plasma frequency), the single-particle
motion in magnetic fields, the magnetohydrodynamic (MHD) fluid model,
the diagnostic techniques, and the industrial and fusion applications.

\section{What plasma is: ionisation, Debye length, plasma frequency}\pathtag{core}

A neutral gas heated above $\sim 10^{4}\,\si{\kelvin}$ begins to
ionise; the ionisation fraction follows the Saha equation,
\[
\frac{n_{i}n_{e}}{n_{n}} = \left(\frac{m_{e}k_{B}T}{2\pi\hbar^{2}}
\right)^{3/2}\exp(-E_{i}/k_{B}T),
\]
where $n_{i}$, $n_{e}$, $n_{n}$ are densities of ions, electrons,
neutrals, and $E_{i}$ is the ionisation energy of the species
($13.6\,\si{\electronvolt}$ for hydrogen). Below this temperature a
plasma can be created at lower energy density by non-equilibrium
discharge (electrical breakdown, RF excitation); fluorescent lamps
and industrial etching plasmas operate at electron temperatures of
$1$ to $10\,\si{\electronvolt}$ but with low ionisation fractions of
$10^{-6}$ to $10^{-3}$.

A plasma is operationally defined by three working criteria:
\begin{itemize}[leftmargin=*]
  \item Collective behaviour dominates individual-particle motion.
  \item The Debye length is much smaller than the plasma's spatial
    extent.
  \item There are many particles in a Debye sphere ($N_{D} \gg 1$).
\end{itemize}
These three criteria are not independent restatements; each rules out a
different way of failing to be a plasma. A gas with too few charges per
Debye sphere screens poorly and behaves as a collection of independent
two-body Coulomb collisions rather than a collective medium. A region
smaller than a Debye length cannot screen at all, so the quasineutrality
that defines plasma behaviour does not hold there. The ionisation fraction
itself enters only through $n_{e}$; a weakly ionised gas can still be a
good plasma if its charged-particle density is high enough, which is why
fluorescent lamps at $10^{-6}$ ionisation are genuine plasmas.

% Figure: Saha ionization fraction of atomic hydrogen versus temperature
% at a fixed reference density (1e20 per cubic metre). Data generated by
% code/saha_ionization.py into data/saha_hydrogen.csv.
\begin{figure}[htbp]
\centering
\begin{tikzpicture}
\begin{axis}[
  width=0.78\textwidth, height=0.5\textwidth,
  xlabel={temperature $T$ (K)},
  ylabel={ionization fraction $n_i/n_{\mathrm{tot}}$},
  xmin=4000, xmax=20000, ymin=0, ymax=1.05,
  grid=both, grid style={enggray!20},
  every axis plot/.append style={thick},
  legend pos=south east, legend cell align=left,
]
  \addplot[engred, mark=*, mark size=1.2pt] table[
    col sep=comma, x=temperature_K, y=ionization_fraction]
    {volumes/04-energy/ch13-plasma-physics/data/saha_hydrogen.csv};
  \addlegendentry{$n_{\mathrm{tot}}=10^{20}\,\mathrm{m^{-3}}$}
  \draw[dashed, enggray] (axis cs:4000,0.5) -- (axis cs:20000,0.5);
\end{axis}
\end{tikzpicture}
\caption[Saha ionization fraction of hydrogen]{Ionization fraction of
atomic hydrogen against temperature at a fixed total density of
$10^{20}\,\si{\per\cubic\meter}$, from the Saha equation. The transition
from neutral gas to nearly fully ionized plasma is sharp: hydrogen passes
from a few percent ionized to almost fully ionized over a temperature
range of only a few thousand kelvin, with half-ionization near
$10^{4}\,\si{\kelvin}$, far below the $13.6\,\si{\electronvolt}\approx
1.6\times10^{5}\,\si{\kelvin}$ that the ionization energy alone would
suggest. The shift comes from the large phase-space factor and the
exponential sensitivity to the tail of the energy distribution. The curve
moves to higher temperature as density rises.}
\label{fig:vol04ch13:saha-fraction}
\end{figure}


The Saha curve in Figure~\ref{fig:vol04ch13:saha-fraction} shows the
sharpness of the ionisation transition. Hydrogen passes from a few percent
ionised to nearly fully ionised over a few thousand kelvin, with
half-ionisation near $10^{4}\,\si{\kelvin}$, far below the $1.6\times
10^{5}\,\si{\kelvin}$ that the $13.6\,\si{\electronvolt}$ ionisation energy
alone would suggest. The phase-space prefactor and the exponential
sensitivity to the high-energy tail of the Maxwell distribution combine to
ionise the gas long before the mean thermal energy reaches the ionisation
threshold. The transition temperature also depends on density, moving
higher as the gas is compressed; the accompanying
\texttt{code/saha\_ionization.py} computes the curve at any reference
density.

The \engterm{Debye length} is the distance over which the plasma
screens an external charge. For a plasma of electron density $n_{e}$
and electron temperature $T_{e}$,
\[
\lambda_{D} = \sqrt{\frac{\varepsilon_{0}k_{B}T_{e}}{n_{e}e^{2}}},
\]
with $k_{B}T_{e}$ in joules. For a fluorescent lamp's plasma
($n_{e}\approx 10^{17}\,\si{\per\meter\cubed}$, $T_{e}\approx 1\,\si{
\electronvolt}$, $k_{B}T_{e} = 1.6\times 10^{-19}\,\si{\joule}$),
$\lambda_{D}\approx 30\,\si{\micro\meter}$; for the solar
photosphere ($n_{e}\approx 10^{19}\,\si{\per\meter\cubed}$,
$T_{e}\approx 6000\,\si{\kelvin}$, $k_{B}T_{e}\approx 0.5\,\si{
\electronvolt}$), $\lambda_{D}\approx 1.7\,\si{\micro\meter}$; for an
ITER fusion plasma ($n_{e}\approx 10^{20}\,\si{\per\meter\cubed}$,
$T_{e}\approx 10\,\si{\kilo\electronvolt}$), $\lambda_{D}\approx
70\,\si{\micro\meter}$. In every working plasma, $\lambda_{D}$ is
much smaller than the device dimension, validating the plasma
definition.

% Figure: Debye shielding. The bare Coulomb potential of a test charge
% and the screened Yukawa potential phi ~ (1/r) exp(-r/lambda_D) it takes
% on inside a plasma. Horizontal axis in units of the Debye length.
\begin{figure}[htbp]
\centering
\begin{tikzpicture}
\begin{axis}[
  width=0.78\textwidth, height=0.5\textwidth,
  xlabel={distance from test charge $r/\lambda_D$},
  ylabel={potential (normalised)},
  domain=0.25:5, samples=120,
  xmin=0, xmax=5, ymin=0, ymax=4.2,
  legend pos=north east,
  legend cell align=left,
  every axis plot/.append style={thick},
  grid=both, grid style={enggray!20},
]
  \addplot[engred] {1/x};
  \addlegendentry{bare Coulomb $\propto 1/r$}
  \addplot[engblue] {exp(-x)/x};
  \addlegendentry{screened $\propto e^{-r/\lambda_D}/r$}
  \draw[dashed, enggray] (axis cs:1,0) -- (axis cs:1,4.2);
  \node[enggray, font=\scriptsize, anchor=south west] at (axis cs:1.05,3.6)
    {$r=\lambda_D$};
\end{axis}
\end{tikzpicture}
\caption[Debye shielding of a test charge]{Debye shielding. A test charge
inserted into a plasma attracts a cloud of opposite charge that screens
its field. The bare Coulomb potential falls as $1/r$; the screened
potential the plasma actually presents falls as
$e^{-r/\lambda_D}/r$, dropping by a factor $1/e$ at one Debye length and
becoming negligible within a few $\lambda_D$. The plasma is electrically
quiet beyond the Debye sphere, which is the working content of quasineutrality.}
\label{fig:vol04ch13:debye-shielding}
\end{figure}


Figure~\ref{fig:vol04ch13:debye-shielding} shows what shielding does to a
charge inserted into a plasma. The bare Coulomb potential falls as $1/r$;
the potential the plasma actually presents falls as
$e^{-r/\lambda_{D}}/r$, dropping by a factor $1/e$ at one Debye length and
becoming negligible within a few $\lambda_{D}$. Beyond the Debye sphere
the plasma is electrically quiet, which is the working content of
quasineutrality: on length scales larger than $\lambda_{D}$, the positive
and negative charge densities track each other so closely that the net
charge density is effectively zero. The derivation of the Debye length
from the linearised Poisson-Boltzmann equation is the first exercise in
the derivation set.

The \engterm{plasma frequency} is the natural oscillation frequency
of electron-density perturbations:
\[
\omega_{p} = \sqrt{\frac{n_{e}e^{2}}{m_{e}\varepsilon_{0}}}.
\]
At $n_{e} = 10^{18}\,\si{\per\meter\cubed}$, $\omega_{p} = 5.6\times
10^{10}\,\si{\radian\per\second}$, or $f_{p} = 9\,\si{\giga\hertz}$.
The plasma frequency separates the regime in which the plasma is
opaque (incident waves below $\omega_{p}$ reflect) from the regime
in which it is transparent (waves above $\omega_{p}$ propagate). The
Earth's ionosphere, with $n_{e}$ from $10^{10}$ to $10^{12}\,\si{
\per\meter\cubed}$, has $f_{p}$ from $1$ to $10\,\si{\mega\hertz}$;
short-wave radio (HF band) below this frequency reflects off the
ionosphere, the working basis of long-distance amateur radio. VHF and
above transmit through.

The cutoff is the content of the plasma's refractive index. For an
electromagnetic wave of frequency $\omega$ in an unmagnetised plasma the
dispersion relation is $\omega^{2} = \omega_p^{2} + c^{2}k^{2}$, so the
refractive index is $n = ck/\omega = \sqrt{1 - \omega_p^{2}/\omega^{2}}$.
Above $\omega_p$ the index is real and the wave propagates with phase speed
above $c$ and group speed below it; below $\omega_p$ the index is imaginary,
the wavenumber is purely imaginary, and the wave decays without
propagating, which is reflection rather than absorption. The decay length is
the skin depth $c/\sqrt{\omega_p^{2}-\omega^{2}}$, which for a static field
($\omega\to 0$) reduces to the collisionless skin depth $c/\omega_p$. At
$n_e = 10^{18}\,\si{\per\meter\cubed}$ the skin depth is $c/\omega_p =
(3\times 10^{8})/(5.6\times 10^{10}) \approx 5\,\si{\milli\meter}$, the
distance an inductively coupled source's field penetrates the bulk before it
is screened, which sets where that source deposits its power. The same
$n=\sqrt{1-\omega_p^2/\omega^2}$ governs radar transmission through
re-entry plasma sheaths, the blackout that interrupts communication with a
spacecraft as it enters the atmosphere wrapped in plasma above the radio
frequency's cutoff density.

% Figure: the collisionless skin depth of an unmagnetised plasma against
% electron density, on log-log axes, with the fluorescent-lamp, processing,
% and fusion densities marked. The skin depth c/omega_p is the distance an
% evanescent wave below the plasma frequency penetrates before it is screened,
% and it sets the depth to which an inductively coupled source deposits power.
\begin{figure}[htbp]
\centering
\begin{tikzpicture}
\begin{loglogaxis}[
  width=0.82\textwidth, height=0.55\textwidth,
  xlabel={electron density $n_e$ (\si{\per\meter\cubed})},
  ylabel={skin depth $c/\omega_p$ (\si{\meter})},
  xmin=1e14, xmax=1e21, ymin=1e-4, ymax=1e0,
  grid=both, grid style={enggray!20},
  every axis plot/.append style={thick},
]
  % c/omega_p = c * sqrt(m_e eps0 / (n e^2)) = 5.31e12 / sqrt(n)
  \addplot[engblue, domain=1e14:1e21, samples=60]
    {5.31e12/sqrt(x)};
  % markers for representative devices
  \addplot[only marks, mark=*, engred, mark size=2pt] coordinates {
    (1e17,1.68e-2) (2e16,3.75e-2) (1e20,5.31e-4) };
  \node[engred, font=\scriptsize, anchor=west] at (axis cs:1e17,1.68e-2)
    {lamp};
  \node[engred, font=\scriptsize, anchor=west] at (axis cs:2e16,3.75e-2)
    {etch};
  \node[engred, font=\scriptsize, anchor=south west] at (axis cs:1e20,5.31e-4)
    {fusion core};
\end{loglogaxis}
\end{tikzpicture}
\caption[Collisionless skin depth versus density]{The collisionless skin
depth $c/\omega_p$ against electron density on log-log axes. A wave below the
plasma frequency does not propagate; it decays over this length, so the skin
depth is the distance an inductively coupled radio-frequency source drives its
field into the bulk before the plasma screens it, and therefore the depth over
which that source deposits its power. The depth falls as $n_e^{-1/2}$, from a
few centimetres in the tenuous corner to well below a millimetre in a fusion
core. The three marked points are the fluorescent lamp, the processing
discharge of the etch case, and the fusion-core density; a source that must
heat the bulk of a dense plasma cannot do it by a surface wave alone, which is
one reason a fusion device heats with neutral beams and resonant waves that
tunnel or mode-convert past the cutoff rather than with a simple induction
coil.}
\label{fig:vol04ch13:skin-depth}
\end{figure}


\paragraph{Worked example: the skin depth and where an RF source deposits its power.}
An inductively coupled plasma source wraps a coil around the chamber and drives
current in the gas by induction, so the question of how deep the field reaches
before the plasma screens it decides where the power lands. Below the plasma
frequency the wave is evanescent and decays over the collisionless skin depth
$\delta = c/\omega_p$, which with $\omega_p = \sqrt{n_e e^{2}/(m_e\varepsilon_0)}$
is $\delta = c\sqrt{m_e\varepsilon_0/(n_e e^{2})} = 5.31\times 10^{6}/\sqrt{n_e}$
in metres. For a processing discharge at $n_e = 2\times 10^{16}\,\si{\per\meter
\cubed}$, $\delta = 5.31\times 10^{12}/\sqrt{2\times 10^{16}} \approx 3.8\times
10^{-2}\,\si{\meter}$, about $38\,\si{\milli\meter}$, comparable to the chamber,
so the induced field fills the discharge and the power deposits through the bulk.
For a fusion-core density $n_e = 10^{20}\,\si{\per\meter\cubed}$ the same formula
gives $\delta \approx 5.3\times 10^{-4}\,\si{\meter}$, half a millimetre, so a
sub-cutoff wave dies in the first skin of the plasma and cannot heat the core at
all. Figure~\ref{fig:vol04ch13:skin-depth} carries the scaling across the whole
density range: the depth falls as $n_e^{-1/2}$, and the collapse of the skin
depth with density is the reason an inductive coil heats a processing plasma
through and through but a fusion device must reach its core with neutral beams or
with waves launched above the local cutoff. The same skin depth sets the
frequency at which a source becomes efficient: the coil must drive the plasma at
a frequency low enough that the skin depth spans a useful fraction of the
chamber, which places industrial inductive sources at the same $13.56\,\si{\mega
\hertz}$ band the capacitive tools use.

The number of particles in a Debye sphere is
\[
N_{D} = \frac{4}{3}\pi \lambda_{D}^{3} n_{e},
\]
and the \engterm{plasma parameter} $\Lambda = n_{e}\lambda_{D}^{3}$.
For a working plasma, $\Lambda \gg 1$; for laboratory plasmas
$\Lambda$ is typically $10^{6}$ to $10^{12}$. The plasma parameter
controls the importance of two-particle collisions versus collective
mean-field behaviour; large $\Lambda$ is the regime of
collisionless plasma dominated by Vlasov-equation dynamics.

The plasma parameter has a partner that reads the same physics from the energy
side, the \engterm{coupling parameter} $\Gamma$, the ratio of the mean Coulomb
interaction energy between neighbours to the mean thermal energy. With a mean
interparticle spacing $a = (4\pi n_e/3)^{-1/3}$, the coupling parameter is
$\Gamma = e^{2}/(4\pi\varepsilon_0 a k_B T_e)$. A plasma is \engterm{weakly
coupled}, the ordinary regime of this chapter, when $\Gamma \ll 1$, so the
thermal energy dominates and particles fly past one another with only small
deflections; it is \engterm{strongly coupled} when $\Gamma \gtrsim 1$, so the
Coulomb energy orders the particles and the gas begins to behave like a liquid
or a crystal. The two parameters are inverses of the same smallness, $\Gamma
\sim \Lambda^{-2/3}$, so a large Debye number and a small coupling parameter are
one statement made twice.

\paragraph{Worked example: coupling parameter of a lamp, a fusion core, and a white dwarf.}
Evaluate $\Gamma$ across three plasmas to see where the weak-coupling picture
holds. For the fluorescent lamp, $n_e = 10^{17}\,\si{\per\meter\cubed}$ and
$T_e = 1\,\si{\electronvolt}$, the spacing is $a = (4\pi\times 10^{17}/3)^{-1/3}
\approx 1.3\times 10^{-6}\,\si{\meter}$, and $\Gamma = (1.6\times 10^{-19})^{2}/
[4\pi(8.85\times 10^{-12})(1.3\times 10^{-6})(1.6\times 10^{-19})] \approx
1.1\times 10^{-3}$, comfortably weakly coupled. For the fusion core, $n_e =
10^{20}$ and $T_e = 10^{4}\,\text{eV}$, the higher density is overwhelmed by the
far higher temperature and $\Gamma \approx 1.1\times 10^{-6}$, weaker still, which
is why a fusion plasma is an almost ideal gas of charges. A white-dwarf interior,
by contrast, sits at $n_e \sim 10^{36}\,\si{\per\meter\cubed}$ and $T \sim
10^{7}\,\si{\kelvin}$, giving $\Gamma$ of order tens to hundreds: the ions there
form a Coulomb crystal, a genuinely strongly coupled plasma outside the reach of
the weak-coupling formulas. The engineer reading a new plasma computes $\Gamma$
first, because every collisionless and mean-field result of this chapter assumes
$\Gamma \ll 1$, and the dusty-plasma Coulomb crystals of a later worked example
are the one laboratory case where that assumption fails and the strong-coupling
physics takes over.

% Figure: plasma-regime chart on the density (x) versus temperature (y)
% plane, log-log. Marks where laboratory, industrial, space, and fusion
% plasmas sit. Each point and label placed explicitly for compile safety.
\begin{figure}[htbp]
\centering
\begin{tikzpicture}
\begin{loglogaxis}[
  width=0.82\textwidth, height=0.58\textwidth,
  xlabel={electron density $n_e$ ($\mathrm{m^{-3}}$)},
  ylabel={electron temperature $T_e$ (eV)},
  xmin=1e6, xmax=1e32, ymin=1e-2, ymax=1e6,
  grid=both, grid style={enggray!20},
]
  \addplot[only marks, mark=*, engred, mark size=2pt] coordinates {
    (1e20,10000) (1e31,5000) (1e23,1)
  };
  \addplot[only marks, mark=square*, engblue, mark size=2pt] coordinates {
    (1e11,0.1) (1e6,100)
  };
  \addplot[only marks, mark=triangle*, englime!60!black, mark size=2.4pt]
    coordinates { (1e16,1) (1e17,3) (1e14,5) };
  % labels
  \node[font=\scriptsize, anchor=south] at (axis cs:1e20,10000) {tokamak};
  \node[font=\scriptsize, anchor=south] at (axis cs:1e31,5000) {inertial};
  \node[font=\scriptsize, anchor=west] at (axis cs:1e23,1) {solar core};
  \node[font=\scriptsize, anchor=west] at (axis cs:1e11,0.1) {ionosphere};
  \node[font=\scriptsize, anchor=west] at (axis cs:1e6,100) {solar wind};
  \node[font=\scriptsize, anchor=east] at (axis cs:1e16,1) {lamp};
  \node[font=\scriptsize, anchor=west] at (axis cs:1e17,3) {etch};
  \node[font=\scriptsize, anchor=east] at (axis cs:1e14,5) {arc};
\end{loglogaxis}
\end{tikzpicture}
\caption[Plasma regimes on the density-temperature plane]{Where plasmas
live on the density-temperature plane. Engineered and natural plasmas span
more than twenty orders of magnitude in density and seven in temperature:
the near-vacuum, room-temperature solar wind; the low-density, cool
ionosphere; the dense, modestly hot discharges of fluorescent lamps and
reactive-ion-etch tools; and the hot, dense plasmas of magnetic-confinement
(tokamak) and inertial-confinement fusion. Round markers are hot or dense
plasmas, square markers are space plasmas, and triangular markers are
low-temperature laboratory and industrial plasmas. Every point obeys the
same operational tests of Section 13.1: a Debye length far below the device
size and many particles per Debye sphere. The chart is the map an engineer
reads to know which physics governs a given device.}
\label{fig:vol04ch13:plasma-regimes}
\end{figure}


Engineered and natural plasmas occupy a vast range of density and
temperature, shown in Figure~\ref{fig:vol04ch13:plasma-regimes}. The
near-vacuum solar wind, the cool ionosphere, the modest discharges of
fluorescent lamps and etch tools, and the hot dense cores of fusion
devices all satisfy the same three operational tests. The chart is the map
an engineer reads to decide which description applies: a collisional fluid
near the dense, cool corner; a collisionless kinetic plasma in the hot,
tenuous corner; the magnetohydrodynamic single-fluid model wherever the
magnetic field organises the motion. The companion
\texttt{code/plasma\_parameters.py} computes $\lambda_{D}$, $\omega_{p}$,
and $N_{D}$ for each entry and confirms $N_{D}\gg 1$ across the whole
range, with the solar core sitting at the degenerate edge where the
classical plasma picture begins to fail.

A plasma carries two electrostatic waves that the same length and time
scales already named, and an engineer meets both. The fast one is the
\engterm{Langmuir wave}, the electron oscillation of the plasma frequency
given a small thermal correction. A cold plasma oscillates at exactly
$\omega_p$ regardless of wavelength, but electron pressure lets a
compression propagate, and the Bohm-Gross relation $\omega^{2} =
\omega_p^{2}(1 + 3k^{2}\lambda_D^{2})$ adds a wavelength dependence that
becomes strong once the wavelength approaches the Debye length. The slow one
is the \engterm{ion-acoustic wave}, a sound wave in which the light electrons
supply the restoring pressure and the heavy ions supply the inertia, so it
travels at the ion-sound speed $c_s = \sqrt{k_B T_e/m_i}$, slower than the
Langmuir wave by the square root of the mass ratio. The two branches,
Figure~\ref{fig:vol04ch13:wave-dispersion}, split the plasma's response into a
fast electron oscillation pinned near $\omega_p$ and a slow compressional ion
wave that rises linearly from the origin. A diagnostic or a heating scheme
couples to whichever branch its frequency matches, which is why the ion-sound
speed reappears below as the Bohm velocity that governs every plasma sheath.

% Figure: dispersion relations of the two electrostatic plasma waves. The
% electron Langmuir wave sits at the plasma frequency and is nearly flat at
% small wavenumber, while the ion-acoustic wave rises linearly from the origin
% at the ion-sound speed. The two branches separate the fast electron
% oscillation from the slow compressional ion wave.
\begin{figure}[htbp]
\centering
\begin{tikzpicture}
\begin{axis}[
  width=0.82\textwidth, height=0.55\textwidth,
  xlabel={wavenumber $k\lambda_D$},
  ylabel={frequency $\omega/\omega_p$},
  xmin=0, xmax=1.0, ymin=0, ymax=1.6,
  grid=both, grid style={enggray!20},
  legend pos=north west,
  legend cell align={left},
  every axis plot/.append style={thick},
]
  % Langmuir (Bohm-Gross): omega^2 = omega_p^2 (1 + 3 (k lambda_D)^2)
  \addplot[engblue, domain=0:1.0, samples=80]
    {sqrt(1 + 3*x*x)};
  \addlegendentry{Langmuir (Bohm-Gross)}
  % ion-acoustic: omega = k c_s, with c_s/(omega_p lambda_D) = sqrt(m_e/m_i)
  % for a light illustrative mass ratio the slope is set to 0.35 for visibility
  \addplot[engorange, domain=0:1.0, samples=80]
    {0.35*x};
  \addlegendentry{ion-acoustic}
  % plasma frequency reference line
  \draw[dashed, enggray] (axis cs:0,1) -- (axis cs:1.0,1);
  \node[enggray, font=\scriptsize, anchor=south east] at (axis cs:1.0,1.0)
    {$\omega_p$};
\end{axis}
\end{tikzpicture}
\caption[Electrostatic wave branches]{The two electrostatic waves a plasma
carries. The electron Langmuir branch (blue) follows the Bohm-Gross relation
$\omega^{2}=\omega_p^{2}(1+3k^{2}\lambda_D^{2})$, so it starts at the plasma
frequency $\omega_p$ at long wavelength and stiffens as the wavelength shrinks
toward the Debye length. The ion-acoustic branch (orange) rises linearly from
the origin at the ion-sound speed $c_s=\sqrt{k_B T_e/m_i}$, a slow
compressional wave in which the light electrons provide the pressure and the
heavy ions provide the inertia. The slope shown is exaggerated for visibility;
the true ion-acoustic branch is slower than the Langmuir branch by the square
root of the mass ratio. The two branches split the plasma's wave response into
a fast electron oscillation near $\omega_p$ and a slow ion sound wave, and a
diagnostic or a heating scheme couples to one or the other by matching its
frequency to the branch it wants.}
\label{fig:vol04ch13:wave-dispersion}
\end{figure}


\paragraph{Worked example: ion-sound speed and the Bohm velocity.}
The ion-acoustic speed and the Bohm velocity that sets the sheath edge are the
same quantity read in two contexts. For an argon processing plasma at $T_e =
3\,\si{\electronvolt}$ ($k_B T_e = 4.8\times 10^{-19}\,\si{\joule}$) and cold
ions, the sound speed is $c_s = \sqrt{k_B T_e/m_i} = \sqrt{(4.8\times
10^{-19})/(6.6\times 10^{-26})} \approx 2.7\times 10^{3}\,\si{\meter\per
\second}$, exactly the Bohm velocity the etch case used. For a deuterium fusion
edge plasma at $T_e = 100\,\si{\electronvolt}$ the same formula with $m_D =
3.34\times 10^{-27}\,\si{\kilogram}$ gives $c_s = \sqrt{(1.6\times
10^{-17})/(3.34\times 10^{-27})} \approx 6.9\times 10^{4}\,\si{\meter\per
\second}$, about twenty-five times faster, since the temperature is up a factor
of thirty and the mass down a factor of twenty. The ion-acoustic speed rises as
$\sqrt{T_e}$, so a hotter edge drives ions into the divertor faster, which sets
the parallel heat flux the divertor must survive. The same $c_s$ is the speed
at which a density perturbation propagates along an unmagnetised discharge, so
the time for a probe or a langmuir sweep to sample a disturbance is the chamber
size divided by $c_s$, a few microseconds in the argon plasma above.

\begin{mastery}
The fluid picture of the two wave branches hides a kinetic effect that a fluid
model cannot produce: a wave can damp on a collisionless plasma with no
dissipation at all, by \engterm{Landau damping}. The mechanism is a resonance
between the wave's phase velocity and the particles moving at that speed. A
particle slightly slower than the wave is accelerated by it and takes energy
from the wave; a particle slightly faster gives energy back. In a Maxwellian
distribution there are more particles just below any given speed than just
above, because the distribution falls with speed, so the net exchange takes
energy out of the wave and damps it. The damping rate depends on the slope of
the velocity distribution at the phase velocity, $\gamma \propto \partial f/
\partial v|_{v=\omega/k}$, which is the sense in which Landau damping is a
kinetic effect invisible to any moment-based fluid closure. Two engineering
consequences follow. A Langmuir wave with phase velocity a few thermal speeds
above the bulk damps weakly and can propagate, which is why the Bohm-Gross
branch is observable, while a wave with phase velocity near the thermal speed
damps in a few oscillations and cannot. And a distribution with a positive
slope, an inverted population such as a beam injected into a plasma, drives the
wave rather than damping it, the bump-on-tail instability that underlies
beam-plasma interactions and some radio-frequency heating schemes. The engineer
who must decide whether a launched wave will reach the plasma core or damp at
the edge computes the phase velocity against the local temperature and reads the
answer off the slope of the distribution, a calculation the fluid model of the
rest of this chapter cannot pose.
\end{mastery}

Landau damping removes energy from a wave when the velocity distribution
falls at the phase velocity. The same resonance run in reverse feeds energy
into a wave when the distribution rises there, and the cleanest case an
engineer meets is two electron populations streaming through each other. A
cold beam fired into a stationary plasma, or two counter-streaming beams,
carries free energy in the relative motion, and a plasma wave sitting near
the resonant wavenumber $k v_0 = \omega_p$ taps that energy and grows. The
cold dispersion relation for two equal beams of speed $\pm v_0$ has a complex
root over a band of wavenumbers below the cutoff, with a peak growth rate
$\gamma_{\max}\approx\omega_p/2\sqrt{2}$, so the disturbance e-folds in a few
plasma periods, faster than any collision. Figure~\ref{fig:vol04ch13:two-stream}
draws the growth rate across the unstable band.

% Figure: the two-stream instability growth rate against normalised
% wavenumber. Two counter-streaming electron populations feed free energy
% into a plasma wave; the growth rate is non-zero over a band of wavenumbers
% below the cutoff k v0 / omega_p = 1, peaking near unity, and it is the
% collisionless free-energy release that Landau damping is the stable
% counterpart of. The curve is the growth rate normalised to the plasma
% frequency for two equal cold beams of speed v0.
\begin{figure}[htbp]
\centering
\begin{tikzpicture}
\begin{axis}[
  width=0.82\textwidth, height=0.52\textwidth,
  xlabel={normalised wavenumber $k v_0/\omega_p$},
  ylabel={growth rate $\gamma/\omega_p$},
  xmin=0, xmax=1.4, ymin=0, ymax=0.45,
  grid=both, grid style={enggray!20},
  every axis plot/.append style={thick},
]
  % Cold two-stream dispersion for two equal beams of density n/2, speed +-v0:
  % the peak growth is gamma_max/omega_p = 1/(2 sqrt2) ~ 0.354 near k v0/omega_p ~ 1.
  % A smooth representative envelope that vanishes at k=0 and at the cutoff x=1.2.
  \addplot[engblue, domain=0:1.155, samples=80]
    {0.354*sqrt(max(0, 1 - (x*x-1)*(x*x-1)/0.25))};
  \addplot[engred, dashed, domain=0:1.4, samples=2] coordinates {(1,0)(1,0.45)};
  \node[engred, font=\scriptsize, anchor=south west, rotate=90]
    at (axis cs:1,0.05) {resonant $k v_0=\omega_p$};
  \node[engblue, font=\scriptsize, anchor=south]
    at (axis cs:1,0.37) {$\gamma_{\max}\approx\omega_p/2\sqrt{2}$};
\end{axis}
\end{tikzpicture}
\caption[Two-stream instability growth rate]{Growth rate of the cold
two-stream instability against wavenumber, normalised to the plasma
frequency, for two equal electron beams counter-streaming at speed $v_0$.
The instability grows only over a band of wavenumbers around the resonant
value $k v_0 = \omega_p$, where a wave sits still in the beam frame and
draws free energy from the relative streaming. The peak growth rate is
$\gamma_{\max}\approx\omega_p/2\sqrt{2}$, so the wave e-folds in a few plasma
periods, faster than any collisional process. The two-stream instability is
the free-energy release that Landau damping is the stable counterpart of: a
distribution with a positive slope at the phase velocity drives the wave, a
distribution that falls with speed damps it, and reading the slope decides
which. The same mechanism sets the electron-beam limit in beam-plasma
amplifiers and seeds the turbulence that carries anomalous transport.}
\label{fig:vol04ch13:two-stream}
\end{figure}


\paragraph{Worked example: the two-stream growth rate and where a beam-plasma device sits.}
Take an electron beam of density $n_b = 10^{15}\,\si{\per\meter\cubed}$
streaming at $v_0 = 3\times 10^{7}\,\si{\meter\per\second}$ through a
background plasma of $n_0 = 10^{17}\,\si{\per\meter\cubed}$. The background
plasma frequency is $\omega_p = \sqrt{n_0 e^{2}/(m_e\varepsilon_0)} \approx
1.8\times 10^{10}\,\si{\radian\per\second}$. For a weak beam the maximum
growth rate is not the equal-beam value but the reduced beam-plasma rate
$\gamma_{\max}\approx (\sqrt{3}/2)(n_b/2n_0)^{1/3}\omega_p$, so with $n_b/2n_0
= 5\times 10^{-3}$ the cube root is $0.17$ and $\gamma_{\max}\approx
(0.87)(0.17)(1.8\times 10^{10}) \approx 2.7\times 10^{9}\,\si{\radian\per
\second}$, an e-folding time of $\sim 0.4\,\si{\nano\second}$. The instability
peaks at the wavenumber where the wave is resonant with the beam,
$k\approx\omega_p/v_0 = 1.8\times 10^{10}/3\times 10^{7} \approx 600\,\si{
\per\meter}$, a wavelength of about a centimetre. A device built to exploit
this, a travelling-wave beam-plasma amplifier, sizes its interaction length so
that the wave grows through several e-foldings across it; a device that must
avoid it, such as a neutral-beam duct or an electron gun, keeps the beam
below the density at which the growth length falls short of the path. The same
arithmetic sets the length over which a runaway-electron beam in a disrupting
tokamak destabilises the plasma it runs through.

\begin{archetype}[Stability and scaling]
A plasma is a system whose stability is the entire engineering
problem. Three scales set the operational physics: $\lambda_{D}$
(charge screening), $\omega_{p}^{-1}$ (collective oscillation time),
and $r_{L}$ (Larmor radius for magnetic confinement, Section 13.2).
The discipline of plasma engineering is to identify which scales
constrain a particular device and to design for stability at every
length and time scale that the plasma populates.
\end{archetype}

\section{Magnetised plasmas: cyclotron motion and drifts}\pathtag{standard}

A charged particle of mass $m$ and charge $q$ in a uniform magnetic
field $\vect{B}$ experiences the Lorentz force $q\vect{v}\times
\vect{B}$, which is perpendicular to $\vect{v}$ and $\vect{B}$. The
particle moves in a circle perpendicular to $\vect{B}$ at the
\engterm{cyclotron frequency} (or gyrofrequency)
\[
\omega_{c} = \frac{|q|B}{m},
\]
with radius (the \engterm{Larmor radius})
\[
r_{L} = \frac{v_{\perp}}{\omega_{c}} = \frac{mv_{\perp}}{|q|B}.
\]
For an electron at $1\,\si{\electronvolt}$ thermal energy in a
$1\,\si{\tesla}$ field, $v_{\perp}\approx 6\times 10^{5}\,\si{\meter
\per\second}$, $\omega_{c}/(2\pi)\approx 28\,\si{\giga\hertz}$, and
$r_{L}\approx 3\,\si{\micro\meter}$. For a proton in the same
conditions ($v_{\perp}\approx 1.4\times 10^{4}\,\si{\meter\per
\second}$), $\omega_{c}/(2\pi)\approx 15\,\si{\mega\hertz}$ and
$r_{L}\approx 150\,\si{\micro\meter}$. The mass ratio $m_{p}/m_{e}
\approx 1836$ means that on any given field, electrons gyrate much
faster and on much smaller radii than ions.

% Figure: charged-particle motion in a magnetic field. Left: pure gyration
% (B out of page), opposite senses for ion and electron. Right: cycloidal
% E x B drift, the guiding centre marching steadily perpendicular to both
% E and B with the same velocity for both signs of charge.
\begin{figure}[htbp]
\centering
\begin{tikzpicture}[>=latex, scale=1.0]
% ===== Left panel: gyration =====
\begin{scope}[shift={(0,0)}]
  \node[font=\small] at (0,2.6) {gyration ($\vect{B}$ out of page)};
  % B field dots
  \foreach \x in {-1.6,-0.8,0,0.8,1.6}{
    \foreach \y in {-1.6,-0.8,0.8,1.6}{
      \node[engblue, font=\tiny] at (\x,\y) {$\odot$};
    }
  }
  % ion orbit (large, clockwise)
  \draw[engred, thick] (0.9,0) circle (0.9);
  \draw[engred, ->, thick] (0.9,0.9) -- (1.0,0.9);
  \node[engred, font=\scriptsize] at (0.9,0) {ion};
  % electron orbit (small, counter-clockwise)
  \draw[englime!60!black, thick] (-1.0,0) circle (0.45);
  \draw[englime!60!black, ->, thick] (-1.0,0.45) -- (-1.1,0.45);
  \node[englime!60!black, font=\scriptsize] at (-1.0,-0.85) {e$^-$};
\end{scope}
% ===== Right panel: E x B drift =====
\begin{scope}[shift={(6.2,0)}]
  \node[font=\small] at (0,2.6) {$\vect{E}\times\vect{B}$ drift};
  % E field arrows pointing up
  \draw[engorange, ->, thick] (-2.2,-1.8) -- (-2.2,1.8);
  \node[engorange, font=\scriptsize] at (-2.5,1.6) {$\vect{E}$};
  \node[engblue, font=\tiny] at (-1.6,-1.8) {$\odot\,\vect{B}$};
  % cycloid (guiding centre drifts to the right)
  \draw[engblue, thick, samples=160, domain=0:9, variable=\t]
    plot ({-1.4 + 0.42*\t - 0.42*sin(deg(\t))}, {0.42 - 0.42*cos(deg(\t))});
  % drift arrow
  \draw[enggray, ->, thick] (-1.4,-1.4) -- (1.9,-1.4)
    node[midway, below, font=\scriptsize] {$v_{E\times B}=E/B$};
\end{scope}
\end{tikzpicture}
\caption[Gyration and $E\times B$ drift]{Single-particle motion in a
magnetic field. Left: with $\vect{B}$ out of the page, ions and electrons
gyrate in opposite senses; the electron orbit is far smaller because the
Larmor radius scales with mass. Right: an added perpendicular electric
field bends each gyro-orbit into a cycloid whose guiding centre drifts at
$v_{E\times B}=E/B$, perpendicular to both $\vect{E}$ and $\vect{B}$ and
independent of charge, mass, and energy, so ions and electrons drift
together with no net current.}
\label{fig:vol04ch13:gyration-drift}
\end{figure}


Figure~\ref{fig:vol04ch13:gyration-drift} shows the two motions side by
side. With $\vect{B}$ out of the page, ions and electrons gyrate in
opposite senses, and the electron orbit is far smaller because the Larmor
radius scales with mass. Add a perpendicular electric field and each
gyro-orbit bends into a cycloid whose guiding centre marches steadily at
the $E\times B$ drift velocity.

The particle's motion parallel to $\vect{B}$ is unaffected by the
magnetic field. The combined motion is a helix: circular gyration
in the plane perpendicular to $\vect{B}$ at $\omega_{c}$, free motion
parallel to $\vect{B}$ at $v_{\parallel}$. The \engterm{guiding
centre} is the average position over a gyro-orbit, which drifts under
external forces.

When an electric field $\vect{E}$ is present perpendicular to
$\vect{B}$, the guiding centre drifts at the \engterm{$E\times B$
drift velocity}
\[
\vect{v}_{E\times B} = \frac{\vect{E}\times \vect{B}}{B^{2}}.
\]
This drift is independent of charge, mass, and energy; electrons and
ions drift together. The $E\times B$ drift is the foundation of MHD
in which the plasma moves as a single fluid. Two further drifts,
\engterm{gradient-$B$} and \engterm{curvature drifts}, are
charge-dependent and act to separate electrons from ions in
inhomogeneous fields, producing currents and self-consistent
electric fields. The drift formulas are the working content of
magnetised-plasma engineering; their detailed derivation is in
\cite{text:chen-plasma-physics, text:goldston-rutherford-plasma}.

% Figure: single-particle motion in crossed E and B fields. An ion and an
% electron, gyrating in opposite senses, both drift in the same direction
% at v = E x B / B^2. The cycloidal guiding-centre march is the content of
% the worked example on crossed-field drift.
\begin{figure}[htbp]
\centering
\begin{tikzpicture}[scale=1.0]
  % axes for orientation
  \draw[->, enggray] (0,-1.6) -- (0,1.9) node[above, font=\scriptsize] {$y$};
  \draw[->, enggray] (-0.3,-1.4) -- (10.5,-1.4) node[right, font=\scriptsize] {$x$};
  % field labels: B out of page, E pointing +y
  \node[engblue, font=\scriptsize] at (10.0,1.4) {$\vect{B}$ out of page};
  \draw[engred, ->, thick] (9.3,-1.0) -- (9.3,0.2)
    node[midway, right, font=\scriptsize] {$\vect{E}$};
  % drift direction arrow (E x B is in +x for E=+y, B=+z)
  \draw[engorange, ->, very thick] (0.4,1.55) -- (3.0,1.55)
    node[midway, above, font=\scriptsize] {$\vect{v}_{E\times B}=\vect{E}\times\vect{B}/B^{2}$};
  % ion cycloid (broad arcs, drifting right)
  \draw[engred, thick, domain=0:9.4, samples=160, smooth, variable=\t]
    plot ({\t},{0.7 - 0.7*cos(\t*60)});
  \node[engred, font=\scriptsize, anchor=north] at (1.0,-0.15) {ion};
  % electron cycloid (small tight arcs, same drift direction)
  \draw[engblue, thick, domain=0:9.4, samples=320, smooth, variable=\t]
    plot ({\t},{-0.9 + 0.22 - 0.22*cos(\t*220)});
  \node[engblue, font=\scriptsize, anchor=north] at (1.0,-1.15) {electron};
\end{tikzpicture}
\caption[Crossed-field drift]{Single-particle motion in crossed electric
and magnetic fields, with $\vect{B}$ out of the page and $\vect{E}$ pointing
upward. The ion (upper, wide arcs) and the electron (lower, tight arcs)
gyrate in opposite senses and on radii that differ by the square root of
the mass ratio, yet both guiding centres march in the same direction at the
$E\times B$ drift velocity $\vect{v}_{E\times B}=\vect{E}\times\vect{B}/B^{2}$.
The drift is independent of charge, mass, and energy, so it carries the
plasma bodily without separating ions from electrons and produces no net
current. This charge-independence is what lets the single-fluid MHD picture
treat the bulk plasma motion as one velocity field.}
\label{fig:vol04ch13:crossed-field-drift}
\end{figure}


Figure~\ref{fig:vol04ch13:crossed-field-drift} carries the crossed-field
picture to its working conclusion. The ion and the electron trace cycloids
of very different size, set by Larmor radii that differ by the square root
of the mass ratio, and they wind in opposite senses because their charges
have opposite sign. The guiding centres march together at one velocity that
does not care about charge, mass, or energy. That single shared velocity is
the bridge to the fluid model: where the $E\times B$ drift dominates the
perpendicular motion, the two-species plasma moves as one velocity field,
which is exactly what the MHD momentum equation assumes.

\paragraph{Worked example: crossed-field drift and the cycloid amplitude.}
Take a $1\,\si{\tesla}$ field out of the page and a perpendicular electric
field $E = 5\,\si{\kilo\volt\per\meter}$. The drift speed is $v_{E\times B}
= E/B = 5\times 10^{3}/1 = 5\times 10^{3}\,\si{\meter\per\second}$, the same
for both species and directed along $\vect{E}\times\vect{B}$. The gyration
superimposed on this drift has, for a cold particle released from rest, a
speed equal to the drift speed itself, so the orbit is a true cycloid whose
loops just touch the starting line. The amplitude of the loop is the Larmor
radius evaluated at $v_\perp = v_{E\times B}$: for a deuteron, $r_{L} =
m_{D}v_{E\times B}/(eB) = (3.34\times 10^{-27})(5\times 10^{3})/[(1.6\times
10^{-19})(1)] \approx 1.0\times 10^{-4}\,\si{\meter}$, about
$0.1\,\si{\milli\meter}$; for an electron the loop is smaller by the mass
ratio, $\sim 2.8\times 10^{-8}\,\si{\meter}$. The two loops differ by the full
mass ratio $m_{D}/m_{e}$ in radius but share the same forward march,
which is why a probe that resolves only the bulk motion sees a single
drifting fluid and not two species.

The general drift is one formula and the rest are special cases of it.
Any force $\vect{F}$ perpendicular to $\vect{B}$ pushes the guiding centre
across the field at
\[
\vect{v}_{F} = \frac{1}{q}\,\frac{\vect{F}\times\vect{B}}{B^{2}}.
\]
The derivation repeats the $E\times B$ argument with $q\vect{E}$ replaced by
$\vect{F}$: write the velocity as a steady drift plus gyration, choose the
drift so the gyration equation carries no net transverse force, cross the
balance $\vect{F} + q\vect{v}_{F}\times\vect{B} = 0$ with $\vect{B}$, and
use $\vect{v}_{F}\cdot\vect{B}=0$. Setting $\vect{F} = q\vect{E}$ recovers
the charge-independent $E\times B$ drift, because the $q$ in the force
cancels the $1/q$ in front. Any force that does not scale with charge, by
contrast, gives a drift that carries the sign of $q$ and so separates ions
from electrons and drives a current.

The \engterm{gradient-$B$ drift} is the first such case. When the field
strength varies across the orbit, the Larmor radius is larger on the
weak-field side of the loop than on the strong-field side, so the loop does
not close and the guiding centre walks sideways. Averaging the Lorentz
force over a gyro-orbit in a field with a transverse gradient gives an
effective force $\vect{F} = -\mu\,\grad B$, with $\mu = mv_{\perp}^{2}/(2B)$
the magnetic moment, and the general drift formula then gives
\[
\vect{v}_{\grad B} = \frac{m v_{\perp}^{2}}{2qB}\,
\frac{\vect{B}\times\grad B}{B^{2}}.
\]
The \engterm{curvature drift} is the second case. A particle streaming
along a curved field line of radius of curvature $R_{c}$ feels a
centrifugal force $\vect{F} = m v_{\parallel}^{2}\hat{\vect{R}}_{c}/R_{c}$
in the guiding-centre frame, and the general formula gives
\[
\vect{v}_{R} = \frac{m v_{\parallel}^{2}}{qB^{2}}\,
\frac{\vect{R}_{c}\times\vect{B}}{R_{c}^{2}}.
\]
Both drifts scale as $1/q$ and so point in opposite directions for ions and
electrons. In the curved field of a torus this charge separation builds a
vertical electric field across the plasma, and the resulting downward
$E\times B$ drift would push the whole plasma into the outer wall in
microseconds. The cure is the helical twist of the tokamak field: a field
line that spirals from the top of the plasma to the bottom and back
short-circuits the charge separation along the field, which is the reason a
purely toroidal field cannot confine a plasma and a poloidal component is
required.

% Figure: gradient-B drift. In a field that weakens from left to right the
% Larmor radius is larger where B is smaller, so successive gyro-loops do not
% close and the guiding centre marches sideways. Ions and electrons drift in
% opposite directions because the gyration sense depends on charge, so the
% drift carries a net current.
\begin{figure}[htbp]
\centering
\begin{tikzpicture}[scale=1.0]
  % field-strength indication: tightly packed dots on the left, sparse right
  \foreach \x in {0,0.45,...,6.0} {
    \foreach \y in {-1.4,-0.7,...,1.5} {
      \node[engblue, font=\tiny] at (\x,\y) {$\cdot$};
    }
  }
  \node[engblue, font=\scriptsize, anchor=south west] at (0,1.7)
    {strong $B$ (small $r_L$)};
  \node[engblue, font=\scriptsize, anchor=south east] at (6,1.7)
    {weak $B$ (large $r_L$)};
  \draw[engblue, ->] (0,-1.9) -- (6,-1.9)
    node[midway, below, font=\scriptsize] {decreasing $|\vect{B}|$, i.e.\ $\grad B$};
  % ion gyro-orbits of growing radius, guiding centre marching right (upper)
  \draw[engred, thick]
    (0.4,0.9) arc (180:-160:0.22)
    (1.1,0.95) arc (180:-150:0.30)
    (2.0,1.0) arc (180:-140:0.40)
    (3.1,1.05) arc (180:-130:0.52);
  \draw[engred, ->, thick] (4.0,0.9) -- (5.2,0.9);
  \node[engred, font=\scriptsize, anchor=south] at (4.6,0.95)
    {ion drift};
  % electron gyro-orbits, guiding centre marching left (lower)
  \draw[engorange, thick]
    (5.6,-0.9) arc (0:340:0.22)
    (4.9,-0.95) arc (0:330:0.30)
    (4.0,-1.0) arc (0:320:0.40);
  \draw[engorange, ->, thick] (2.9,-0.9) -- (1.7,-0.9);
  \node[engorange, font=\scriptsize, anchor=south] at (2.3,-0.85)
    {electron drift};
\end{tikzpicture}
\caption[Gradient-B drift]{The gradient-$B$ drift in a field that weakens
from left to right. A gyrating particle has a larger Larmor radius on the
weak-field part of its orbit than on the strong-field part, so the loop does
not close on itself and the guiding centre steps sideways, perpendicular to
both $\vect{B}$ and $\grad B$. The step size scales with $r_L$ and with the
fractional change in field over a Larmor radius. Ions (upper, red) and
electrons (lower, orange) gyrate in opposite senses, so they drift in
opposite directions, and the gradient drift therefore carries a net
current, in contrast to the charge-independent $E\times B$ drift. The same
charge separation, produced by gradient and curvature drifts in a toroidal
field, is what a tokamak must short out with its helical field twist to keep
the plasma confined.}
\label{fig:vol04ch13:gradb-drift}
\end{figure}


Figure~\ref{fig:vol04ch13:gradb-drift} shows the gradient drift building from
the growing Larmor radius across the orbit. The companion picture of the
torus is the working content of the next chapter on confinement: every
toroidal device must arrange its field so that the gradient and curvature
drifts of ions and electrons do not accumulate into a confinement-destroying
charge separation.

\paragraph{Worked example: gradient and curvature drifts in a tokamak.}
Take a deuteron at $T = 10\,\si{\kilo\electronvolt}$ in the $B = 5\,\si{
\tesla}$ field of a tokamak with major radius $R_{c} = 6\,\si{\meter}$, the
scale of the field curvature. Its thermal speed components are $v_{\perp}
\approx v_{\parallel} \approx \sqrt{2k_{B}T/m_{D}} \approx 9.8\times
10^{5}\,\si{\meter\per\second}$ (from the worked example below). The field
gradient across the torus is of order $\grad B/B \approx 1/R_{c}$, so the
gradient-drift speed is
\[
v_{\grad B} = \frac{m_{D}v_{\perp}^{2}}{2eB}\frac{|\grad B|}{B}
\approx \frac{(3.34\times 10^{-27})(9.8\times 10^{5})^{2}}
{2(1.6\times 10^{-19})(5)}\frac{1}{6}
\approx 3.3\times 10^{2}\,\si{\meter\per\second}.
\]
The curvature drift has the same form with $v_{\parallel}^{2}$ in place of
$\tfrac12 v_{\perp}^{2}$, so it is about twice the gradient drift here,
$\approx 6.7\times 10^{2}\,\si{\meter\per\second}$, and the two add because
both point vertically for a given charge. A drift of $\sim 10^{3}\,\si{
\meter\per\second}$ across a $2\,\si{\meter}$ plasma would empty it in a few
milliseconds if the helical twist did not cancel it, which sets the
quantitative requirement on the poloidal field: the field line must wind
poloidally fast enough that a particle samples the top and bottom of the
plasma before its vertical drift carries it to the wall.

The \engterm{magnetic mirror} effect arises from conservation of the
magnetic moment $\mu = mv_{\perp}^{2}/(2B)$ along a field line of
varying intensity. As a charged particle moves into a region of
higher $B$, its perpendicular kinetic energy increases, slowing its
parallel motion; the particle reflects if $v_{\parallel}^{2}/v^{2} <
1 - B_{\min}/B_{\max}$. Magnetic mirrors are the basis of the early
fusion machines (the mirror at Lawrence Livermore) and of the
geomagnetic trapping that produces the Earth's Van Allen radiation
belts. The condition $v_{\parallel}^{2}/v^{2} < 1 - B_{\min}/B_{\max}$
defines a loss cone in velocity space: particles whose velocity vector
points too nearly along the field escape through the throat, while
particles with enough perpendicular energy reflect. The unavoidable
scattering of particles into the loss cone is the reason a simple mirror
leaks and never became a fusion reactor on its own.

% Figure: magnetic mirror and loss cone. Left: field lines pinching at two
% throats with a trapped particle bouncing between mirror points. Right: the
% velocity-space loss cone whose half-angle is set by the mirror ratio. The
% mirror confinement worked example carries the numbers.
\begin{figure}[htbp]
\centering
\begin{tikzpicture}[scale=1.0]
  % === left panel: mirror geometry ===
  % converging field lines (upper and lower), pinched at x=0 and x=6
  \draw[engblue, thick] (0,0.5) .. controls (1.5,1.4) and (4.5,1.4) .. (6,0.5);
  \draw[engblue, thick] (0,-0.5) .. controls (1.5,-1.4) and (4.5,-1.4) .. (6,-0.5);
  \draw[engblue, thick] (0,0.25) .. controls (1.5,0.9) and (4.5,0.9) .. (6,0.25);
  \draw[engblue, thick] (0,-0.25) .. controls (1.5,-0.9) and (4.5,-0.9) .. (6,-0.25);
  % coils at the throats
  \fill[engred] (0,0.55) circle (0.07); \fill[engred] (0,-0.55) circle (0.07);
  \fill[engred] (6,0.55) circle (0.07); \fill[engred] (6,-0.55) circle (0.07);
  \node[engred, font=\scriptsize, anchor=south] at (0,0.7) {$B_{\max}$};
  \node[engred, font=\scriptsize, anchor=south] at (6,0.7) {$B_{\max}$};
  \node[enggray, font=\scriptsize, anchor=south] at (3,1.0) {$B_{\min}$};
  % trapped particle bounce path
  \draw[engorange, ->, thick] (1.2,0) .. controls (3,0.55) .. (4.8,0)
    .. controls (3,-0.55) .. (1.3,0.02);
  \node[engorange, font=\scriptsize, anchor=north] at (3,-0.6) {trapped, bouncing};

  % === right panel: velocity-space loss cone ===
  \begin{scope}[shift={(9.0,0)}]
    \draw[->, enggray] (-1.7,0) -- (1.9,0)
      node[right, font=\scriptsize] {$v_\parallel$};
    \draw[->, enggray] (0,-1.5) -- (0,1.7)
      node[above, font=\scriptsize] {$v_\perp$};
    % loss cones (half-angle ~ 35 deg)
    \fill[engred!18] (0,0) -- (1.5,1.05) -- (1.5,-1.05) -- cycle;
    \fill[engred!18] (0,0) -- (-1.5,1.05) -- (-1.5,-1.05) -- cycle;
    \draw[engred] (0,0) -- (1.5,1.05); \draw[engred] (0,0) -- (1.5,-1.05);
    \draw[engred] (0,0) -- (-1.5,1.05); \draw[engred] (0,0) -- (-1.5,-1.05);
    \node[engred, font=\scriptsize, anchor=west] at (0.7,0.15) {loss cone};
    \draw[engblue, thick] (0,0) circle (1.3);
    \node[engblue, font=\scriptsize, anchor=south] at (0,1.32) {trapped};
  \end{scope}
\end{tikzpicture}
\caption[Magnetic mirror and loss cone]{The magnetic-mirror confinement
scheme. Left: field lines converge to a higher intensity $B_{\max}$ at each
throat, with $B_{\min}$ at the centre. Conservation of the magnetic moment
$\mu=mv_\perp^{2}/(2B)$ converts parallel into perpendicular energy as a
particle moves into stronger field, reflecting it before it reaches the
throat, so the particle bounces between the two mirror points. Right: in
velocity space the reflection condition $v_\parallel^{2}/v^{2}<1-B_{\min}/
B_{\max}$ traps particles outside a loss cone whose half-angle is set by the
mirror ratio $R_m=B_{\max}/B_{\min}$. Particles scattered into the cone
stream out through the throat, and that unavoidable leak is why a simple
mirror never closed the fusion balance on its own.}
\label{fig:vol04ch13:magnetic-mirror}
\end{figure}


Figure~\ref{fig:vol04ch13:magnetic-mirror} pairs the field geometry with
the velocity-space picture the geometry implies. The mirror ratio $R_m =
B_{\max}/B_{\min}$ is the one number that sets confinement: a larger ratio
opens a wider band of trapped pitch angles and narrows the loss cone. The
loss-cone half-angle $\theta_{\mathrm{lc}}$ follows from the reflection
condition with the equality $\sin^{2}\theta_{\mathrm{lc}} = B_{\min}/
B_{\max} = 1/R_m$, so the fraction of an isotropic velocity distribution
that sits inside the cone and escapes is fixed by $R_m$ alone, independent
of energy.

\paragraph{Worked example: mirror ratio and confined fraction.}
A mirror coil set produces $B_{\min} = 0.5\,\si{\tesla}$ at the centre and
$B_{\max} = 4\,\si{\tesla}$ at each throat, a mirror ratio $R_m = 8$. The
loss-cone half-angle is $\theta_{\mathrm{lc}} = \arcsin(\sqrt{1/R_m}) =
\arcsin(\sqrt{1/8}) = \arcsin(0.354) \approx 20.7^{\circ}$. The fraction of
an isotropic distribution lost through both throats is the solid angle of
the two cones over the full sphere,
\[
f_{\mathrm{loss}} = 1 - \cos\theta_{\mathrm{lc}}
= 1 - \cos(20.7^{\circ}) \approx 1 - 0.936 = 0.064,
\]
so about $6\%$ of particles sit in the loss cone at any instant. That number
understates the problem: Coulomb collisions continually scatter trapped
particles into the empty cone, and the cone refills on the collision time,
so the confinement time of a mirror is set by the scattering time rather
than by the instantaneous loss fraction. Raising $R_m$ from $8$ to $20$
shrinks the half-angle only to $\arcsin(\sqrt{0.05}) \approx 12.9^{\circ}$
and the loss fraction to $\approx 0.025$, a modest gain for a large jump in
peak field. The weak return on mirror ratio is one reason the mirror line
gave way to closed toroidal confinement.

\paragraph{Worked example: cyclotron radiation from a magnetised electron.}
A gyrating electron is an accelerated charge, so it radiates, and in a strong
field the loss is large enough to enter the fusion power balance. The
single-electron cyclotron power is $P = \mu_0 e^{2}\omega_c^{2}v_\perp^{2}/(6\pi
c)$, the Larmor formula with the centripetal acceleration $a = \omega_c
v_\perp$. For a fusion electron at $T_e = 10\,\si{\kilo\electronvolt}$ in a
$B = 5\,\si{\tesla}$ field, $v_\perp \approx 5.9\times 10^{7}\,\si{\meter\per
\second}$ and $\omega_c = eB/m_e \approx 8.8\times 10^{11}\,\si{\radian\per
\second}$, so the per-electron power is $P = (4\pi\times 10^{-7})(1.6\times
10^{-19})^{2}(8.8\times 10^{11})^{2}(5.9\times 10^{7})^{2}/[6\pi(3\times
10^{8})] \approx 1.5\times 10^{-14}\,\si{\watt}$. Multiplied by the electron
density $10^{20}\,\si{\per\meter\cubed}$ this is $\sim 1.5\times 10^{6}\,\si{\watt
\per\meter\cubed}$, of the same order as the fusion power density, so the
cyclotron loss would be crippling were it not that the plasma is optically thick
to most of its own cyclotron radiation and reabsorbs it. The emission scales as
$B^{2}T_e$, so it climbs faster than fusion power with field and temperature and
sets a soft ceiling on how hot and how strongly magnetised a reactor can run
before the escaping cyclotron radiation, at the harmonics where the plasma turns
transparent, drains the energy balance. The same emission, read as a signal
rather than a loss, is the basis of electron-cyclotron-emission thermometry,
which reads $T_e$ from the intensity at each cyclotron harmonic across the
plasma.

\paragraph{Worked example: the banana orbit width of a trapped particle.}
The mirror physics of a straight field reappears inside a torus, because the
field strength varies as $1/R$ along a field line and traps the particles with
small parallel velocity between the strong-field inner points. A trapped
particle does not follow the flux surface; it drifts off it by a banana-shaped
excursion whose width is the step size of the neoclassical transport of the
later mastery box. The banana width is the trapped-particle bounce averaged
against the vertical drift, and it works out to $w_b \approx q\,r_L/\sqrt{\epsilon}$,
where $q$ is the safety factor, $r_L$ the poloidal-field Larmor radius, and
$\epsilon = r/R$ the inverse aspect ratio. Take a deuteron at $T = 10\,\si{\kilo
\electronvolt}$ in the case-study tokamak, $R = 6\,\si{\meter}$, $r = 1\,\si{
\meter}$ so $\epsilon = 1/6$, with a poloidal field $B_\theta \approx 0.6\,\si{
\tesla}$ and safety factor $q \approx 2$. The poloidal Larmor radius is $r_L =
m_D v_\perp/(eB_\theta) = (3.34\times 10^{-27})(9.8\times 10^{5})/[(1.6\times
10^{-19})(0.6)] \approx 3.4\times 10^{-2}\,\si{\meter}$, so the banana width is
$w_b \approx (2)(3.4\times 10^{-2})/\sqrt{1/6} \approx 0.17\,\si{\meter}$, about
seventeen centimetres, far larger than the few-millimetre gyroradius on the full
field. A collision that scatters a trapped particle steps it across the field by
this banana width rather than by the gyroradius, which is why neoclassical
transport is enhanced by orders of magnitude over the naive classical estimate,
and why the trapped-particle population, not the passing one, dominates
cross-field diffusion in a torus.

\paragraph{Worked example: ion-cyclotron resonance heating frequency.}
A magnetised plasma absorbs a wave efficiently when the wave frequency matches a
species' cyclotron frequency, and the ion-cyclotron branch heats the ions
directly, unlike the electron-cyclotron scheme of the later worked example. The
resonance condition is $\omega = \omega_{ci} = ZeB/m_i$. For the case-study
tokamak field $B = 5.3\,\si{\tesla}$ and a deuterium majority, $\omega_{ci} =
(1.6\times 10^{-19})(5.3)/(3.34\times 10^{-27}) \approx 2.5\times 10^{8}\,\si{
\radian\per\second}$, a frequency $f_{ci} = \omega_{ci}/2\pi \approx 40\,\si{
\mega\hertz}$, in the radio band a conventional transmitter reaches. Heating the
deuterium majority directly at this frequency couples poorly, so the standard
scheme adds a small hydrogen or helium-3 minority whose cyclotron frequency
falls at a chosen radius and dumps the launched power into the minority ions,
which then share it with the bulk by collisions. The resonance is a vertical
surface in the torus, the locus where $B$ takes the value that tunes $\omega_{ci}$
to the transmitter, so aiming the resonance is a matter of choosing the
frequency for the field at the target radius, exactly as the electron-cyclotron
worked example aimed its gyrotron. The megahertz ion-cyclotron band and the
tens-of-gigahertz electron-cyclotron band bracket the radio-frequency heating
toolkit a tokamak carries.

\begin{mastery}
The magnetic moment $\mu = mv_{\perp}^{2}/(2B)$ is the first
\engterm{adiabatic invariant}: it is conserved when the field changes
slowly compared to the gyration period, $\frac{1}{\omega_{c}}\left|
\frac{1}{B}\frac{\dd B}{\dd t}\right| \ll 1$. Two further adiabatic
invariants govern slower motions: the longitudinal invariant $J = \oint
v_{\parallel}\,\dd\ell$ for a particle bouncing between mirror points, and
the flux invariant $\Phi$ enclosed by the slow drift of the guiding centre
around a closed device. The three invariants order the timescales of
magnetised-plasma motion, and the controlled violation of each is how
energy and particles are transported across field lines. In the Van Allen
belts the first two invariants are well conserved over many bounce
periods, which is why energetic particles remain trapped for months;
violation of the third invariant by magnetic storms is what fills and
empties the belts.
\end{mastery}

\section{Fluid models of plasma: MHD}\pathtag{core}

The kinetic description of plasma (Vlasov equation for each species)
is necessary at the deepest level but is operationally cumbersome.
The \engterm{magnetohydrodynamic} (MHD) approximation treats the
plasma as a single conducting fluid with mass density $\rho$,
velocity $\vect{v}$, pressure $p$, current density $\vect{J}$, and
electric and magnetic fields satisfying Maxwell's equations. The
ideal-MHD equations are
\begin{align}
\pdv{\rho}{t} + \divg(\rho\vect{v}) &= 0,\\
\rho\dv{\vect{v}}{t} &= -\grad p + \vect{J}\times \vect{B},\\
\vect{E} + \vect{v}\times \vect{B} &= 0,\\
\pdv{\vect{B}}{t} &= \curl(\vect{v}\times \vect{B}),\\
\curl \vect{B} &= \mu_{0}\vect{J}.
\end{align}
The first equation is mass continuity. The second is Newton's second
law with the $\vect{J}\times \vect{B}$ force on the plasma; this
force is perpendicular to both $\vect{J}$ and $\vect{B}$ and is the
core of plasma confinement in magnetic devices. The third is the
ideal-MHD Ohm's law (zero resistivity); the fourth is its consequence,
the \engterm{frozen-in flux theorem}: magnetic field lines are tied
to the fluid and move with it. The fifth is Amp\`ere's law in the
slow-flow limit.

Two operating regimes follow:
\begin{itemize}[leftmargin=*]
  \item \engterm{Magnetic pressure} $p_{B} = B^{2}/(2\mu_{0})$.
    Compares directly to the plasma's thermal pressure
    $p = nk_{B}T$. The ratio $\beta = p/p_{B}$ measures whether the
    plasma is magnetically dominated ($\beta\ll 1$, the working
    regime of fusion devices and the solar corona) or
    pressure-dominated ($\beta\gtrsim 1$, the working regime of the
    solar wind and the chromosphere).
  \item \engterm{Alfv\'en speed} $v_{A} = B/\sqrt{\mu_{0}\rho}$. The
    propagation speed of magnetic-tension waves along field lines.
    At ITER conditions ($B = 5\,\si{\tesla}$, $\rho \sim 10^{-6}\,\si{
    \kilogram\per\meter\cubed}$), $v_{A}\approx 4\times 10^{6}\,\si{
    \meter\per\second}$.
\end{itemize}

The frozen-in flux theorem has dramatic operational consequences. In
a perfectly conducting plasma, magnetic-field topology is preserved
on convective time scales: field lines cannot reconnect or break. In
real plasmas the small but finite resistivity allows reconnection on
the resistive time scale $\tau_{R} = \mu_{0}L^{2}/\eta$, where $\eta$
is resistivity and $L$ is the characteristic length. The Spitzer
resistivity scales as $T_{e}^{-3/2}$; hot plasmas are extraordinarily
good conductors. Magnetic reconnection in the solar corona ($\tau_{R}
\sim 10^{12}\,\si{\second}$ by Spitzer alone, but $\sim 100\,\si{
\second}$ observationally) is one of the open problems in plasma
astrophysics, with implications for solar flares and coronal heating.

\paragraph{Worked example: Spitzer resistivity and ohmic heating.}
The Spitzer parallel resistivity of a hydrogenic plasma is
\[
\eta_{\parallel} \approx 5.2\times 10^{-5}\,
\frac{Z\ln\Lambda}{[T_e/\text{eV}]^{3/2}}\ \si{\ohm\meter},
\]
with $Z$ the ion charge and $\ln\Lambda$ the Coulomb logarithm, typically
$15$ to $20$ in fusion plasmas. At $T_e = 1\,\si{\kilo\electronvolt} =
10^{3}\,\text{eV}$, $Z = 1$, $\ln\Lambda = 17$, the resistivity is
$\eta_{\parallel} \approx 5.2\times 10^{-5}(17)/(10^{3})^{3/2} \approx
2.8\times 10^{-8}\,\si{\ohm\meter}$, close to that of copper at room
temperature. At $T_e = 10\,\si{\kilo\electronvolt}$ it falls by $10^{3/2}
\approx 32$ to $\sim 9\times 10^{-10}\,\si{\ohm\meter}$, an order of
magnitude better than copper, which is the operational meaning of the
$T_e^{-3/2}$ scaling: a hot plasma is an excellent conductor. The ohmic
heating power density is $p_{\Omega} = \eta_{\parallel}J^{2}$; for a tokamak
running a current density $J = 1\,\si{\mega\ampere\per\meter\squared}$ at
$1\,\si{\kilo\electronvolt}$, $p_{\Omega} = (2.8\times 10^{-8})(10^{6})^{2}
\approx 2.8\times 10^{4}\,\si{\watt\per\meter\cubed}$, a useful start-up
heat. The same scaling defeats ohmic heating as the path to ignition: once
the plasma reaches a few $\si{\kilo\electronvolt}$ the resistivity has
fallen so far that the deposited power can no longer climb the temperature
toward the $10\,\si{\kilo\electronvolt}$ fusion window, and auxiliary
heating (neutral beams, radio-frequency power) must take over.

\paragraph{Worked example: collision frequency and the mean free path.}
The resistivity above hides a collision frequency an engineer often needs
directly. The electron-ion collision frequency of a hydrogenic plasma is
$\nu_{ei} \approx 2.9\times 10^{-12}\,n_e\ln\Lambda\,[T_e/\text{eV}]^{-3/2}\,
\si{\per\second}$ with $n_e$ in $\si{\per\meter\cubed}$. For the case-study tokamak at $n_e = 10^{20}\,\si{\per\meter
\cubed}$, $T_e = 12\,\si{\kilo\electronvolt} = 1.2\times 10^{4}\,\text{eV}$,
and $\ln\Lambda = 17$, this is $\nu_{ei} \approx 2.9\times 10^{-12}(10^{20})(17)
/(1.2\times 10^{4})^{3/2} \approx 3.8\times 10^{3}\,\si{\per\second}$, so an
electron collides only a few thousand times a second. The thermal electron
speed is $v_{te} = \sqrt{k_B T_e/m_e} = \sqrt{(1.92\times 10^{-15})/(9.11\times
10^{-31})} \approx 4.6\times 10^{7}\,\si{\meter\per\second}$, so the mean free
path is $\lambda_{\mathrm{mfp}} = v_{te}/\nu_{ei} \approx 4.6\times 10^{7}/
3.8\times 10^{3} \approx 1.2\times 10^{4}\,\si{\meter}$, twelve kilometres,
thousands of times the machine size. A fusion-core electron circles the torus
hundreds of times between collisions, which is why the plasma is collisionless
on the confinement scale and the parallel transport is set by the field
geometry rather than by local collisions. The same formula at the fluorescent
lamp, $n_e = 10^{17}$, $T_e = 1\,\text{eV}$, $\ln\Lambda = 10$, gives
$\nu_{ei} \approx 2.9\times 10^{-12}(10^{17})(10)/1 \approx 2.9\times
10^{6}\,\si{\per\second}$ and, with $v_{te} \approx 4.2\times 10^{5}\,\si{
\meter\per\second}$, an electron-ion mean free path of $\lambda_{\mathrm{mfp}}
\approx 0.14\,\si{\meter}$, comparable to the tube. Coulomb collisions alone
are marginal at that scale; what makes the lamp collisional is the weakly
ionised gas, where each electron strikes a neutral argon or mercury atom
every millimetre or so. Those electron-neutral collisions dominate the
transport, a fluid description with local coefficients applies, and the lamp
sits at the opposite extreme from the collisionless core. The single group
$n_e T_e^{-3/2}$ carries the electron-ion collision rate from one regime to
the other.

% Figure: a representative safety-factor profile q(r) across the minor radius
% of a tokamak, rising from a central value near one to an edge value of three
% to four. Rational surfaces where q is a low-order fraction are marked, since
% these are where tearing and kink modes resonate. The Kruskal-Shafranov
% stability line q_edge=1 is drawn as the floor the profile must stay above.
\begin{figure}[htbp]
\centering
\begin{tikzpicture}
\begin{axis}[
  width=0.82\textwidth, height=0.55\textwidth,
  xlabel={normalised minor radius $r/a$},
  ylabel={safety factor $q(r)$},
  xmin=0, xmax=1, ymin=0, ymax=4.5,
  grid=both, grid style={enggray!20},
  every axis plot/.append style={thick},
]
  % monotonic q profile rising from ~0.9 on axis to ~3.8 at edge
  \addplot[engblue, domain=0:1, samples=80]
    {0.9 + 2.9*x*x};
  % Kruskal-Shafranov floor q=1
  \draw[dashed, engred] (axis cs:0,1) -- (axis cs:1,1);
  \node[engred, font=\scriptsize, anchor=south west] at (axis cs:0,1)
    {Kruskal-Shafranov floor $q=1$};
  % rational surfaces q=2 and q=3
  \draw[dotted, enggray] (axis cs:0,2) -- (axis cs:0.62,2);
  \draw[dotted, enggray] (axis cs:0.62,0) -- (axis cs:0.62,2);
  \node[enggray, font=\scriptsize, anchor=south] at (axis cs:0.62,2.05)
    {$q=2$ surface};
  \draw[dotted, enggray] (axis cs:0,3) -- (axis cs:0.84,3);
  \draw[dotted, enggray] (axis cs:0.84,0) -- (axis cs:0.84,3);
  \node[enggray, font=\scriptsize, anchor=south] at (axis cs:0.84,3.05)
    {$q=3$ surface};
\end{axis}
\end{tikzpicture}
\caption[Safety-factor profile and rational surfaces]{A representative
safety-factor profile across the minor radius of a tokamak. The safety factor
$q$ counts how many times a field line wraps the long way around the torus for
each time it wraps the short way; a large $q$ means a gentle helical pitch and
a stiff, stable column. The profile must stay above the Kruskal-Shafranov
floor $q=1$ to avoid the fast current-driven kink, and it rises to an edge
value of three to four in normal operation. The surfaces where $q$ passes
through a low-order fraction, drawn here at $q=2$ and $q=3$, are the rational
surfaces where a field line closes on itself after a few transits; these are
the resonant locations at which tearing modes grow magnetic islands and where
a disruption is most often seeded. A designer shapes the current profile to
keep $q$ above one on axis and to place the rational surfaces where they do
the least harm.}
\label{fig:vol04ch13:safety-factor}
\end{figure}


\paragraph{Worked example: the safety factor and the kink stability limit.}
The single number that decides whether a tokamak column holds together is the
\engterm{safety factor} $q$, the number of times a field line wraps the long way
around the torus for each time it wraps the short way. For a large-aspect-ratio
circular plasma, $q(r) = r B_\phi/(R B_\theta)$, the ratio of the toroidal field
$B_\phi$ to the poloidal field $B_\theta$ scaled by the geometry. Take the
case-study tokamak at its edge, $r = a = 2\,\si{\meter}$, $R = 6\,\si{\meter}$,
toroidal field $B_\phi = 5.3\,\si{\tesla}$, carrying a plasma current $I_p =
15\,\si{\mega\ampere}$. The edge poloidal field from Amp\`ere's law is
$B_\theta = \mu_0 I_p/(2\pi a) = (4\pi\times 10^{-7})(1.5\times 10^{7})/[2\pi
(2)] \approx 1.5\,\si{\tesla}$, so the edge safety factor is $q_a = a B_\phi/
(R B_\theta) = (2)(5.3)/[(6)(1.5)] \approx 1.2$. The Kruskal-Shafranov limit
requires $q_a > 1$ to hold off the fast current-driven kink, and the profile of
Figure~\ref{fig:vol04ch13:safety-factor} must clear that floor everywhere while
keeping $q$ above one on axis. The edge value near $1.2$ sits uncomfortably
close to the limit, which is the quantitative statement of why a tokamak cannot
be pushed to arbitrarily high current: raising $I_p$ raises $B_\theta$ and drops
$q_a$ toward one, and once the edge safety factor crosses the Kruskal-Shafranov
floor the column kinks and disrupts. The current limit and the beta limit
together bound the operating space, and a designer works the plasma at the
highest current that keeps $q_a$ safely above two in normal operation.

% Figure: the friction force on an electron against its speed, showing the
% collisional drag rising then falling as 1/v^2, crossed by the constant
% accelerating force of an applied electric field. Where the applied force
% exceeds the peak drag, an electron above the crossing speed accelerates
% without bound: it runs away. The Dreicer field is the field at which the
% peak drag equals the applied force and the whole distribution runs away.
\begin{figure}[htbp]
\centering
\begin{tikzpicture}
\begin{axis}[
  width=0.82\textwidth, height=0.55\textwidth,
  xlabel={electron speed $v/v_{te}$},
  ylabel={force on electron (arb.\ units)},
  xmin=0.2, xmax=5, ymin=0, ymax=1.2,
  grid=both, grid style={enggray!20},
  legend pos=north east,
  legend cell align={left},
  every axis plot/.append style={thick},
]
  % collisional drag: rises to a peak near v_te then falls as 1/v^2
  % model F_drag = A v / (1 + v^3) scaled; peak near v ~ 1
  \addplot[engblue, domain=0.2:5, samples=120]
    {1.4*x/(0.3 + x*x*x)};
  \addlegendentry{collisional drag}
  % applied field force below the peak: partial runaway above crossing
  \addplot[engorange, domain=0.2:5, samples=2]
    {0.45};
  \addlegendentry{applied field $eE$}
  % crossing / critical speed marker
  \addplot[only marks, mark=*, engred, mark size=2pt] coordinates {(2.05,0.45)};
  \node[engred, font=\scriptsize, anchor=south west] at (axis cs:2.05,0.45)
    {critical speed $v_c$};
  \node[enggray, font=\scriptsize, anchor=south] at (axis cs:1.0,1.02)
    {peak drag};
\end{axis}
\end{tikzpicture}
\caption[Runaway electrons and the critical field]{The collisional drag on an
electron rises with speed to a peak near the thermal speed and then falls off
as the inverse square of the speed, because a faster electron spends less time
near each scattering centre. A steady electric field exerts a constant force
$eE$, drawn as the horizontal line. An electron faster than the crossing speed
$v_c$, where the drag has fallen below the applied force, is accelerated
faster than it is slowed and runs away to relativistic energy. Lowering the
field raises $v_c$ and shrinks the runaway population; raising the field toward
the Dreicer value, where the applied force meets the peak of the drag curve,
lets the entire distribution run away. In a tokamak disruption the loop
voltage induced by the collapsing current drives exactly this runaway, and a
multi-mega-electronvolt electron beam that can drill through the vessel wall is
the result the mitigation system must prevent.}
\label{fig:vol04ch13:runaway}
\end{figure}


\paragraph{Worked example: the Dreicer field and runaway electrons.}
The falling drag of a fast electron, Figure~\ref{fig:vol04ch13:runaway}, means a
strong enough electric field accelerates part of the distribution without bound.
The threshold is the \engterm{Dreicer field}, the field at which the accelerating
force equals the collisional drag on a thermal electron,
\[
E_D = \frac{n_e e^{3}\ln\Lambda}{4\pi\varepsilon_0^{2}k_B T_e}.
\]
For the case-study tokamak at $n_e = 10^{20}\,\si{\per\meter\cubed}$, $T_e =
12\,\si{\kilo\electronvolt} = 1.92\times 10^{-15}\,\si{\joule}$, and $\ln\Lambda
= 17$, the numerator is $(10^{20})(1.6\times 10^{-19})^{3}(17) \approx 6.96\times
10^{-36}$ and the denominator is $4\pi(8.85\times 10^{-12})^{2}(1.92\times
10^{-15}) \approx 1.89\times 10^{-36}$, so $E_D \approx 3.7\,\si{
\volt\per\meter}$. The loop voltage that drives the normal plasma current is a
few volts around the $\sim 38\,\si{\meter}$ circumference, an electric field of
order $0.1\,\si{\volt\per\meter}$, far below the Dreicer field, so ordinary
operation makes only a negligible runaway population. In a disruption, however,
the current quenches in milliseconds and the induced loop voltage spikes by
orders of magnitude, driving the field toward $E_D$; the fraction of electrons
above the critical speed then runs away to tens of mega-electronvolts and forms
a beam that carries a large part of the pre-disruption current. That beam, if it
strikes the wall in a narrow spot, drills deep local damage, which is why
runaway avoidance and dissipation are a named requirement of the disruption
mitigation system rather than a secondary concern.

\begin{mastery}
The dimensionless number that governs reconnection is the
\engterm{Lundquist number} $S = \mu_{0}L v_{A}/\eta$, the ratio of the
resistive diffusion time to the Alfv\'en transit time. The Sweet-Parker
model predicts a reconnection rate scaling as $S^{-1/2}$, which is far too
slow to explain observed solar flares at $S \sim 10^{12}$. The resolution,
established over decades of plasma-astrophysics work, is that thin current
sheets become unstable to the plasmoid instability above $S \sim 10^{4}$,
breaking into chains of magnetic islands that reconnect at a rate nearly
independent of $S$. The practical lesson for the engineer is that the
ideal-MHD frozen-in picture is an approximation that fails precisely where
it matters most, in thin current sheets where field topology is rearranged
and stored magnetic energy is released suddenly. The same physics underlies
the sawtooth crash in tokamaks, in which the core field reconnects on a
fast timescale and flattens the central temperature profile.
\end{mastery}

\section{Foundations: the equations behind the working formulas}\pathtag{core}

The operational formulas used so far were quoted with a sketch of their
origin. An engineer who must trust them at the margin, where a $20\%$
confinement shortfall decides ignition or a factor-of-two ion mass moves an
inferred density, needs the derivations in full. Four results carry most of
the weight: the Saha equation that fixes the ionisation, the decomposition
of the magnetic force into a pressure and a tension that gives the beta
limit its meaning, the reduction of the two-species kinetic plasma to the
single conducting fluid of MHD, and the frozen-in flux theorem that makes
field topology a conserved quantity. We derive each from its starting point
and read off the engineering consequence.

\subsection{The Saha equation from statistical mechanics}

The Saha equation is the law of mass action applied to the ionisation
reaction $\mathrm{A} \rightleftharpoons \mathrm{A}^{+} + e^{-}$ in thermal
equilibrium. The derivation is the chemical-equilibrium argument with the
electron treated as one of the products, and it shows where every factor in
the quoted formula comes from.

In thermal and chemical equilibrium the chemical potentials balance,
$\mu_{n} = \mu_{i} + \mu_{e}$, where $\mu_{n}$, $\mu_{i}$, $\mu_{e}$ are the
chemical potentials of the neutral atom, the ion, and the electron. For an
ideal gas of structureless particles of mass $m$ and internal partition
function $g$, the chemical potential is
\[
\mu = -k_{B}T\ln\!\left(\frac{g}{n}\left(
\frac{m k_{B}T}{2\pi\hbar^{2}}\right)^{3/2}\right)
+ \varepsilon_{0},
\]
where $n$ is the number density, the bracketed factor is the inverse cube of
the thermal de Broglie wavelength $\lambda_{\mathrm{th}} =
(2\pi\hbar^{2}/m k_{B}T)^{1/2}$, and $\varepsilon_{0}$ is the rest energy of
the internal ground state. The thermal wavelength is the quantum length over
which a particle's wavefunction is spread at temperature $T$; the quantity
$n\lambda_{\mathrm{th}}^{3}$ measures whether the gas is classical
($n\lambda_{\mathrm{th}}^{3}\ll 1$) or degenerate.

Write the equilibrium condition with the rest energies measured from the
neutral ground state, so the ion-plus-electron state lies an ionisation
energy $E_{i}$ above the neutral. Equating the chemical potentials and
collecting terms gives the law of mass action
\[
\frac{n_{i}n_{e}}{n_{n}} = \frac{g_{i}g_{e}}{g_{n}}
\left(\frac{m_{e}k_{B}T}{2\pi\hbar^{2}}\right)^{3/2}
\exp\!\left(-\frac{E_{i}}{k_{B}T}\right).
\]
The electron mass appears in the prefactor, not the ion mass, because the
heavy ion and the neutral have nearly equal thermal wavelengths and their
ratio cancels; the surviving factor is the electron's thermal wavelength
cubed. The statistical weights $g_{i}$, $g_{e} = 2$ (two electron spin
states), and $g_{n}$ are the degeneracies of the internal states; for
hydrogen, taking $g_{i}/g_{n} = 1/2$ leaves the prefactor in the form quoted
at the start of the chapter, $g_{e}g_{i}/g_{n} = 1$. The exponential is the
Boltzmann factor for paying the ionisation energy out of the thermal bath.

Two engineering lessons follow from the structure rather than the value. The
prefactor has dimensions of density and is enormous at any laboratory
temperature, of order $10^{27}\,\si{\per\meter\cubed}$ at $10^{4}\,\si{
\kelvin}$, so the ionisation runs to substantial fractions long before the
mean thermal energy $k_{B}T$ reaches $E_{i}$. The transition is sharp
because the right-hand side multiplies a slowly varying power of $T$ by a
fast exponential; the half-ionisation temperature solves $n_{e} \approx
n_{n}$ and lands near $E_{i}/k_{B}T \approx 10$, an order of magnitude below
the naive $E_{i}/k_{B}T = 1$. The density dependence is hidden in the ratio:
at fixed temperature, raising the total density pushes the equilibrium back
toward the neutral side, since recombination scales as $n_{i}n_{e}$ while
ionisation scales as $n_{n}$, so a denser gas ionises at a higher
temperature. The companion \texttt{code/saha\_ionization.py} solves the
quadratic for the ionisation fraction at any density and reproduces the
curve of Figure~\ref{fig:vol04ch13:saha-fraction}.

\paragraph{Worked example: half-ionisation temperature of hydrogen.}
Solve the Saha equation for the temperature at which hydrogen is half
ionised at a total density $n_{\mathrm{tot}} = n_{n} + n_{i} = 10^{20}\,\si{
\per\meter\cubed}$, a representative low-density laboratory plasma. At half
ionisation $n_{i} = n_{e} = n_{n} = n_{\mathrm{tot}}/2 = 5\times 10^{19}\,
\si{\per\meter\cubed}$, so the Saha ratio is $n_{i}n_{e}/n_{n} = n_{e} =
5\times 10^{19}\,\si{\per\meter\cubed}$. Setting the right-hand side equal to
this and writing $\theta = E_{i}/k_{B}T$ with $E_{i} = 13.6\,\text{eV}$,
\[
5\times 10^{19} = \left(\frac{m_{e}k_{B}T}{2\pi\hbar^{2}}\right)^{3/2}
e^{-\theta}.
\]
The prefactor at $T$ near $10^{4}\,\si{\kelvin}$ is $\approx 3.0\times
10^{27}\,\si{\per\meter\cubed}$, so $e^{-\theta} \approx 5\times 10^{19}/
3\times 10^{27} = 1.7\times 10^{-8}$, giving $\theta = \ln(6\times 10^{7})
\approx 17.9$. Then $T = E_{i}/(k_{B}\theta) = (13.6)(1.6\times 10^{-19})/
[(1.38\times 10^{-23})(17.9)] \approx 8.8\times 10^{3}\,\si{\kelvin}$. The
gas is half ionised near $9000\,\si{\kelvin}$, far below the $1.6\times
10^{5}\,\si{\kelvin}$ that $E_{i}/k_{B}T = 1$ would give, which is the
quantitative content of the sharp Saha transition. A self-consistent
solution iterates the prefactor at the new temperature, but one pass already
lands within a few percent because the exponential dominates.

\subsection{The magnetic force as a pressure and a tension}

The force per unit volume on a plasma is $\vect{J}\times\vect{B}$. Using
Amp\`ere's law $\vect{J} = (\curl\vect{B})/\mu_{0}$ and the vector identity
$(\curl\vect{B})\times\vect{B} = (\vect{B}\cdot\grad)\vect{B} - \grad(B^{2}
/2)$, the force splits into two pieces,
\[
\vect{J}\times\vect{B}
= \frac{(\vect{B}\cdot\grad)\vect{B}}{\mu_{0}}
- \grad\!\left(\frac{B^{2}}{2\mu_{0}}\right).
\]
The second term is the gradient of the scalar $p_{B} = B^{2}/(2\mu_{0})$, the
\engterm{magnetic pressure}: it pushes plasma from strong-field to weak-field
regions exactly as a gas pressure pushes from high to low. The first term is
the \engterm{magnetic tension}: it is non-zero only where the field lines
curve. Writing $\vect{B} = B\hat{\vect{b}}$ and using $(\hat{\vect{b}}\cdot
\grad)\hat{\vect{b}} = \hat{\vect{n}}/R_{c}$, where $\hat{\vect{n}}$ points
toward the centre of curvature and $R_{c}$ is the radius of curvature, the
tension term becomes
\[
\frac{(\vect{B}\cdot\grad)\vect{B}}{\mu_{0}}
= \frac{B^{2}}{\mu_{0}R_{c}}\hat{\vect{n}}
+ \hat{\vect{b}}\,\frac{\partial}{\partial s}\!
\left(\frac{B^{2}}{2\mu_{0}}\right),
\]
where $s$ is arc length along the field line. The first piece is a force
$B^{2}/(\mu_{0}R_{c})$ per unit volume directed toward the centre of
curvature, the tension of a stretched elastic string of tension
$B^{2}/\mu_{0}$ per unit area; the second cancels the parallel part of the
pressure gradient, so the net force along the field comes only from the
plasma pressure. Figure~\ref{fig:vol04ch13:magnetic-stress} draws the two
contributions.

% Figure: the magnetic force decomposed into pressure and tension. A bundle
% of curved field lines exerts a pressure across the lines (B^2/2mu0) and a
% tension along them (B^2/mu0 over a radius of curvature R_c). The net
% J x B force is the sum, derived in the foundations section.
\begin{figure}[htbp]
\centering
\begin{tikzpicture}[>=Stealth, scale=1.0]
  % curved field lines (arcs), curving downward, centre of curvature below
  \foreach \off in {0,0.5,1.0,1.5} {
    \draw[engblue, thick]
      (0,\off) .. controls (2.5,{\off-0.7}) and (4.5,{\off-0.7}) .. (7,\off);
  }
  % arrowheads on the middle line
  \draw[engblue, thick, ->]
    (0,1.0) .. controls (2.5,0.3) and (4.5,0.3) .. (7,1.0);
  \node[engblue, anchor=west] at (7.1,1.0) {$\vect{B}$};

  % tension along the line: arrows pulling the apex toward straightening
  \draw[engred, very thick, ->] (3.5,0.43) -- (2.6,0.55);
  \draw[engred, very thick, ->] (3.5,0.43) -- (4.4,0.55);
  \node[engred, anchor=north] at (3.5,0.30)
    {tension $\dfrac{B^{2}}{\mu_{0}}$};

  % pressure across the lines: arrows pushing apart vertically
  \draw[engorange, very thick, ->] (1.2,0.6) -- (1.2,1.3);
  \draw[engorange, very thick, ->] (1.2,0.45) -- (1.2,-0.25);
  \node[engorange, anchor=west, font=\footnotesize] at (1.35,1.15)
    {pressure $\dfrac{B^{2}}{2\mu_{0}}$};

  % radius of curvature indication
  \draw[enggray, dashed, ->] (3.5,-2.0) -- (3.5,0.40);
  \node[enggray, anchor=east, font=\footnotesize] at (3.4,-1.0) {$R_{c}$};
  \fill[enggray] (3.5,-2.0) circle (1.2pt);
  \node[enggray, anchor=north, font=\footnotesize] at (3.5,-2.1)
    {centre of curvature};

  % net inward force note
  \draw[englime, line width=1.2pt, ->] (5.6,0.55) -- (5.6,-0.15);
  \node[englime, anchor=west, font=\footnotesize] at (5.7,0.2)
    {net $\vect{J}\times\vect{B}$};
\end{tikzpicture}
\caption[Magnetic pressure and tension]{The magnetic force on a plasma,
written as the divergence of the Maxwell stress, separates into a pressure
$B^{2}/2\mu_{0}$ acting across the field lines and a tension $B^{2}/\mu_{0}$
acting along them. Where the lines are straight the two combine to a pure
isotropic-minus-along pressure; where they curve with radius $R_{c}$ the
tension contributes a force $B^{2}/(\mu_{0}R_{c})$ directed toward the
centre of curvature. The curvature term is the restoring force of the
shear-Alfv\'en wave and the confining force that a toroidal field exerts on
a plasma column. The decomposition is derived in the foundations section
from $\vect{J}\times\vect{B} = (\curl\vect{B})\times\vect{B}/\mu_{0}$.}
\label{fig:vol04ch13:magnetic-stress}
\end{figure}


The decomposition is the engineer's tool for reading any magnetic
confinement geometry. A straight field line carries pressure but no tension,
so it confines only across itself; a curved field line adds a restoring
force that is the source of the shear-Alfv\'en wave, whose speed
$v_{A} = B/\sqrt{\mu_{0}\rho}$ is the propagation speed of a transverse
twang on a field-line string of tension $B^{2}/\mu_{0}$ and mass per length
$\rho$ times the cross-section, exactly the string-wave result $\sqrt{T/
\mu_{\mathrm{lin}}}$. The Alfv\'en speed quoted earlier is therefore not a
separate fact but the tension wave of this very decomposition.

The static force balance of a confined plasma is $\grad p = \vect{J}\times
\vect{B}$, which with the decomposition reads $\grad(p + B^{2}/2\mu_{0}) =
(\vect{B}\cdot\grad)\vect{B}/\mu_{0}$. Where the field is straight the
right-hand side vanishes and the total pressure $p + B^{2}/(2\mu_{0})$ is
constant across the plasma: the thermal pressure that the plasma carries is
bought from the magnetic pressure that the field gives up.
Figure~\ref{fig:vol04ch13:pressure-balance} shows the trade across a plasma
column.

% Figure: radial pressure balance in a confined plasma column. The thermal
% pressure p falls from the axis to the edge while the magnetic pressure
% B^2/2mu0 rises to compensate, so the total p + B^2/2mu0 is flat. Plasma
% beta is the ratio of the two at the axis.
\begin{figure}[htbp]
\centering
\begin{tikzpicture}
\begin{axis}[
  width=0.82\textwidth, height=0.50\textwidth,
  xlabel={normalised radius $r/a$},
  ylabel={pressure (arb.\ units)},
  xmin=0, xmax=1, ymin=0, ymax=1.25,
  grid=both, grid style={enggray!20},
  every axis plot/.append style={thick},
  legend pos=outer north east, legend cell align=left,
]
  % thermal pressure: high on axis, falls to ~0 at edge
  \addplot[engred, domain=0:1, samples=80] {0.7*(1 - x^2)};
  \addlegendentry{thermal $p$}
  % magnetic pressure: low on axis, rises to confine at edge
  \addplot[engblue, domain=0:1, samples=80] {0.3 + 0.7*x^2};
  \addlegendentry{magnetic $B^{2}/2\mu_{0}$}
  % total: flat (pressure balance)
  \addplot[englime, line width=1.2pt, domain=0:1, samples=80] {1.0};
  \addlegendentry{total (constant)}
  % beta marker on axis
  \draw[enggray, dashed] (axis cs:0,0.3) -- (axis cs:0.18,0.3);
  \draw[enggray, dashed] (axis cs:0,0.7) -- (axis cs:0.18,0.7);
  \node[enggray, font=\scriptsize, anchor=west]
    at (axis cs:0.05,0.5) {$\beta = p/p_{B}$};
\end{axis}
\end{tikzpicture}
\caption[Radial pressure balance and plasma beta]{Radial equilibrium of a
confined plasma column. The thermal pressure $p = nk_{B}T$ peaks on the
axis and falls to the edge, while the magnetic pressure $B^{2}/2\mu_{0}$
rises outward to compensate, so that the total pressure
$p + B^{2}/2\mu_{0}$ stays constant across the column, the statement of
radial force balance in a low-current screw pinch. The plasma beta is the
ratio of thermal to magnetic pressure evaluated on the axis; a fusion
device runs at a few percent beta (magnetically dominated), while the solar
wind and the chromosphere run at beta of order one or above. A higher beta
extracts more thermal pressure from a given field and so more fusion power
per unit of magnet cost, which is why raising the beta limit is a standing
goal of magnetic-confinement design.}
\label{fig:vol04ch13:pressure-balance}
\end{figure}


The ratio that names the trade is the \engterm{plasma beta},
\[
\beta = \frac{p}{B^{2}/(2\mu_{0})} = \frac{2\mu_{0}nk_{B}T}{B^{2}}.
\]
A device that confines a thermal pressure $p$ with a field $B$ runs at this
beta, and the beta the field can hold against the instabilities the pressure
drives is the figure of merit of a confinement scheme. Fusion power density
scales as $n^{2}\langle\sigma v\rangle \sim p^{2}/T^{2}\langle\sigma v
\rangle$ at fixed temperature, so it scales as $\beta^{2}B^{4}$: the fusion
power a magnet of a given field can produce rises as the fourth power of the
field and the square of the achievable beta, which is why both high-field
magnets and high-beta operation are pursued together. The Troyon scaling
sets an empirical beta limit of a few percent in a tokamak before the
pressure-driven ballooning modes break confinement, the same number the
case-study plasma sat at.

\paragraph{Worked example: beta and the field a plasma needs.}
The case-study tokamak ran at $n = 10^{20}\,\si{\per\meter\cubed}$,
$T = 12\,\si{\kilo\electronvolt}$, and $B = 5.3\,\si{\tesla}$, giving
$\beta \approx 0.034$. Ask the inverse question a designer faces: to hold
the same plasma at the Troyon-limited $\beta = 0.05$, how strong a field is
allowed at most? Since $\beta = 2\mu_{0}nk_{B}T/B^{2}$, holding $p$ fixed and
raising the allowed beta to $0.05$ lowers the required field to $B =
\sqrt{2\mu_{0}p/\beta}$. With $p = 2nk_{B}T = 2(10^{20})(1.92\times
10^{-15}) = 3.84\times 10^{5}\,\si{\pascal}$, $B = \sqrt{2(4\pi\times
10^{-7})(3.84\times 10^{5})/0.05} = \sqrt{1.93\times 10^{1}} \approx
4.4\,\si{\tesla}$. Running at the beta limit rather than below it lets the
same plasma be confined by a $4.4\,\si{\tesla}$ field instead of
$5.3\,\si{\tesla}$, a $30\%$ saving in $B^{2}$ and so in stored magnet
energy. The designer reads this as the payoff of operating near the beta
limit: every point of beta won is magnet cost saved or, at fixed field,
fusion power gained.

The beta a machine may run is not a single number but a boundary in a plane
of operating parameters, and reading that plane is the equilibrium engineer's
daily work. The Troyon scaling writes the achievable beta as $\beta_{\max}
[\%] = \beta_N I_p/(aB)$, tying the pressure the field can hold to the
normalised plasma current through a constant $\beta_N$ of order $2.8$ that the
stability calculations and the experiments agree on. The scaling has a
partner on the current side, since the same current that buys beta drops the
edge safety factor toward the Kruskal-Shafranov floor, and a third boundary
on the density side, the Greenwald limit, beyond which the edge cools and the
plasma disrupts. Figure~\ref{fig:vol04ch13:beta-troyon} draws the accessible
region as the wedge below all three limits. A design does not sit at the
centre of the wedge but pushes toward the high-current, high-beta corner where
the fusion power is largest, holding back only enough margin to survive the
scatter in the limits and the transients of a real discharge. The case-study
plasma of the later section sits inside this wedge at a beta of a few percent,
short of the Troyon corner, which is the room a designer keeps against the
confinement uncertainty that section closes on.

% Figure: the tokamak operating space in the plane of normalised plasma
% current against normalised beta, bounded on two sides. The Troyon beta
% limit caps the pressure the field can hold before ballooning modes break
% confinement; the Greenwald density limit and the Kruskal-Shafranov current
% limit cap the current. The accessible operating region is the wedge between
% the limits, and a design works the plasma at the corner that maximises
% fusion power while staying clear of every boundary.
\begin{figure}[htbp]
\centering
\begin{tikzpicture}
\begin{axis}[
  width=0.80\textwidth, height=0.56\textwidth,
  xlabel={normalised current $I_p/aB$ (\si{\mega\ampere}/\si{\meter}/\si{\tesla})},
  ylabel={plasma beta $\beta$ (\%)},
  xmin=0, xmax=3.0, ymin=0, ymax=6,
  grid=both, grid style={enggray!20},
  every axis plot/.append style={thick},
]
  % Troyon: beta_max (%) = beta_N * I_p/(a B), with beta_N ~ 2.8 (conservative 2 shown)
  \addplot[engblue, domain=0:2.6, samples=2] {2.0*x};
  % current limit (Kruskal-Shafranov / q>2) as a vertical line
  \addplot[engred, dashed, domain=0:6, samples=2] coordinates {(2.4,0)(2.4,6)};
  % beta ceiling from disruption experience
  \addplot[engamber, dotted, domain=0:3, samples=2] coordinates {(0,5)(3,5)};
  \node[engblue, font=\scriptsize, anchor=south east, rotate=63]
    at (axis cs:2.1,4.2) {Troyon limit $\beta=\beta_N I_p/aB$};
  \node[engred, font=\scriptsize, anchor=north, rotate=90]
    at (axis cs:2.4,3.0) {current limit $q_a>2$};
  \node[engamber, font=\scriptsize, anchor=south west]
    at (axis cs:0.1,5.05) {disruption ceiling};
  \node[font=\scriptsize, anchor=center] at (axis cs:1.5,1.5) {operating region};
  \addplot[only marks, mark=*, engblack, mark size=2pt] coordinates {(1.2,3.4)};
  \node[font=\scriptsize, anchor=west] at (axis cs:1.25,3.4) {case study};
\end{axis}
\end{tikzpicture}
\caption[Tokamak operating space and the beta limit]{The tokamak operating
space in the plane of normalised plasma current against plasma beta. The
Troyon scaling makes the achievable beta rise in proportion to the normalised
current $I_p/aB$ through a constant $\beta_N$ of a few percent, the sloped
line, so a designer buys beta by driving current. The current cannot rise
without bound: the Kruskal-Shafranov and safety-factor limits cap it where
the edge safety factor falls toward two, the vertical boundary, and
accumulated disruption experience sets a further ceiling on beta itself, the
horizontal boundary. The accessible region is the wedge below all three
limits, and the fusion power a machine can make is largest at the upper
corner, so a design pushes the operating point, marked here for the chapter's
case study, toward that corner while holding enough margin to survive the
scatter in the limits.}
\label{fig:vol04ch13:beta-troyon}
\end{figure}


\paragraph{Worked example: the Troyon limit and the current a beta target demands.}
Read Figure~\ref{fig:vol04ch13:beta-troyon} as a design constraint. The
case-study tokamak has minor radius $a = 2\,\si{\meter}$, field $B =
5.3\,\si{\tesla}$, and plasma current $I_p = 15\,\si{\mega\ampere}$, so the
normalised current is $I_p/(aB) = 15/[(2)(5.3)] \approx 1.42\,\si{\mega
\ampere}/\si{\meter}/\si{\tesla}$. With the Troyon constant $\beta_N = 2.8$,
the beta limit is $\beta_{\max} = 2.8\times 1.42 \approx 4.0\%$, so the
operating beta of $3.4\%$ computed above sits at about $85\%$ of the Troyon
limit, a deliberate margin rather than an accident. Ask the inverse question:
to run at a target $\beta = 5\%$ with the same field and size, the normalised
current would have to rise to $I_p/(aB) = 5/2.8 \approx 1.79$, so $I_p =
1.79(2)(5.3) \approx 19\,\si{\mega\ampere}$, a $27\%$ current increase that
drops the edge safety factor and moves the operating point toward the
current-limit boundary of the figure. The two limits pull against each other,
which is why a machine cannot chase beta by current alone and why higher field
and larger size, which raise the current a given safety factor allows, are the
levers a reactor design actually pulls.

\paragraph{Worked example: the theta-pinch and pressure balance across a column.}
The pressure-tension decomposition is easiest to read where the field is
straight, as in a theta-pinch: a cylinder of plasma with an axial field $B_z$
outside and a lower field inside, the difference carrying an azimuthal
(theta-directed) current in the plasma skin. With straight field lines the
tension term vanishes and the balance $\grad(p + B^{2}/2\mu_0) = 0$ says the
total pressure is uniform across the column. Take a plasma that must hold a
thermal pressure $p = 2\times 10^{5}\,\si{\pascal}$ with no field inside, so the
outside field alone balances it: $B_z^{2}/(2\mu_0) = p$, giving $B_z =
\sqrt{2\mu_0 p} = \sqrt{2(4\pi\times 10^{-7})(2\times 10^{5})} = \sqrt{0.503}
\approx 0.71\,\si{\tesla}$. A modest field confines a substantial pressure
because the magnetic pressure grows as $B^{2}$; doubling the field to $1.4\,\si{
\tesla}$ would hold four times the plasma pressure. The azimuthal current the
skin carries follows from Amp\`ere's law across the skin of thickness $\delta$:
$J_\theta \delta = B_z/\mu_0 = 0.71/(4\pi\times 10^{-7}) \approx 5.6\times
10^{5}\,\si{\ampere\per\meter}$, so a $1\,\si{\milli\meter}$ skin carries a
current density of $\sim 5.6\times 10^{8}\,\si{\ampere\per\meter\squared}$. The
theta-pinch was an early fusion concept precisely because the arithmetic is so
favourable on paper; its downfall was that a straight column has ends, and the
plasma streams out of them on the ion-acoustic time, which is why closed
toroidal geometry, at the cost of curvature and its drifts, replaced the open
pinch.

\subsection{From two fluids to one: the MHD reduction}

The single-fluid MHD equations quoted earlier are a reduction of a more
honest two-fluid description, one momentum equation for the electrons and
one for the ions. The reduction shows which terms were dropped and so where
MHD fails. Each species $s$ (electrons $e$, ions $i$) obeys a continuity and
a momentum equation,
\begin{align}
\pdv{n_{s}}{t} + \divg(n_{s}\vect{u}_{s}) &= 0,\\
m_{s}n_{s}\dv{\vect{u}_{s}}{t} &= -\grad p_{s}
+ q_{s}n_{s}(\vect{E} + \vect{u}_{s}\times\vect{B})
+ \vect{R}_{s},
\end{align}
where $\vect{u}_{s}$ is the species flow, $p_{s}$ its pressure, and
$\vect{R}_{s}$ the collisional friction between species, with
$\vect{R}_{e} = -\vect{R}_{i}$ by momentum conservation. The single-fluid
variables are the mass density $\rho = n_{i}m_{i} + n_{e}m_{e} \approx
n_{i}m_{i}$, the centre-of-mass velocity $\vect{v} = (n_{i}m_{i}\vect{u}_{i}
+ n_{e}m_{e}\vect{u}_{e})/\rho \approx \vect{u}_{i}$, the current density
$\vect{J} = e(n_{i}\vect{u}_{i} - n_{e}\vect{u}_{e})$, and the total pressure
$p = p_{i} + p_{e}$.

Add the two momentum equations. The electric force cancels under
quasineutrality $n_{i} \approx n_{e}$, the friction terms cancel because
$\vect{R}_{e} = -\vect{R}_{i}$, and the magnetic terms combine into
$\vect{J}\times\vect{B}$, leaving the single-fluid momentum equation
\[
\rho\dv{\vect{v}}{t} = -\grad p + \vect{J}\times\vect{B},
\]
the second ideal-MHD equation. The Ohm's law comes from the electron
momentum equation alone. Divide it by $-e n_{e}$ and keep the leading terms;
the electron inertia $m_{e}\,\dd\vect{u}_{e}/\dd t$ is smaller than the
magnetic term by the ratio of the dynamical frequency to the electron
gyrofrequency and is dropped, giving the \engterm{generalised Ohm's law}
\[
\vect{E} + \vect{v}\times\vect{B}
= \eta\vect{J}
+ \frac{\vect{J}\times\vect{B}}{en_{e}}
- \frac{\grad p_{e}}{en_{e}}
+ \frac{m_{e}}{e^{2}n_{e}}\pdv{\vect{J}}{t}.
\]
The four terms on the right are the resistive drag, the \engterm{Hall term},
the electron-pressure term, and the electron-inertia term. Ideal MHD keeps
only the left side: it drops resistivity (infinite conductivity), the Hall
term (valid when the length scale is large compared to the ion skin depth
$c/\omega_{pi}$), the pressure term (same condition), and the inertia term
(slow compared to the electron plasma frequency). The reduction is therefore
the long-wavelength, low-frequency, high-conductivity limit, the slow window
of Figure~\ref{fig:vol04ch13:timescale-ordering}.

% Figure: the ordering of plasma timescales that justifies the MHD
% reduction. From the fast electron plasma period and electron gyro-period
% up through the ion equivalents, the MHD/Alfven time, and the slow
% resistive diffusion and confinement times. MHD lives in the slow window
% where the fast scales have been averaged out.
\begin{figure}[htbp]
\centering
\begin{tikzpicture}[>=Stealth]
  % horizontal log axis of time
  \draw[thick,->] (0,0) -- (12,0);
  \node[anchor=west] at (12.1,0) {$t$ (s), log scale};
  % tick marks with labels (representative tokamak numbers)
  \foreach \x/\lab in {0.6/$10^{-11}$, 2.4/$10^{-9}$, 4.2/$10^{-7}$,
                       6.0/$10^{-5}$, 7.8/$10^{-3}$, 9.6/$10^{-1}$,
                       11.4/$10^{1}$} {
    \draw[thick] (\x,0.1) -- (\x,-0.1);
    \node[anchor=north, font=\scriptsize] at (\x,-0.15) {\lab};
  }
  % event markers above the axis
  \draw[engblue, thick] (0.6,0) -- (0.6,0.7);
  \node[engblue, anchor=south, font=\scriptsize, rotate=40, anchor=west]
    at (0.6,0.7) {$\omega_{pe}^{-1}$};
  \draw[engblue, thick] (1.2,0) -- (1.2,1.4);
  \node[engblue, anchor=south, font=\scriptsize, rotate=40, anchor=west]
    at (1.2,1.4) {$\omega_{ce}^{-1}$};
  \draw[engorange, thick] (3.6,0) -- (3.6,0.7);
  \node[engorange, anchor=south, font=\scriptsize, rotate=40, anchor=west]
    at (3.6,0.7) {$\omega_{ci}^{-1}$};
  \draw[englime, line width=1.2pt] (6.0,0) -- (6.0,1.4);
  \node[englime, anchor=south, font=\scriptsize, rotate=40, anchor=west]
    at (6.0,1.4) {$\tau_{A}$ (Alfv\'en)};
  \draw[engred, thick] (8.4,0) -- (8.4,0.7);
  \node[engred, anchor=south, font=\scriptsize, rotate=40, anchor=west]
    at (8.4,0.7) {$\tau_{E}$ (confinement)};
  \draw[engred, thick] (10.8,0) -- (10.8,1.4);
  \node[engred, anchor=south, font=\scriptsize, rotate=40, anchor=west]
    at (10.8,1.4) {$\tau_{R}$ (resistive)};
  % MHD validity band
  \draw[enggray, dashed] (5.4,-0.9) rectangle (11.4,-0.45);
  \node[enggray, font=\scriptsize] at (8.4,-0.675)
    {ideal-MHD window: faster scales averaged out};
\end{tikzpicture}
\caption[Plasma timescale ordering]{The hierarchy of plasma timescales for
representative tokamak parameters, on a logarithmic time axis. The electron
plasma period $\omega_{pe}^{-1}$ and the electron gyro-period
$\omega_{ce}^{-1}$ are the fastest; the ion gyro-period $\omega_{ci}^{-1}$
is slower by the mass ratio. The single-fluid magnetohydrodynamic
description is valid in the slow window to the right of the ion scales,
where the fast electron and gyration motions have been averaged into
fluid moments. The Alfv\'en time $\tau_{A}$ sets the speed of ideal-MHD
instabilities, the confinement time $\tau_{E}$ sets the energy loss, and
the resistive time $\tau_{R}$ governs the slow rearrangement of field
topology. Ideal MHD assumes $\tau_{R}\to\infty$; finite resistivity
re-enters only on the longest scale.}
\label{fig:vol04ch13:timescale-ordering}
\end{figure}


The dropped terms are the ones an engineer reaches for when MHD fails. The
resistive term restores reconnection and the resistive instabilities. The
Hall term matters in reconnection layers and in the Hall thruster, where the
ion skin depth is comparable to the channel, and it is the reason that
device carries the name it does. The electron-pressure term drives the
diamagnetic drift and the gradients that feed drift-wave turbulence, the
turbulence that sets the anomalous transport of the earlier mastery box. The
inertia term carries the electron plasma oscillation and the high-frequency
waves. Reading which term has been dropped tells the engineer which physics
the MHD model cannot see.

\subsection{The frozen-in flux theorem, in full}

The most consequential result of ideal MHD is that magnetic field lines move
with the fluid. The proof combines the ideal Ohm's law with Faraday's law
and the kinematics of a moving surface. Start from the ideal Ohm's law
$\vect{E} = -\vect{v}\times\vect{B}$ and substitute into Faraday's law
$\partial\vect{B}/\partial t = -\curl\vect{E}$ to get the induction equation
\[
\pdv{\vect{B}}{t} = \curl(\vect{v}\times\vect{B}).
\]
Consider the magnetic flux $\Phi = \int_{S(t)}\vect{B}\cdot\dd\vect{A}$
through a surface $S(t)$ whose boundary is a loop of fluid elements, so the
surface is carried along by the flow. The total rate of change of the flux
through a comoving surface has two contributions, the explicit time change
of $\vect{B}$ and the sweeping of the boundary through the field,
\[
\dv{\Phi}{t}
= \int_{S}\pdv{\vect{B}}{t}\cdot\dd\vect{A}
+ \oint_{\partial S}\vect{B}\cdot(\vect{v}\times\dd\vect{\ell}).
\]
The boundary term rewrites with the scalar triple product as $-\oint(\vect{v}
\times\vect{B})\cdot\dd\vect{\ell}$, which by Stokes' theorem is
$-\int_{S}\curl(\vect{v}\times\vect{B})\cdot\dd\vect{A}$. Adding it to the
first term and using the induction equation,
\[
\dv{\Phi}{t}
= \int_{S}\left[\pdv{\vect{B}}{t}
- \curl(\vect{v}\times\vect{B})\right]\cdot\dd\vect{A} = 0.
\]
The flux through every comoving loop is constant in time. Two fluid elements
that share a field line at one instant share it for all time, so the field
lines are advected by the flow as if painted on the fluid. This is the
\engterm{frozen-in} condition, and it is the reason a magnetically confined
plasma cannot simply slide off its field: to move across the field the
plasma would have to break the flux conservation, which ideal MHD forbids.

The theorem fails exactly where the dropped resistive term re-enters. With
finite resistivity the induction equation gains a diffusion term,
\[
\pdv{\vect{B}}{t} = \curl(\vect{v}\times\vect{B})
+ \frac{\eta}{\mu_{0}}\nabla^{2}\vect{B},
\]
and the field now diffuses through the fluid on the resistive time
$\tau_{R} = \mu_{0}L^{2}/\eta$. The ratio of the advection term to the
diffusion term is the \engterm{magnetic Reynolds number} $R_{m} = \mu_{0}Lv/
\eta$, the same group as the Lundquist number with the flow speed in place of
the Alfv\'en speed. Where $R_{m}\gg 1$ the field is frozen and topology is
conserved; where $R_{m}\sim 1$, in thin current sheets and at small scales,
the field slips and reconnects. The engineering reading is that flux is
conserved on the large scale that sets the confinement and violated on the
small scale that sets the disruptions, the two regimes the chapter's failure
section returns to.

\paragraph{Worked example: magnetic Reynolds number of a tokamak and a lamp.}
For the case-study tokamak at $T_{e} = 12\,\si{\kilo\electronvolt}$, the
Spitzer resistivity is $\eta \approx 5.2\times 10^{-5}(17)/(1.2\times
10^{4})^{3/2} \approx 6.8\times 10^{-10}\,\si{\ohm\meter}$. Taking a flow
speed of order the Alfv\'en speed $v_{A}\approx 7\times 10^{6}\,\si{\meter
\per\second}$ and a length $L = 2\,\si{\meter}$, the magnetic Reynolds
number is $R_{m} = \mu_{0}Lv/\eta = (4\pi\times 10^{-7})(2)(7\times 10^{6})/
(6.8\times 10^{-10}) \approx 2.6\times 10^{10}$, so the field is frozen to
better than ten significant figures on the device scale and ideal MHD is an
excellent description of the bulk. Contrast a fluorescent lamp at $T_{e} =
1\,\si{\electronvolt}$, where the resistivity is $\eta \approx 5.2\times
10^{-5}(10)/(1)^{3/2} \approx 5\times 10^{-4}\,\si{\ohm\meter}$, a flow speed
of $\sim 10^{3}\,\si{\meter\per\second}$, and a length $L = 0.01\,\si{\meter}$:
$R_{m} = (4\pi\times 10^{-7})(0.01)(10^{3})/(5\times 10^{-4}) \approx
2.5\times 10^{-5}$, far below one. The lamp plasma does not freeze flux at
all; its field diffuses through the gas almost instantly, which is why
low-temperature discharges are described by collisional transport and
chemistry rather than by MHD. The same plasma physics gives opposite
descriptions at the two ends of the temperature range because the single
group $R_{m}$ moves by fifteen orders of magnitude between them.

\begin{mastery}
The reduction above produced \engterm{ideal} MHD, but a magnetised plasma
supports a richer set of transport processes that survive even when the
collisional friction is small, collected under the name
\engterm{neoclassical transport}. In a torus a fraction of the particles are
trapped between the magnetic mirrors formed by the $1/R$ variation of the
field strength along a field line, and these trapped particles execute
banana-shaped orbits whose width is much larger than the gyroradius. Their
collisional scattering across the field steps by the banana width rather than
the gyroradius, so the diffusion coefficient is enhanced over the naive
classical estimate by a factor of order $(R/a)^{3/2}q^{2}$, the
neoclassical factor, which can be one to two orders of magnitude. The
trapped-particle population also drives the \engterm{bootstrap current}, a
self-generated toroidal current proportional to the pressure gradient that a
tokamak can use to supplement or replace the transformer-driven current,
relaxing the pulsed-operation limit that the transformer imposes. A reactor
design that aims at steady state leans on a large bootstrap fraction, so the
neoclassical theory that the ideal reduction left out is, for the
steady-state engineer, the part that matters most. Even neoclassical
transport, though, is dwarfed in practice by the turbulent anomalous
transport of the earlier mastery box, so the working confinement time
remains an empirical input rather than a first-principles output.
\end{mastery}

\begin{mastery}
The static force balance $\grad p = \vect{J}\times\vect{B}$ becomes, in the
axisymmetric geometry of a tokamak, a single scalar partial differential
equation for the poloidal flux function $\psi$, the \engterm{Grad-Shafranov
equation}. Writing the poloidal field as $\vect{B}_{\mathrm{pol}} = \grad\psi
\times\grad\phi$ with $\phi$ the toroidal angle, and the toroidal field through a
current function $F(\psi) = R B_\phi$, the balance reduces to
\[
R\pdv{}{R}\!\left(\frac{1}{R}\pdv{\psi}{R}\right)
+ \frac{\partial^{2}\psi}{\partial Z^{2}}
= -\mu_0 R^{2}\dv{p}{\psi} - F\dv{F}{\psi},
\]
where $R$ and $Z$ are the major-radius and vertical coordinates. The two free
functions $p(\psi)$ and $F(\psi)$, the pressure and the current profiles, are
inputs a designer chooses; the equation then returns the shape of the nested
flux surfaces the plasma settles into. The whole equilibrium engineering of a
tokamak, the elongation and triangulation of the cross-section that raise the
achievable beta and current, the position of the magnetic axis, the shift of the
flux surfaces outward under pressure known as the Shafranov shift, is the
solution of this one equation for chosen profiles and coil currents. A control
system reconstructs $\psi$ in real time from magnetic measurements outside the
plasma by solving the Grad-Shafranov equation as an inverse problem, and it
adjusts the shaping-coil currents to hold the plasma on its intended shape. The
equation is the bridge from the abstract force balance of the foundations
section to the concrete engineering of the coil set that surrounds every
tokamak.
\end{mastery}

\begin{mastery}
The frozen-in theorem conserves the flux through every comoving loop, and a
weighted integral of it survives even when the theorem itself fails a little,
the \engterm{magnetic helicity} $K = \int\vect{A}\cdot\vect{B}\,\dd V$, with
$\vect{A}$ the vector potential. Helicity measures the linkage and twist of the
field lines, and it is conserved far better than the magnetic energy when a
plasma reconnects, because reconnection rearranges topology locally while the
total linkage is nearly untouched. Woltjer and Taylor drew the consequence: a
turbulent, reconnecting plasma relaxes to the state of lowest magnetic energy
at fixed total helicity, and that state solves $\curl\vect{B} = \lambda\vect{B}$
with a single constant $\lambda$, a force-free field in which the current runs
everywhere parallel to the field so the $\vect{J}\times\vect{B}$ force vanishes.
This \engterm{Taylor relaxation} is why a reversed-field pinch, left to itself,
settles into a reproducible field profile in which the toroidal field actually
reverses at the edge, a configuration no external coil imposes and that the
plasma finds on its own by relaxing at constant helicity. The engineering
payoff is twofold. A device can be brought to a wanted magnetic state by
injecting helicity rather than by driving current directly, the basis of
helicity-injection current drive in spheromaks and spherical tokamaks; and the
same conservation law bounds how much magnetic energy a solar flare or a
disruption can release, since the field cannot relax below the lowest-energy
state its frozen-in helicity allows. Reading which quantity a plasma conserves
under reconnection, energy poorly and helicity well, is how an engineer
predicts the state a turbulent plasma will choose when it is left to find its
own equilibrium.
\end{mastery}

\section{Plasma diagnostics}\pathtag{standard}

Measuring the parameters of a plasma without disturbing it is the
discipline of plasma diagnostics. Three working techniques recur:

A \engterm{Langmuir probe} is a small conducting electrode inserted
into the plasma; its $I$-$V$ characteristic reveals the electron
density, electron temperature, and plasma potential through three
identifiable regions: ion saturation (probe biased strongly negative,
collecting only ions), electron-temperature region (probe slightly
negative, exponential I-V), and electron saturation (probe biased
positive, collecting electrons). The slope of the exponential region
gives $T_{e}$; the saturation values give $n_{e}$. Langmuir probes
are the workhorse of low-temperature plasma diagnostics (industrial
processing plasmas, edge plasmas in tokamaks).

% Figure: idealised Langmuir-probe I-V characteristic, showing the three
% regions: ion saturation, the exponential electron-retardation region
% whose slope gives T_e, and electron saturation. Data computed inline in
% code/langmuir_probe.py and saved to data/langmuir_iv.csv.
\begin{figure}[htbp]
\centering
\begin{tikzpicture}
\begin{axis}[
  width=0.78\textwidth, height=0.5\textwidth,
  xlabel={probe bias $V$ (V)},
  ylabel={probe current $I$ (mA)},
  xmin=-40, xmax=25, ymin=-3, ymax=12,
  grid=both, grid style={enggray!20},
  every axis plot/.append style={thick},
]
  \addplot[engblue, mark=*, mark size=0.7pt] table[
    col sep=comma, x=bias_V, y=current_mA]
    {volumes/04-energy/ch13-plasma-physics/data/langmuir_iv.csv};
  \draw[dashed, enggray] (axis cs:-40,0) -- (axis cs:25,0);
  \node[engred, font=\scriptsize, anchor=west] at (axis cs:-38,-1.6)
    {ion saturation};
  \node[font=\scriptsize, anchor=south] at (axis cs:-8,2.5)
    {exponential: slope $\to T_e$};
  \node[engred, font=\scriptsize, anchor=east] at (axis cs:24,10.4)
    {electron saturation};
  \draw[dashed, enggray] (axis cs:5,-3) -- (axis cs:5,12);
  \node[enggray, font=\scriptsize, anchor=south west] at (axis cs:5.3,8.5)
    {$V_p$};
\end{axis}
\end{tikzpicture}
\caption[Langmuir-probe characteristic]{Idealised current-voltage
characteristic of a Langmuir probe. Strongly negative bias repels all
electrons and the probe draws only the small ion-saturation current.
Raising the bias lets a growing fraction of the electron distribution
reach the probe; in this electron-retardation region the current rises
exponentially, and the slope of $\ln(I-I_{\mathrm{ion}})$ against $V$
gives the electron temperature $T_e$. Above the plasma potential $V_p$
the probe collects the full electron flux and the current saturates;
the saturation values fix the density $n_e$. The data are generated by
\texttt{code/langmuir\_probe.py}.}
\label{fig:vol04ch13:langmuir-iv}
\end{figure}


Figure~\ref{fig:vol04ch13:langmuir-iv} shows the three regions of the
characteristic. The standard analysis fits a straight line to
$\ln(I - I_{\mathrm{ion}})$ against bias in the retardation region; the
reciprocal slope is the electron temperature in volts.
The companion \texttt{code/langmuir\_probe.py} generates a synthetic
characteristic for a $3\,\si{\electronvolt}$ plasma and recovers the input
temperature from the slope, which is the calculation a probe-diagnostic
program runs on every sweep.

The numbers a probe returns come out of the same Bohm and saturation
physics the sheath section developed. The electron temperature is read from
the retardation slope: in that region the collected electron current is
$I_e = I_{e0}\exp(e(V - V_p)/k_B T_e)$, so a plot of $\ln I_e$ against bias
$V$ is a straight line of slope $e/k_B T_e$, and the reciprocal slope read
in volts is $T_e$ in electronvolts directly. A decade of current change over
$6.9\,\si{\volt}$ of bias, for instance, gives $k_B T_e/e = 6.9/\ln 10
\approx 3\,\si{\volt}$, so $T_e = 3\,\si{\electronvolt}$. The density comes
from the ion saturation current, which is the Bohm flux times the probe
area, $I_{\mathrm{sat}} = 0.6\,n_e e u_B A_{\mathrm{probe}}$ with $u_B =
\sqrt{k_B T_e/m_i}$. Inverting for $n_e$ requires the same effective ion
mass the sheath analysis used, which is why an honest probe report states
the assumed ion species: a factor-of-two error in $m_i$ moves the inferred
density by $\sqrt{2}$. The plasma potential $V_p$ is read from the knee
where the characteristic bends from the exponential region into electron
saturation, and the floating potential $V_f$, where the net current is
zero, sits below $V_p$ by the same sheath drop a floating wall carries. A
probe thus returns $T_e$, $n_e$, $V_p$, and $V_f$ from a single voltage
sweep, which is why it remains the workhorse of low-temperature plasma
diagnostics despite the perturbation it introduces by drawing current.

\engterm{Thomson scattering} uses a laser pulse fired into the plasma;
the scattered light is Doppler-broadened by the electron-velocity
distribution and Doppler-shifted by the bulk plasma motion. The
width gives $T_{e}$ directly through Doppler broadening; the
amplitude gives $n_{e}$ through the scattering cross-section
($\sigma_{T} = 6.65\times 10^{-29}\,\si{\meter\squared}$, the
classical Thomson cross-section). Thomson scattering is non-
intrusive (the laser pulse perturbs the plasma negligibly) and gives
local measurements; the price is that the scattered signal is weak
and the technique requires either a high-energy pulsed laser or a
high-power CW laser with a sensitive detector.

\engterm{Spectroscopy} measures the line and continuum radiation
emitted by the plasma. Line ratios (intensities of specific atomic
or ionic transitions) give electron temperature; line broadening
gives ion temperature (through Doppler) and density (through Stark
effect in dense plasmas). Spectroscopy is integrated along the line
of sight and requires Abel-inversion or tomographic reconstruction
for spatially resolved measurements. The advantage is that
spectroscopy is purely passive; the plasma is not perturbed at all.

The three techniques cover the density-temperature map between them. A
Langmuir probe survives only where the plasma is cool enough not to melt it,
so it serves low-temperature processing plasmas and the cool edge of a fusion
device but cannot enter a fusion core. Thomson scattering reads the core
directly because a laser can be fired into any plasma, but its signal is faint:
the scattered power is the incident power times the electron density times the
Thomson cross-section times the scattering length, and for a fusion plasma that
product is of order $10^{-13}$ of the laser pulse, so the technique demands a
joule-class laser, a large collection solid angle, and photon counting.
Spectroscopy is the cheapest and least intrusive, needing only a light path and
a spectrometer, but it returns a line-of-sight integral that must be inverted,
and its temperature reading rests on an atomic model of the emitting species
whose accuracy sets the accuracy of the result. An honest diagnostic programme
runs at least two of the three so that each checks the others: a probe and a
spectrometer at the edge, Thomson scattering and spectroscopy in the core, with
the density from one method and the temperature from another cross-validated
before either is trusted in the power balance.

A fourth technique reads the line-integrated density through the plasma's
refractive index rather than through any emitted or scattered light.
\engterm{Interferometry} splits a laser or microwave beam, passes one arm
through the plasma and holds the other as a reference, and recombines them; the
plasma arm accumulates a phase shift because the plasma refractive index $n =
\sqrt{1-\omega_p^{2}/\omega^{2}}$ is below one, so the wave travels faster than
in vacuum and arrives advanced. For a probing frequency well above the plasma
frequency the phase shift is $\Delta\phi = (\omega/c)\int(1-n)\,\dd\ell
\approx (e^{2}/2m_e\varepsilon_0 c\omega)\int n_e\,\dd\ell$, proportional to the
line-integrated electron density along the beam. The measurement is a fringe
count, so it is absolute and needs no calibration against a source strength,
which is why an interferometer is the primary density diagnostic on most
tokamaks; its limits are that it returns a line integral that must be inverted
for a profile and that a fast density change can slip a fringe unnoticed if the
detector is too slow.

\paragraph{Worked example: the interferometer fringe shift across a fusion plasma.}
Probe the case-study plasma with a far-infrared laser at $\lambda = 119\,\si{
\micro\meter}$ ($\omega = 2\pi c/\lambda \approx 1.58\times 10^{13}\,\si{
\radian\per\second}$) along a chord of length $\ell = 2\,\si{\meter}$ through
$n_e = 10^{20}\,\si{\per\meter\cubed}$. The phase shift is $\Delta\phi =
(e^{2}/2m_e\varepsilon_0 c\omega)\,n_e\ell$. The prefactor is $(1.6\times
10^{-19})^{2}/[2(9.11\times 10^{-31})(8.85\times 10^{-12})(3\times 10^{8})
(1.58\times 10^{13})] \approx 3.35\times 10^{-19}\,\si{\meter\squared}$, so
$\Delta\phi = (3.35\times 10^{-16})(10^{20})(2) \approx 67\,\si{\radian}$, about
$67/2\pi \approx 11$ fringes. A shift of eleven whole fringes is easy to count,
which is why $119\,\si{\micro\meter}$ is a standard tokamak interferometer line:
it gives several fringes of signal, enough to resolve the density well, while
staying far enough above the plasma frequency that refraction does not bend the
beam off the detector. Choosing the wavelength is the design trade the phase
formula makes plain, since $\Delta\phi\propto\lambda$ favours a long wavelength
for signal while refraction, which scales as $\lambda^{2}$, favours a short one,
so the working lines sit in the far infrared where the two demands balance.

\paragraph{Worked example: the Thomson scattering signal and why it is hard.}
Fire a $J_L = 1\,\si{\joule}$ laser pulse at $\lambda = 1064\,\si{\nano\meter}$
through a fusion plasma of $n_e = 10^{20}\,\si{\per\meter\cubed}$ and collect
the light scattered from a $L = 1\,\si{\centi\meter}$ length of the beam into a
detector subtending a solid angle $\Omega = 10^{-2}\,\si{\steradian}$. The
fraction of the beam scattered per unit length is $n_e\sigma_T = (10^{20})
(6.65\times 10^{-29}) \approx 6.65\times 10^{-9}\,\si{\per\meter}$, so over the
$1\,\si{\centi\meter}$ scattering length the scattered fraction is $6.65\times
10^{-11}$. Of that, the fraction reaching the detector is $\Omega/4\pi \approx
10^{-2}/12.6 \approx 8\times 10^{-4}$, so the collected energy is $J_L\times
6.65\times 10^{-11}\times 8\times 10^{-4} \approx 4\times 10^{-14}\,\si{\joule}$.
A $1064\,\si{\nano\meter}$ photon carries $hc/\lambda \approx 1.9\times
10^{-19}\,\si{\joule}$, so the detector sees $\approx 4\times 10^{-14}/1.9
\times 10^{-19} \approx 2\times 10^{5}$ photons per pulse, a countable but small
number that must be spread across a spectrometer to resolve the Doppler width.
The width itself gives the temperature: a $1\,\si{\kilo\electronvolt}$ electron
population has a thermal speed $v_{te} \approx 1.3\times 10^{7}\,\si{\meter\per
\second}$, so the scattered light is Doppler-broadened by $\Delta\lambda/\lambda
\approx 2v_{te}/c \approx 0.09$, a broad $\sim 100\,\si{\nano\meter}$ smear that
a modest spectrometer resolves. The calculation shows the trade the diagnostic
lives under: the density sets the amplitude and the temperature sets the width,
so a single Thomson measurement returns both, but the faint signal forces a
powerful laser and a sensitive, well-shielded detector, which is why the
technique is reserved for the measurements a probe cannot make.

% Figure: oblique reflection of a radio wave from the ionosphere and the
% secant law. A ray launched vertically reflects only if its frequency is below
% the critical frequency; a ray launched obliquely at zenith angle phi reflects
% at a higher frequency, the maximum usable frequency, larger than the critical
% frequency by sec(phi). The geometry is why long HF links exploit low takeoff
% angles.
\begin{figure}[htbp]
\centering
\begin{tikzpicture}[>=Latex, scale=1.0]
  % ground
  \draw[engblack, thick] (-4.5,0) -- (4.5,0);
  \node[font=\scriptsize, anchor=north west] at (-4.5,0) {ground};
  % ionosphere layer (band)
  \fill[engblue!12] (-4.5,2.6) rectangle (4.5,3.2);
  \draw[engblue!50] (-4.5,2.6) -- (4.5,2.6);
  \draw[engblue!50] (-4.5,3.2) -- (4.5,3.2);
  \node[engblue, font=\scriptsize, anchor=south west] at (-4.5,3.2)
    {ionospheric layer, peak $n_e$};
  % transmitter and receiver
  \fill[engblack] (-3.5,0) circle (1.5pt);
  \node[font=\scriptsize, anchor=north] at (-3.5,0) {TX};
  \fill[engblack] (3.5,0) circle (1.5pt);
  \node[font=\scriptsize, anchor=north] at (3.5,0) {RX};
  % vertical ray up and back (critical frequency, reflects at f_c)
  \draw[engorange, thick, ->] (0,0) -- (0,2.6);
  \draw[engorange, thick, ->] (0,2.6) -- (0,0.1);
  \node[engorange, font=\scriptsize, anchor=west] at (0.1,1.3)
    {vertical: reflects at $f_c$};
  % oblique ray path TX -> reflection point -> RX
  \draw[engblue, thick, ->] (-3.5,0) -- (0,2.6);
  \draw[engblue, thick, ->] (0,2.6) -- (3.5,0.05);
  % zenith angle marker at reflection point
  \draw[engblue, dashed] (0,2.6) -- (0,3.4);
  \draw[engblue] (0,3.2) arc (90:145:0.6);
  \node[engblue, font=\scriptsize] at (-0.55,3.05) {$\phi$};
  \node[engblue, font=\scriptsize, anchor=south] at (-1.9,1.25)
    {oblique: $f_{\mathrm{MUF}}=f_c\sec\phi$};
\end{tikzpicture}
\caption[Ionospheric reflection and the secant law]{Oblique reflection of a
radio wave from the ionosphere. A ray launched straight up (orange) reflects
only if its frequency lies below the layer's critical frequency $f_c$, the
plasma frequency at the density peak. A ray launched obliquely (blue) at zenith
angle $\phi$ meets the layer at a shallow angle and reflects even when its
frequency exceeds $f_c$, because the wave only has to be cut off by the reduced
component of density along its path. The highest frequency a given path can
return is the maximum usable frequency, $f_{\mathrm{MUF}}=f_c\sec\phi$, larger
than the vertical critical frequency by the secant of the zenith angle. A long
link uses a low takeoff angle to make $\phi$ large and push the usable frequency
well above $f_c$, which is why the working frequency of an HF circuit rises with
its length and falls at night as the ionosphere thins.}
\label{fig:vol04ch13:ionospheric-refraction}
\end{figure}


\paragraph{Worked example: the maximum usable frequency and the secant law.}
Radio propagation is a diagnostic of the ionosphere read backward, and the
secant law that governs it, Figure~\ref{fig:vol04ch13:ionospheric-refraction},
follows from the same cutoff the plasma frequency sets. A wave reflects from a
layer when the local plasma frequency matches the wave frequency resolved along
the density gradient. At vertical incidence the critical frequency is the peak
plasma frequency, $f_c \approx 9\sqrt{n_e/10^{12}}\,\si{\mega\hertz}$, so a
daytime F2 peak of $n_e = 1.5\times 10^{12}\,\si{\per\meter\cubed}$ gives $f_c
\approx 9\sqrt{1.5} \approx 11\,\si{\mega\hertz}$. A wave launched obliquely at
zenith angle $\phi$ reflects at a higher frequency, the maximum usable
frequency $f_{\mathrm{MUF}} = f_c\sec\phi$, because only the component of the
wave normal to the layer has to be cut off. For a $2000\,\si{\kilo\meter}$ hop
reflecting near $300\,\si{\kilo\meter}$ altitude, the takeoff geometry gives a
zenith angle of about $\phi = 74^{\circ}$ at the reflection point, so $\sec\phi
\approx 3.6$ and $f_{\mathrm{MUF}} \approx 11\times 3.6 \approx 40\,\si{\mega
\hertz}$, four times the vertical critical frequency. A long-range HF circuit
therefore works far above the vertical cutoff by using a low takeoff angle, and
its usable frequency drops at night as the F2 density falls and $f_c$ with it.
The same physics that makes a fusion plasma opaque to microwaves below $\omega_p$
lets a shortwave broadcaster reach a continent by choosing a frequency just
below the path's MUF.

\section{Industrial plasmas: etching, deposition, lighting, propulsion}\pathtag{standard}

Industrial plasmas are the bridge from plasma physics to
manufacturing.

Every industrial discharge begins by breaking down a neutral gas, and the
voltage it takes to strike the plasma is not a material constant but a
function of the pressure and the gap together. A free electron in the gap
gains energy from the field and ionises a neutral if it collects enough
between collisions; each ionisation frees another electron, and when the
avalanche multiplication just replaces the electrons lost to the walls the gas
breaks down. Paschen's law captures the balance: the breakdown voltage depends
on the gas only through the product $pd$ of pressure and gap, and it has a
minimum. At small $pd$ an electron crosses the gap with too few collisions to
multiply, so the voltage must climb; at large $pd$ the electron loses its
energy to frequent collisions before it reaches the ionisation threshold, so
again the voltage must climb. Between the two extremes a discharge strikes at
a few hundred volts. Figure~\ref{fig:vol04ch13:paschen} draws the curve for
two gases, and the minimum is where a low-pressure processing tool is run,
because it is the pressure at which the least applied voltage sustains the
plasma.

% Figure: Paschen curves, the breakdown voltage of a gas gap against the
% product of pressure and electrode spacing, for two gases. The curve has a
% minimum: too small a pd gives too few ionising collisions, too large a pd
% loses electron energy to collisions before ionisation, so a discharge
% strikes most easily at an intermediate pd. The minimum sets the operating
% pressure of every low-pressure discharge tool.
\begin{figure}[htbp]
\centering
\begin{tikzpicture}
\begin{semilogxaxis}[
  width=0.82\textwidth, height=0.55\textwidth,
  xlabel={$pd$ (\si{\pascal}\,\si{\meter})},
  ylabel={breakdown voltage $V_b$ (\si{\volt})},
  xmin=0.1, xmax=1000, ymin=100, ymax=3000,
  grid=both, grid style={enggray!20},
  legend style={font=\scriptsize, at={(0.98,0.98)}, anchor=north east},
  every axis plot/.append style={thick},
]
  % Paschen law: Vb = B pd / ln( A pd / ln(1+1/gamma) )
  % argon-like constants: A=12, B=180 (per Pa m, scaled), gamma-term ~ 2.5
  \addplot[engblue, domain=0.3:1000, samples=120]
    {180*x/ln(max(1.01, 12*x/2.5))};
  \addlegendentry{argon-like}
  % air-like: A=11, B=340
  \addplot[engred, domain=0.6:1000, samples=120]
    {340*x/ln(max(1.01, 11*x/2.5))};
  \addlegendentry{air-like}
  \addplot[only marks, mark=*, engblue, mark size=1.8pt] coordinates {(1.0,270)};
  \node[engblue, font=\scriptsize, anchor=north] at (axis cs:1.0,255) {minimum};
\end{semilogxaxis}
\end{tikzpicture}
\caption[Paschen breakdown curves]{Breakdown voltage of a gas gap against the
product of pressure and electrode spacing, on a semi-log scale, for an
argon-like and an air-like gas. Each curve has a minimum because breakdown is
a competition: at small $pd$ an electron crosses the gap with too few
ionising collisions to build an avalanche, so the voltage must rise to
compensate, while at large $pd$ the electron loses its energy in frequent
collisions before it reaches the ionisation threshold, so again the voltage
must rise. The discharge strikes most easily at the intermediate $pd$ of the
minimum, a few hundred volts near $pd\approx 1\,\si{\pascal}\,\si{\meter}$.
The curve is the reason a low-pressure processing tool is run near the
minimum, where the least voltage sustains the plasma, and the reason a
high-voltage vacuum insulator fails not in hard vacuum but at the modest
pressure where its gap sits at the Paschen minimum.}
\label{fig:vol04ch13:paschen}
\end{figure}


\paragraph{Worked example: the Paschen minimum and the strike voltage of an etch tool.}
Paschen's law writes the breakdown voltage as $V_b = B\,pd/\ln[A\,pd/\ln(1+
1/\gamma)]$, with $A$ and $B$ gas constants and $\gamma$ the secondary-electron
yield at the cathode. The minimum is found by differentiation and lands at
$(pd)_{\min} = (e/A)\ln(1+1/\gamma)$ with the Euler number $e = 2.718$, giving
a minimum voltage $V_{b,\min} = (B/A)\,e\ln(1+1/\gamma)$. For argon, take $A
\approx 12\,\si{\per\pascal\per\meter}$, $B \approx 180\,\si{\volt\per\pascal
\per\meter}$, and a secondary yield $\gamma \approx 0.1$, so $\ln(1+1/\gamma)
= \ln 11 \approx 2.4$. The minimum $pd$ is $(pd)_{\min} = (2.718/12)(2.4)
\approx 0.54\,\si{\pascal}\,\si{\meter}$ and the minimum voltage is
$V_{b,\min} = (180/12)(2.718)(2.4) \approx 98\,\si{\volt}$, a hundred volts at
a $pd$ near half a pascal-metre. The etch reactor of the second case study
runs at $5\,\si{\pascal}$ across a $30\,\si{\milli\meter}$ gap, so its $pd =
(5)(0.03) = 0.15\,\si{\pascal}\,\si{\meter}$, a little below the minimum on
the steep left branch, which is why its strike voltage sits well above the
$98\,\si{\volt}$ floor and why lowering the pressure further would make the
discharge harder, not easier, to ignite. The same curve read at high $pd$
explains a failure of high-voltage vacuum hardware: an insulator that holds
off a megavolt in hard vacuum can flash over at a modest leaked pressure if
its gap slides onto the Paschen minimum, so a vacuum designer keeps every gap
either far below or far above the dangerous $pd$.

\engterm{Plasma etching} uses a low-pressure RF discharge to produce
reactive species (radicals, ions) that etch a substrate selectively.
A typical reactive-ion etching (RIE) chamber operates at $1$ to
$100\,\si{\pascal}$ (around $10^{-2}$ to $1\,\text{torr}$), $13.56\,\si{
\mega\hertz}$ RF excitation, $100$ to $1000\,\si{\watt}$ power, and
process gases like CF$_{4}$ (for silicon etching), Cl$_{2}$ (for
metal etching), or O$_{2}$ (for photoresist removal). The directional
ion bombardment from the plasma sheath produces anisotropic etching
(vertical walls of trenches and lines) at sub-micron dimensions. The
process is the working tool of every modern semiconductor fab; its
control is a multi-billion-dollar engineering discipline.

% Figure: plasma sheath. Potential drops from the quasineutral bulk to a
% negative wall over a few Debye lengths; the schematic shows the bulk,
% the sheath edge, and the ion-acceleration region that produces the
% directional bombardment used in reactive-ion etching.
\begin{figure}[htbp]
\centering
\begin{tikzpicture}
\begin{axis}[
  width=0.78\textwidth, height=0.46\textwidth,
  xlabel={distance from wall (Debye lengths $x/\lambda_D$)},
  ylabel={potential $\phi$ (V)},
  domain=0:8, samples=120,
  xmin=0, xmax=8, ymin=-22, ymax=3,
  grid=both, grid style={enggray!20},
  every axis plot/.append style={thick},
]
  % sheath: potential rises from wall (-20 V) to ~0 in the bulk
  \addplot[engblue] {-20*exp(-x)};
  \draw[dashed, enggray] (axis cs:0,0) -- (axis cs:8,0);
  \node[engblue, font=\scriptsize, anchor=west] at (axis cs:3.2,-1.6)
    {quasineutral bulk $\phi\approx 0$};
  \draw[engred, ->, thick] (axis cs:1.6,-15) -- (axis cs:0.2,-19)
    node[midway, above, sloped, font=\scriptsize] {ion acceleration};
  \node[font=\scriptsize, anchor=south] at (axis cs:0.0,-21.5) {wall};
  \node[font=\scriptsize, engred, anchor=south west] at (axis cs:0.05,-21.5)
    {$\;V_{\mathrm{wall}}<0$};
\end{axis}
\end{tikzpicture}
\caption[Plasma sheath potential]{The plasma sheath. Because electrons are
far more mobile than ions, any surface in contact with a plasma charges
negative until it repels enough electrons to balance the ion and electron
fluxes. The resulting potential drop is confined to a sheath a few Debye
lengths thick at the wall, while the bulk plasma stays quasineutral at
nearly constant potential. Ions entering the sheath are accelerated
toward the wall along the field gradient, producing the directional
bombardment that gives reactive-ion etching its anisotropy. The
exponential profile shown is a schematic; the exact sheath solution
follows the Child-Langmuir law in the high-voltage limit.}
\label{fig:vol04ch13:sheath-potential}
\end{figure}


The anisotropy comes from the sheath, shown in
Figure~\ref{fig:vol04ch13:sheath-potential}. Electrons are far more mobile
than ions, so any surface in contact with a plasma charges negative until
it repels enough electrons to balance the two fluxes. The resulting
potential drop is confined to a sheath a few Debye lengths thick at the
surface, while the bulk plasma stays quasineutral at nearly constant
potential. Ions crossing the sheath are accelerated normal to the surface
and arrive with a tightly directed velocity, so they etch straight down
rather than sideways. Control of the sheath voltage, through the RF bias
power, is the control of etch directionality; the same physics sets the
ion energy that drives sputter damage and the substrate heating an etch
process must manage.

\paragraph{Worked example: the Bohm criterion and ion bombardment energy.}
A stable sheath forms only if ions enter it already moving at or above the
\engterm{Bohm velocity}, the ion-acoustic speed $u_{B} = \sqrt{k_{B}T_e/
m_i}$. For an argon plasma ($m_i = 40\,\text{u} = 6.6\times 10^{-26}\,\si{
\kilogram}$) at $T_e = 3\,\si{\electronvolt}$ ($k_{B}T_e = 4.8\times
10^{-19}\,\si{\joule}$), $u_{B} = \sqrt{(4.8\times 10^{-19})/(6.6\times
10^{-26})} \approx 2.7\times 10^{3}\,\si{\meter\per\second}$. To reach this
speed the plasma drops a potential of about $k_{B}T_e/2e \approx 1.5\,\si{
\volt}$ across the quasineutral presheath. The ion then falls through the
full sheath potential, which for a floating wall is
\[
V_{\mathrm{sheath}} = \frac{k_{B}T_e}{2e}
+ \frac{k_{B}T_e}{2e}\ln\!\left(\frac{m_i}{2\pi m_e}\right).
\]
For argon, $\ln(m_i/2\pi m_e) = \ln(40\times 1836/2\pi) \approx \ln(1.17
\times 10^{4}) \approx 9.4$, so the floating sheath drop is $\approx
1.5(1 + 9.4) \approx 16\,\si{\volt}$ and an ion arrives at the wall with
roughly $16\,\si{\electronvolt}$. An applied RF bias deepens the sheath and
raises the bombardment energy to the hundreds of electronvolts that drive
anisotropic etching; the directional ion energy is the single knob that
separates a vertical-wall etch from an isotropic chemical one.

\engterm{Plasma deposition} (plasma-enhanced chemical vapour
deposition, PECVD) uses a similar low-pressure discharge to deposit
thin films of silicon nitride, silicon dioxide, amorphous silicon,
and other materials at lower temperature than conventional CVD. The
plasma activates the precursor molecules, allowing deposition at
$200$ to $400\,\si{\celsius}$ on substrates that could not survive
$800\,\si{\celsius}$ thermal CVD. PECVD is the standard process for
silicon-nitride passivation layers in integrated circuits and for
the active layers of thin-film solar cells. The plasma does the work that
temperature does in thermal CVD: electron impact cracks the stable precursor
into reactive radicals in the gas phase, so the film grows from a flux of
radicals that would need a much hotter surface to form thermally. The film
composition and stress follow from the ion bombardment the substrate sheath
imposes, so the same RF-bias knob that controls etch directionality here tunes
the density, refractive index, and residual stress of the deposited film, which
is why PECVD silicon nitride can be run compressive or tensile by choice of
bias and pressure.

\paragraph{Worked example: PECVD deposition rate from the radical flux.}
A silane-ammonia PECVD process grows silicon nitride from a radical flux that,
for a low-density capacitive discharge, is a small fraction of the neutral flux
striking the wafer. The neutral flux from a gas of density $n_g$ at thermal
speed $\bar{v}$ is $\Gamma_n = \tfrac14 n_g\bar{v}$. At a deposition pressure of
$60\,\si{\pascal}$ and a gas temperature of $600\,\si{\kelvin}$, the density is
$n_g = p/(k_B T) = 60/[(1.38\times 10^{-23})(600)] \approx 7.2\times
10^{21}\,\si{\per\meter\cubed}$, and the mean speed of a $\sim 30\,\text{u}$
precursor fragment is $\bar{v} = \sqrt{8k_B T/(\pi m)} \approx \sqrt{8(1.38
\times 10^{-23})(600)/[\pi(5\times 10^{-26})]} \approx 6.5\times 10^{2}\,\si{
\meter\per\second}$, so $\Gamma_n = \tfrac14(7.2\times 10^{21})(6.5\times
10^{2}) \approx 1.2\times 10^{24}\,\si{\per\meter\squared\per\second}$. If a
fraction $f \approx 10^{-3}$ of that flux is reactive film-forming radicals and
each deposits one silicon-nitride formula unit that adds a volume $\Omega
\approx 4\times 10^{-29}\,\si{\meter\cubed}$ to the film, the growth rate is
$\dot{z} = f\Gamma_n\Omega = (10^{-3})(1.2\times 10^{24})(4\times 10^{-29})
\approx 4.8\times 10^{-8}\,\si{\meter\per\second}$, about $48\,\si{\nano\meter}$
per second, or a few nanometres per second at the more typical percent-level
radical fraction reduced by re-emission. A $200\,\si{\nano\meter}$ passivation
layer therefore lands in tens of seconds to a minute, the throughput a
production tool is built around, and the rate is tuned through the RF power that
sets the radical fraction rather than through the substrate temperature that a
thermal process would raise.

\engterm{Plasma lighting} uses the radiation from excited atoms or
molecules in a discharge tube. Fluorescent lamps use a low-pressure
mercury discharge that emits $254\,\si{\nano\meter}$ ultraviolet
radiation; a phosphor coating on the tube wall converts the UV to
visible light. Modern compact fluorescent and tube-fluorescent lamps
achieve $\sim 80\,\si{\lumen\per\watt}$, compared to $\sim 15\,\si{
\lumen\per\watt}$ for incandescent. High-intensity discharge lamps
(HID: metal halide, high-pressure sodium) operate at higher pressure
($10$ to $100\,\si{\kilo\pascal}$) and produce broader spectra; they
are the standard for street lighting, stadiums, and industrial
spaces. LED lighting (Chapter 12) is displacing fluorescent and HID
in most applications since the 2010s; lighting is currently in a
technology transition from gas discharge to solid state.

The efficiency of a discharge lamp is an energy-budget calculation the
chapter's tools already support. The input electrical power heats the
electrons; the electrons excite the mercury atoms; the excited atoms
radiate the $254\,\si{\nano\meter}$ resonance line; the phosphor absorbs a
$254\,\si{\nano\meter}$ photon of $4.9\,\si{\electronvolt}$ and re-emits a
visible photon of about $2.3\,\si{\electronvolt}$. The energy lost in that
last down-conversion, the \engterm{Stokes deficit}, is unavoidable and
fixes a ceiling on the efficiency: even if every input joule reached the
phosphor as ultraviolet, the visible output could be at most $2.3/4.9
\approx 47\%$ of it in energy, and the eye's response then converts the
rest of the way to the $\sim 80\,\si{\lumen\per\watt}$ the lamp delivers
against the $\sim 683\,\si{\lumen\per\watt}$ of an ideal monochromatic
green source. The remaining losses are the electron-heating inefficiency,
the fraction of electron energy that goes to elastic collisions and
non-radiating states rather than to the resonance line, and the
absorption of the ultraviolet in the mercury vapour before it reaches the
wall, which is why the mercury pressure is tuned to an optimum: too little
and the excitation rate falls, too much and the resonance radiation is
trapped and reabsorbed. The high-intensity discharge lamp trades this
line-source physics for a high-pressure arc whose broadened and
continuum-filled spectrum gives better colour at the cost of a hotter,
higher-pressure envelope. The same Stokes-deficit accounting explains why
the light-emitting diode wins: a well-designed diode emits its visible
photon directly from a semiconductor transition matched to the desired
colour, with no ultraviolet intermediary and no down-conversion deficit,
so it escapes the ceiling the phosphor imposes on every gas discharge.

\engterm{Plasma propulsion} (electric propulsion, ion thrusters)
accelerates ionised propellant by electric or electromagnetic fields
to produce thrust. The Hall-effect thruster, the standard for modern
satellites, accelerates xenon ions to $\sim 30\,\si{\kilo\meter\per
\second}$ (specific impulse $\sim 3000\,\si{\second}$, compared to
$\sim 450\,\si{\second}$ for chemical bipropellants). The trade-off
is that thrust is small ($\sim 100\,\si{\milli\newton}$ for a
$5\,\si{\kilo\watt}$ Hall thruster) compared to chemical rockets;
the high specific impulse pays off only for long-duration missions
where propellant mass dominates the spacecraft mass budget. Hall
thrusters are now standard on geostationary satellites for
station-keeping and on deep-space probes; the NASA Psyche mission
(2023) uses them for primary propulsion. The third case study of this
chapter designs one end to end.

\section{Fusion engineering: tokamaks, stellarators, inertial confinement}\pathtag{standard}

Controlled fusion is the canonical long-horizon engineering project
of the late twentieth and early twenty-first centuries. The basic
requirement is to confine a hydrogen-isotope plasma at $\sim 10\,\si{
\kilo\electronvolt}$ temperature and $\sim 10^{20}\,\si{\per\meter
\cubed}$ density long enough that the fusion power exceeds the
losses, $n_{e}\tau_{E}T \ge 3\times 10^{21}\,\si{\meter\cubed\per
\kilo\electronvolt\second}$ (the Lawson criterion for D-T fusion).
Three confinement approaches are being engineered:

% Figure: Maxwell-averaged fusion reactivity <sigma v> versus ion
% temperature for D-T and D-D, on a log-log axis. Data in
% data/fusion_reactivity.csv; the triple-product analysis is in
% code/lawson_triple_product.py.
\begin{figure}[htbp]
\centering
\begin{tikzpicture}
\begin{loglogaxis}[
  width=0.78\textwidth, height=0.52\textwidth,
  xlabel={ion temperature $T$ (keV)},
  ylabel={reactivity $\langle\sigma v\rangle$ ($\mathrm{m^3\,s^{-1}}$)},
  xmin=1, xmax=100, ymin=1e-28, ymax=2e-21,
  grid=both, grid style={enggray!20},
  every axis plot/.append style={thick},
  legend pos=south east, legend cell align=left,
]
  \addplot[engred, mark=*, mark size=1.2pt] table[
    col sep=comma, x=temperature_keV, y=DT_sigmav_m3_s]
    {volumes/04-energy/ch13-plasma-physics/data/fusion_reactivity.csv};
  \addlegendentry{D--T}
  \addplot[engblue, mark=square*, mark size=1.2pt] table[
    col sep=comma, x=temperature_keV, y=DD_sigmav_m3_s]
    {volumes/04-energy/ch13-plasma-physics/data/fusion_reactivity.csv};
  \addlegendentry{D--D}
  \draw[dashed, enggray] (axis cs:15,1e-28) -- (axis cs:15,2e-21);
  \node[enggray, font=\scriptsize, anchor=south, rotate=90]
    at (axis cs:15,1e-25) {operating window};
\end{loglogaxis}
\end{tikzpicture}
\caption[Fusion reactivity versus temperature]{Maxwell-averaged fusion
reactivity $\langle\sigma v\rangle$ against ion temperature for the
deuterium-tritium (D--T) and deuterium-deuterium (D--D) reactions. D--T
exceeds D--D by roughly two orders of magnitude across the practical range
and peaks near $70\,\si{\kilo\electronvolt}$, which is why every
near-term fusion device burns D--T. The required Lawson triple product
$n\tau_E T$ for ignition is lowest near $10$ to $20\,\si{\kilo\electronvolt}$,
not at the reactivity peak, because the bremsstrahlung and confinement
penalties of running hotter outweigh the reactivity gain; this is the
operating window marked. The triple-product calculation is carried out in
\texttt{code/lawson\_triple\_product.py}.}
\label{fig:vol04ch13:fusion-reactivity}
\end{figure}


The choice of fuel and operating temperature follows from
Figure~\ref{fig:vol04ch13:fusion-reactivity}. The deuterium-tritium
reaction has a reactivity roughly two orders of magnitude above
deuterium-deuterium across the practical range, which is why every
near-term device burns D-T despite the cost and handling burden of
tritium. The reactivity peaks near $70\,\si{\kilo\electronvolt}$, but the
required triple product is lowest near $10$ to $20\,\si{\kilo\electronvolt}$,
because running hotter increases bremsstrahlung losses and stresses
confinement faster than it raises the reaction rate. The
\texttt{code/lawson\_triple\_product.py} tabulates the required triple
product against temperature and confirms the minimum of about $3\times
10^{21}\,\si{\meter\cubed\per\kilo\electronvolt\second}$ in that window.

A \engterm{tokamak} confines the plasma in a toroidal magnetic
geometry: a strong toroidal field produced by external coils, plus a
poloidal field produced by a plasma current driven by transformer
action. The combination produces helical field lines that the plasma
follows, with the $\vect{J}\times \vect{B}$ confinement keeping the
plasma off the walls. ITER, under construction at Cadarache, France,
with first plasma expected $\sim 2034$ and full D-T operation $\sim
2039$, will be the largest tokamak ever built: major radius $6.2\,\si{
\meter}$, plasma current $15\,\si{\mega\ampere}$, design fusion power
$500\,\si{\mega\watt}$ for $400\,\si{\second}$ pulses at $Q = 10$
(fusion power output / external heating power).

% Figure: two magnetic-confinement geometries. Left: a tokamak torus with
% toroidal coils and the helical field line. Right: a magnetic mirror with
% two high-field coils and a trapped particle bouncing between the throats.
\begin{figure}[htbp]
\centering
\begin{tikzpicture}[>=latex, scale=1.0]
% ===== Left: tokamak =====
\begin{scope}[shift={(0,0)}]
  \node[font=\small] at (0,2.7) {tokamak (toroidal)};
  % outer and inner torus ellipses
  \draw[enggray, thick] (0,0) ellipse (2.0 and 1.1);
  \draw[enggray, thick] (0,0) ellipse (0.9 and 0.45);
  % plasma ring (between)
  \draw[engred, very thick] (0,0) ellipse (1.45 and 0.78);
  % helical field hint on plasma ring
  \foreach \a in {20,70,120,160,200,250,300,340}{
    \node[engblue, font=\tiny] at ({1.45*cos(\a)},{0.78*sin(\a)}) {$\bullet$};
  }
  % a toroidal coil (small loop crossing the ring)
  \draw[engorange, thick] (1.45,0) ellipse (0.28 and 0.55);
  \draw[engorange, thick] (-1.45,0) ellipse (0.28 and 0.55);
  \node[engred, font=\scriptsize] at (0,-1.45) {plasma};
  \node[engorange, font=\scriptsize] at (1.95,-0.85) {coil};
\end{scope}
% ===== Right: magnetic mirror =====
\begin{scope}[shift={(6.3,0)}]
  \node[font=\small] at (0,2.7) {magnetic mirror};
  % field lines bulging out in the middle, pinched at the throats
  \foreach \s in {-1,1}{
    \draw[engblue, thick]
      (-2.0,{\s*0.35}) .. controls (-0.8,{\s*1.0}) and (0.8,{\s*1.0}) ..
      (2.0,{\s*0.35});
  }
  % coils at the throats
  \draw[engorange, very thick] (-2.0,-0.7) -- (-2.0,0.7);
  \draw[engorange, very thick] (2.0,-0.7) -- (2.0,0.7);
  \node[engorange, font=\scriptsize] at (-2.0,-1.05) {$B_{\max}$};
  \node[engorange, font=\scriptsize] at (2.0,-1.05) {$B_{\max}$};
  \node[engblue, font=\scriptsize] at (0,-1.25) {$B_{\min}$};
  % trapped particle bouncing
  \draw[engred, thick, ->] (-1.1,0) .. controls (-0.3,0.5) and (0.3,0.5) .. (1.1,0);
  \draw[engred, thick, ->] (1.1,0.15) .. controls (0.3,-0.4) and (-0.3,-0.4) .. (-1.1,0.15);
  \node[engred, font=\scriptsize] at (0,1.0) {trapped};
\end{scope}
\end{tikzpicture}
\caption[Tokamak and magnetic-mirror geometries]{Two magnetic-confinement
geometries. Left: a tokamak confines the plasma in a torus; external coils
make the strong toroidal field, and a driven plasma current adds the
poloidal field, so the net field lines spiral helically and the
$\vect{J}\times\vect{B}$ force holds the plasma off the wall. Right: a
magnetic mirror traps particles between two high-field throats. A particle
moving toward rising $B$ converts parallel into perpendicular energy by
conservation of the magnetic moment $\mu$, and reflects if its pitch angle
is steep enough; particles inside the loss cone escape.}
\label{fig:vol04ch13:tokamak-mirror}
\end{figure}


Figure~\ref{fig:vol04ch13:tokamak-mirror} contrasts the toroidal
confinement of a tokamak with the open-ended trapping of a magnetic
mirror. The tokamak closes the field lines on themselves, eliminating the
loss-cone leak of the mirror at the cost of needing a driven plasma
current and the disruption risk that current carries.

A \engterm{stellarator} confines the plasma in a similar toroidal
geometry but with the helical field lines produced entirely by
external coils, eliminating the need for a plasma current. The
Wendelstein 7-X at Greifswald, Germany, achieved first plasma in
2015 and is the largest operating stellarator. The trade-off versus
tokamaks is that stellarators have more complex coil geometry but no
disruption risk (no plasma current to interrupt) and steadier
operation. The design cost is paid up front in geometry. A stellarator
generates its rotational transform by twisting and shaping the plasma
column in three dimensions, so its coils are not the simple planar rings
of a tokamak but complex non-planar shapes that must be built and
positioned to millimetre tolerance over a device metres across; the
Wendelstein 7-X coils were the product of a decade of optimisation
against a computed magnetic configuration, and a fabrication error of a
few millimetres would have spoiled the confinement the optimisation
bought. The reward is that the confinement geometry is fixed by the
hardware rather than sustained by a driven current, so the machine runs
in steady state without the transformer flux swing that pulses a
tokamak, and it carries no free current energy to drive a disruption.
The bootstrap current the neoclassical mastery box named is a
complication here rather than a resource: a stellarator is designed to
minimise the self-driven current that a tokamak seeks to maximise,
because any large current would reintroduce the current-driven
instabilities the stellarator exists to avoid. The two lines are a
standing trade in fusion engineering between the tokamak's simpler coils
and harder operation and the stellarator's harder coils and simpler
operation.

\engterm{Inertial confinement} uses laser or particle-beam pulses to
compress a fuel pellet (a millimetre-scale capsule of frozen
deuterium-tritium) to thousand-fold density and ignition temperature
on the nanosecond time scale. The National Ignition Facility at
Lawrence Livermore achieved net-positive fusion gain in December
2022 (fusion energy out exceeding laser energy on target), the first
demonstration of any approach reaching this milestone. The confinement
in this scheme is not magnetic but inertial: the fuel is compressed and
heated so fast that its own mass holds it together for the nanosecond it
takes to burn, before it flies apart. The relevant confinement time is
the time a compressed sphere of radius $r$ holds together against its own
expansion at the ion-sound speed, $\tau \sim r/c_s$, so the Lawson
product $n\tau$ is met not by a long confinement time at modest density,
as in a tokamak, but by an enormous density for a vanishingly short time,
the opposite corner of the same $n\tau T$ requirement. Reaching the
thousand-fold compression demands that the drive be applied with
extraordinary symmetry, because any asymmetry seeds a Rayleigh-Taylor
instability at the imploding shell that mixes cold fuel into the hot spot
and quenches the burn, which is why the drive is smoothed either by
firing the lasers into a radiation-filled cavity that bathes the pellet
in uniform x-rays, the indirect-drive approach NIF uses, or by exquisite
beam balancing in direct drive. Production of electricity from inertial
confinement remains an open engineering problem: the physics
demonstration of gain on a single shot is separated from a power plant by
the need to fire several such shots a second, to manufacture and inject a
cryogenic target on every shot, and to convert the pulsed neutron and
debris output into steady electricity, none of which NIF, a single-shot
physics facility, was built to do.

% Figure: the plant recirculating-power fraction against fusion gain Q. A
% power plant must run its own magnets, heating, and balance of plant from a
% slice of its own output; that recirculating fraction falls as 1/Q (times a
% conversion and heating-efficiency factor), so a plant needs Q well above the
% break-even Q=1 to sell net electricity. The curve marks the ITER Q=10 point
% and the plant-relevant Q~30-40 band.
\begin{figure}[htbp]
\centering
\begin{tikzpicture}
\begin{axis}[
  width=0.82\textwidth, height=0.55\textwidth,
  xlabel={fusion gain $Q=P_{\mathrm{fus}}/P_{\mathrm{heat}}$},
  ylabel={recirculating power fraction $f_{\mathrm{rec}}$},
  xmin=1, xmax=50, ymin=0, ymax=1.0,
  grid=both, grid style={enggray!20},
  every axis plot/.append style={thick},
]
  % f_rec = 1/(eta_th * (1 + Q/5) * eta_heat) approx; use eta_th=0.4,
  % eta_heat=0.7, and gain in electricity M ~ (1 + Q/5) for alpha+neutron.
  % f_rec = 1 / (0.4 * 0.7 * (1 + Q/5))
  \addplot[engblue, domain=1:50, samples=80]
    {1/(0.4*0.7*(1 + x/5))};
  % ITER Q=10 marker
  \addplot[only marks, mark=*, engred, mark size=2pt] coordinates {(10,1.19)};
  \draw[dashed, enggray] (axis cs:1,1) -- (axis cs:50,1);
  \node[enggray, font=\scriptsize, anchor=south east] at (axis cs:50,1.0)
    {break-even in electricity};
  \node[engred, font=\scriptsize, anchor=west] at (axis cs:10,0.85)
    {ITER $Q=10$};
  % plant band
  \draw[engorange, thick, |-|] (axis cs:30,0.15) -- (axis cs:40,0.15);
  \node[engorange, font=\scriptsize, anchor=north] at (axis cs:35,0.15)
    {plant band};
\end{axis}
\end{tikzpicture}
\caption[Recirculating power against fusion gain]{The fraction of its own
output a fusion plant must recirculate to run its magnets, heating, and
balance of plant, drawn against the fusion gain $Q$. The estimate takes a
thermal-conversion efficiency near $0.4$, a heating-injection efficiency near
$0.7$, and an electricity multiplication set by the alpha and neutron energy,
so the recirculating fraction falls roughly as the inverse of the gain. At the
$Q=10$ ITER point the plant would still spend more than its whole output on
itself, which is why ITER is a physics-gain demonstration and not a power
plant; a plant that sells electricity needs $Q$ in the thirties or forties,
where the recirculating fraction drops to a fifth or less of the output. The
gap between the $Q=10$ that demonstrates the burn and the $Q\sim 40$ that
sells power is the engineering distance the fusion programme still has to
cover.}
\label{fig:vol04ch13:fusion-gain}
\end{figure}


\paragraph{Worked example: fusion gain and the recirculating power a plant must pay.}
The gain $Q = P_{\mathrm{fus}}/P_{\mathrm{heat}}$ that a physicist quotes is not
the number that sells electricity; a plant must run its own magnets, heating,
and cooling from a slice of its own output, and that recirculating fraction is
the figure that decides commercial viability. Trace the power flow. A plant with
fusion power $P_{\mathrm{fus}}$ converts the thermal output to electricity at an
efficiency $\eta_{\mathrm{th}} \approx 0.4$; the neutrons multiply the thermal
yield in the blanket, so the electricity generated is roughly $P_e \approx
\eta_{\mathrm{th}}(P_{\mathrm{fus}} + P_{\mathrm{heat}})$. The heating itself is
injected at a wall-plug efficiency $\eta_{\mathrm{heat}} \approx 0.7$, so the
electricity it consumes is $P_{\mathrm{heat}}/\eta_{\mathrm{heat}} =
P_{\mathrm{fus}}/(Q\eta_{\mathrm{heat}})$. The recirculating fraction is the
ratio of the power the plant spends on its own heating to the power it makes,
$f_{\mathrm{rec}} \approx 1/[\eta_{\mathrm{th}}\eta_{\mathrm{heat}}(1 + Q/5)]$,
using the alpha-plus-neutron accounting that puts a fifth of the fusion energy
into the plasma. At the ITER point $Q = 10$, $f_{\mathrm{rec}} \approx 1/[(0.4)
(0.7)(3)] \approx 1.2$, greater than one, so ITER at $Q = 10$ would spend more
than its whole output on itself, which is why it is a gain demonstration and not
a power plant. To drop the recirculating fraction to a fifth of output demands
$1 + Q/5 \approx 1/(0.28\times 0.2) \approx 18$, so $Q \approx 85$ under this
pessimistic accounting, and the plant-relevant band of $Q \sim 30$ to $40$ in
Figure~\ref{fig:vol04ch13:fusion-gain} rests on better conversion and heating
efficiencies. The distance from the $Q = 10$ that demonstrates the burn to the
$Q$ a plant needs is the engineering gap the fusion programme still has to close.

% Figure: the scrape-off-layer flux tube striking the divertor at a shallow
% angle. The parallel heat flux carried along the field is enormous, but the
% field meets the target at a grazing angle so the flux is spread over a
% larger footprint, cutting the surface load by the sine of the angle. The
% sketch shows the last closed flux surface, the narrow scrape-off layer of
% width lambda_q, and the tilted target that spreads the load.
\begin{figure}[htbp]
\centering
\begin{tikzpicture}[>=Latex, scale=1.0]
  % last closed flux surface (plasma edge)
  \draw[engblue, thick] (0,0.6) .. controls (2.5,1.4) and (5.0,1.2) .. (7.0,0.4);
  \node[engblue, font=\scriptsize, anchor=south] at (2.0,1.25)
    {last closed flux surface};
  % scrape-off layer, a thin band just outside
  \draw[engorange, thick] (0,0.35) .. controls (2.5,1.15) and (5.0,0.95) .. (7.0,0.15);
  \draw[<->, enggray] (7.15,0.15) -- (7.15,0.4);
  \node[enggray, font=\scriptsize, anchor=west] at (7.2,0.28)
    {$\lambda_q$};
  \node[engorange, font=\scriptsize, anchor=south] at (5.4,0.95)
    {scrape-off layer};
  % field line spiralling down into the divertor along the SOL
  \draw[engred, thick, ->] (6.6,0.25) .. controls (7.2,-0.4) and (7.0,-1.1) .. (6.2,-1.7);
  \node[engred, font=\scriptsize, anchor=west] at (7.0,-0.8)
    {$q_\parallel$};
  % divertor target plate, tilted at a shallow angle
  \draw[engblack, very thick] (4.3,-2.3) -- (7.7,-1.3);
  \node[engblack, font=\scriptsize, anchor=north] at (6.0,-2.0)
    {divertor target};
  % grazing angle arc
  \draw[enggray] (6.2,-1.7) -- (6.9,-2.15);
  \draw[enggray, ->] (6.7,-1.55) arc (250:290:0.6);
  \node[enggray, font=\scriptsize] at (6.55,-1.35) {$\alpha$};
  % footprint bracket on the plate
  \draw[decorate, decoration={brace, amplitude=4pt}, enggray]
    (5.9,-2.35) -- (7.1,-1.98);
  \node[enggray, font=\scriptsize, anchor=north] at (6.5,-2.45)
    {spread footprint};
\end{tikzpicture}
\caption[Divertor heat load and the grazing target]{The scrape-off layer
carries the plasma exhaust along the field to the divertor target. The heat
flux parallel to the field, $q_\parallel$, is enormous because the power that
crosses the last closed flux surface is funnelled into a scrape-off layer only
a few millimetres wide, $\lambda_q$. The target is set at a shallow grazing
angle $\alpha$ to the field so that the parallel flux is projected onto a
larger surface footprint, and the load the material sees is
$q_\parallel\sin\alpha$, smaller than the parallel flux by the sine of the
grazing angle. Tilting the target and widening the wetted footprint are the
two geometric knobs a divertor designer has; when they are not enough, the
exhaust power is radiated away volumetrically before it reaches the plate by
seeding an impurity that lights up in the scrape-off layer, the detached
divertor regime.}
\label{fig:vol04ch13:divertor-flux}
\end{figure}


\paragraph{Worked example: the divertor parallel heat flux and the grazing angle.}
The exhaust power that crosses the plasma edge is funnelled into a scrape-off
layer only millimetres wide and carried along the field to the divertor, where it
becomes the single most demanding heat-transfer problem in a fusion reactor.
Take the case-study tokamak exhausting $P_{\mathrm{exh}} = 100\,\si{\mega\watt}$
into a scrape-off layer of width $\lambda_q = 3\,\si{\milli\meter}$ wrapped once
around the machine at major radius $R = 6\,\si{\meter}$, so the area through
which the parallel flux flows, projected onto the poloidal circumference, is $A_\parallel
= 2\pi R\lambda_q\sin\vartheta$ with a field-line pitch $\sin\vartheta \approx
B_\theta/B_\phi \approx 0.1$. The wetted parallel area is then $A_\parallel
\approx 2\pi(6)(3\times 10^{-3})(0.1) \approx 1.1\times 10^{-2}\,\si{\meter
\squared}$, so the parallel heat flux is $q_\parallel = P_{\mathrm{exh}}/
A_\parallel \approx 10^{8}/1.1\times 10^{-2} \approx 9\times 10^{9}\,\si{\watt
\per\meter\squared}$, nine gigawatts a square metre, a flux no solid material
can accept. The rescue is geometry: the target is tilted to a grazing angle
$\alpha \approx 2^{\circ}$ to the field, Figure~\ref{fig:vol04ch13:divertor-flux},
so the flux the surface sees is $q_\perp = q_\parallel\sin\alpha \approx
(9\times 10^{9})(0.035) \approx 3\times 10^{8}\,\si{\watt\per\meter\squared}$,
still above the $\sim 10\,\si{\mega\watt\per\meter\squared}$ a water-cooled
tungsten target can survive by an order of magnitude. Closing the remaining
factor forces the detached-divertor regime, in which a seeded impurity radiates
most of the exhaust power volumetrically in the scrape-off layer before it
reaches the plate, spreading the load over the whole divertor chamber rather
than the narrow strike line. The divertor heat-exhaust problem, not the core
confinement, is widely held to be the hardest engineering obstacle between a
burning plasma and a power plant.

\begin{mastery}
The Lawson criterion follows from a power balance that any reactor engineer
can reconstruct. In a deuterium-tritium plasma of density $n$ (taking
$n_D = n_T = n/2$) at temperature $T$, the volumetric fusion power released
in charged products that heat the plasma is $P_\alpha = \tfrac14 n^{2}
\langle\sigma v\rangle E_\alpha$, where $E_\alpha = 3.5\,\si{\mega
\electronvolt}$ is the alpha-particle energy and the neutron carries the
remaining $14.1\,\si{\mega\electronvolt}$ out of the plasma. The plasma
stores thermal energy $W = 3nk_{B}T$ (a factor of three from two species
each contributing $\tfrac32 k_{B}T$) and loses it on the energy confinement
time $\tau_E$, so the conduction loss is $P_{\mathrm{cond}} = 3nk_{B}T/
\tau_E$. Ignition is the condition that alpha self-heating balances the
loss, $P_\alpha \ge P_{\mathrm{cond}}$:
\[
\tfrac14 n^{2}\langle\sigma v\rangle E_\alpha \ge \frac{3nk_{B}T}{\tau_E}
\quad\Longrightarrow\quad
n\tau_E \ge \frac{12 k_{B}T}{\langle\sigma v\rangle E_\alpha}.
\]
Multiplying both sides by $T$ gives the \engterm{triple product}
$n T \tau_E \ge 12 k_{B}T^{2}/(\langle\sigma v\rangle E_\alpha)$. Over the
$10$ to $20\,\si{\kilo\electronvolt}$ window the reactivity $\langle\sigma
v\rangle$ rises faster than $T^{2}$, so the right-hand side has a minimum
there, $n T \tau_E \approx 3\times 10^{21}\,\si{\meter\cubed\per\kilo
\electronvolt\second}$ for D-T. The triple product, not any single one of
$n$, $T$, or $\tau_E$, is the figure of merit a confinement scheme must
raise, and the choice of operating temperature follows from minimising it
rather than from maximising the bare reactivity.
\end{mastery}

\section{An end-to-end case study: a tokamak power balance}\pathtag{standard}

The chapter's apparatus pays off when an engineer must judge whether a
proposed magnetic-confinement reactor can reach a self-sustaining burn. We
take a tokamak sized between ITER and a notional power plant and carry it
from its density, temperature, and confinement time through the Lawson
criterion and the triple product, then through the competition between
fusion heating and the two dominant loss channels, to an ignition margin
and a verdict on whether the machine as specified would burn. The case is a
good vehicle because every operational formula of the chapter appears in it
once: the plasma parameters, the fusion reactivity, the magnetic and
thermal pressures, the confinement scaling, and the power balance that
decides the project.

\subsection{The machine and its operating point}

Take a tokamak with major radius $R = 6\,\si{\meter}$, minor radius $a =
2\,\si{\meter}$, on-axis toroidal field $B = 5.3\,\si{\tesla}$, and a
plasma volume approximated as a torus, $V = 2\pi^{2}R a^{2} = 2\pi^{2}(6)
(2)^{2} \approx 470\,\si{\meter\cubed}$. The design fills this with a 50:50
deuterium-tritium mixture at electron density $n = 1.0\times 10^{20}\,\si{
\per\meter\cubed}$ and volume-averaged temperature $T = 12\,\si{\kilo
\electronvolt}$, with an energy confinement time $\tau_E = 3.0\,\si{
\second}$ taken from the device's confinement scaling. These four numbers,
density, temperature, confinement time, and volume, fix the entire balance.

A first check confirms the plasma is magnetically dominated. The thermal
pressure is $p = 2nk_{B}T$ (electrons plus ions), with $k_{B}T = 12\times
10^{3}\times 1.6\times 10^{-19} = 1.92\times 10^{-15}\,\si{\joule}$, so $p =
2(1.0\times 10^{20})(1.92\times 10^{-15}) \approx 3.8\times 10^{5}\,\si{
\pascal}$. The magnetic pressure is $p_B = B^{2}/(2\mu_0) = (5.3)^{2}/
(2\times 4\pi\times 10^{-7}) \approx 1.1\times 10^{7}\,\si{\pascal}$, so the
plasma beta is $\beta = p/p_B \approx 0.034$, a few percent, the expected
magnetically confined regime. The field holds the plasma with two orders of
magnitude of pressure margin.

\subsection{Lawson criterion and triple product}

The triple product of the operating point is
\[
n T \tau_E = (1.0\times 10^{20})(12)(3.0)
= 3.6\times 10^{21}\,\si{\meter\cubed\per\kilo\electronvolt\second}.
\]
The ignition threshold for D-T near this temperature is $\approx 3\times
10^{21}\,\si{\meter\cubed\per\kilo\electronvolt\second}$ (the minimum of the
curve in Figure~\ref{fig:vol04ch13:fusion-reactivity} and the mastery box
above). The operating point sits a factor $1.2$ above threshold, so the
machine clears the Lawson bar with a small margin. A reactor designer reads
this margin as thin: a $20\%$ shortfall in confinement time, well within the
scatter of confinement scaling laws, would drop the plasma below ignition.

\subsection{Fusion power against the loss channels}

The fusion power balance carries the verdict. At $T = 12\,\si{\kilo
\electronvolt}$ the D-T reactivity is $\langle\sigma v\rangle \approx
2.6\times 10^{-22}\,\si{\meter\cubed\per\second}$ (from the reactivity
table). With $n_D = n_T = n/2 = 5\times 10^{19}\,\si{\per\meter\cubed}$, the
volumetric fusion reaction rate is
\[
\mathcal{R} = n_D n_T \langle\sigma v\rangle
= (5\times 10^{19})^{2}(2.6\times 10^{-22})
\approx 6.5\times 10^{17}\,\si{\per\meter\cubed\per\second}.
\]
Each reaction releases $17.6\,\si{\mega\electronvolt} = 2.82\times 10^{-12}
\,\si{\joule}$, so the total fusion power density is $\mathcal{R}\times
2.82\times 10^{-12} \approx 1.8\times 10^{6}\,\si{\watt\per\meter\cubed}$,
and over the $470\,\si{\meter\cubed}$ volume the gross fusion power is
$P_{\mathrm{fus}} \approx 8.6\times 10^{8}\,\si{\watt}$, about
$860\,\si{\mega\watt}$. The alpha particles retain one fifth of this,
$P_\alpha \approx 170\,\si{\mega\watt}$, available to heat the plasma; the
neutrons carry $\sim 690\,\si{\mega\watt}$ to the blanket for electricity
and tritium breeding.

The losses compete against the alpha heating. The conduction loss is set by
the confinement time,
\[
P_{\mathrm{cond}} = \frac{3nk_{B}T\,V}{\tau_E}
= \frac{3(1.0\times 10^{20})(1.92\times 10^{-15})(470)}{3.0}
\approx 9.0\times 10^{7}\,\si{\watt},
\]
about $90\,\si{\mega\watt}$. The bremsstrahlung loss for a hydrogenic plasma
is $P_{\mathrm{br}} = 1.7\times 10^{-38}\,n^{2}\sqrt{T/\text{eV}}\,V$ in
watts, $= 1.7\times 10^{-38}(1.0\times 10^{20})^{2}\sqrt{1.2\times 10^{4}}
(470) \approx 8.8\times 10^{5}(470) \approx 9\times 10^{6}\,\si{\watt}$,
about $9\,\si{\mega\watt}$, an order of magnitude below conduction at this
temperature. Total loss is $P_{\mathrm{loss}} \approx 99\,\si{\mega\watt}$.

% Figure: tokamak power balance. Fusion heating power and the two main loss
% channels (bremsstrahlung and conduction) versus ion temperature at fixed
% density and confinement time. The ignition point is where fusion self-
% heating overtakes total losses; the analysis is in the end-to-end case
% study and in code/lawson_triple_product.py.
\begin{figure}[htbp]
\centering
\begin{tikzpicture}
\begin{semilogyaxis}[
  width=0.80\textwidth, height=0.54\textwidth,
  xlabel={ion temperature $T$ (keV)},
  ylabel={volumetric power density ($\mathrm{MW\,m^{-3}}$)},
  xmin=2, xmax=40, ymin=0.02, ymax=5,
  grid=both, grid style={enggray!20},
  every axis plot/.append style={thick},
  legend pos=south east, legend cell align=left,
]
  % alpha self-heating: rises steeply then rolls over with reactivity
  \addplot[engred, domain=2:40, samples=80]
    {0.9*exp(-19/(x+3))*x^0.3};
  \addlegendentry{alpha self-heating $P_\alpha$}
  % bremsstrahlung: scales as sqrt(T)
  \addplot[engblue, domain=2:40, samples=80]
    {0.05*sqrt(x)};
  \addlegendentry{bremsstrahlung $P_{\mathrm{br}}$}
  % conduction loss: fixed by confinement time, roughly linear in T
  \addplot[engorange, domain=2:40, samples=80]
    {0.045*x};
  \addlegendentry{conduction loss $P_{\mathrm{cond}}$}
  \draw[dashed, enggray] (axis cs:10,0.02) -- (axis cs:10,5);
  \node[enggray, font=\scriptsize, anchor=south, rotate=90]
    at (axis cs:10,0.6) {operating point};
\end{semilogyaxis}
\end{tikzpicture}
\caption[Tokamak power balance]{Volumetric power densities in a
deuterium-tritium tokamak at fixed density and confinement time, plotted
against ion temperature. The alpha self-heating power $P_\alpha$ tracks the
fusion reactivity and rises steeply through the $5$ to $15\,\si{\kilo
\electronvolt}$ range. Bremsstrahlung radiation $P_{\mathrm{br}}$ grows
only as $\sqrt{T}$, and the conduction loss $P_{\mathrm{cond}}$ set by the
energy confinement time grows roughly linearly with $T$. Ignition is
reached where $P_\alpha$ overtakes the sum of the two loss channels; below
that crossing the plasma needs external heating, above it the fusion
alphas sustain the burn. The operating window near $10\,\si{\kilo
\electronvolt}$ marked here is where the required triple product is lowest.
The curves are schematic; the worked power balance is carried in the
end-to-end case study and in
\texttt{code/lawson\_triple\_product.py}.}
\label{fig:vol04ch13:power-balance}
\end{figure}


Figure~\ref{fig:vol04ch13:power-balance} shows why the operating point is
chosen where it is. Alpha self-heating climbs steeply through the working
window while bremsstrahlung grows only as $\sqrt{T}$ and conduction grows
linearly, so the crossing where self-heating overtakes total loss falls
near $10\,\si{\kilo\electronvolt}$, exactly where the triple product is
cheapest.

\subsection{Ignition margin and verdict}

The alpha heating $P_\alpha \approx 170\,\si{\mega\watt}$ exceeds the total
loss $P_{\mathrm{loss}} \approx 99\,\si{\mega\watt}$, so the plasma is
self-heating with a margin $P_\alpha/P_{\mathrm{loss}} \approx 1.7$. The
fusion gain referenced to external heating is the engineering figure that
decides the project. If the auxiliary heating needed to hold the operating
point against the loss not covered by alphas is $P_{\mathrm{aux}} =
P_{\mathrm{loss}} - P_\alpha$, this is negative, meaning the alphas more
than cover the losses and the plasma would run away toward higher
temperature until a transport or stability limit caps it; the machine is
nominally ignited. Referenced instead to the heating that brought it to the
operating point, the ratio $Q = P_{\mathrm{fus}}/P_{\mathrm{aux}}$ formally
diverges at ignition, which is the limiting case of the $Q = 10$ target
quoted for ITER.

The verdict is favourable but not comfortable. The machine ignites at the
nominal point, yet the Lawson margin of $1.2$ and the alpha-to-loss margin
of $1.7$ both sit close enough to unity that the realistic uncertainties in
confinement scaling, impurity radiation, and profile shape could erase
them. A reactor designer would respond by raising the confinement time
through size or field rather than by running hotter, since the
bremsstrahlung and conduction penalties of higher temperature climb faster
than the reactivity gain. The case study reproduces, at textbook scale, the
argument that sets the size of every burning-plasma device: the triple
product must clear threshold with enough margin to survive the scatter in
the confinement scaling, and that margin, not the peak reactivity, sets the
reactor scale.

\section{A second case study: a capacitively coupled etch reactor}\pathtag{standard}

The tokamak case carried the chapter's apparatus through the
highest-energy plasma an engineer designs. The opposite corner of the
density-temperature map, the low-temperature processing plasma that etches
every integrated circuit, exercises the same operational formulas at room
scale and decides a different project: whether a proposed reactor recipe
will cut the vertical, selective, uniform trenches a fabrication line
requires. We take a capacitively coupled plasma (CCP) reactor end to end,
from the radio-frequency power and the self-bias it builds, through the
sheath voltage, the ion energy and flux, to the etch rate, the selectivity,
the anisotropy, and the loading effect that ties the etch rate to the open
area on the wafer. Every number is carried so that a process engineer could
reproduce the chain and find the knob that fixes a failing etch. The
reactor geometry and the etch trends are drawn in
Figure~\ref{fig:vol04ch13:ccp-reactor} and
Figure~\ref{fig:vol04ch13:etch-bias}.

% Figure: capacitively coupled plasma (CCP) etch reactor. Schematic of the
% parallel-plate chamber, the RF-driven powered (lower) electrode carrying
% the wafer, the grounded (upper) showerhead electrode, the quasineutral
% bulk plasma, and the two sheaths. The labels carry the operating numbers
% used in the end-to-end CCP case study.
\begin{figure}[htbp]
\centering
\begin{tikzpicture}[>=Stealth, font=\small]
  % chamber wall
  \draw[enggray, thick] (0,0) rectangle (8,5.4);
  % upper grounded electrode (showerhead)
  \fill[enggray!35] (0.6,4.7) rectangle (7.4,5.0);
  \node[font=\scriptsize, anchor=south] at (4.0,5.0)
    {grounded showerhead electrode (gas inlet)};
  % gas-feed arrows
  \foreach \x in {1.6,3.2,4.8,6.4}{
    \draw[engblue, ->] (\x,4.7) -- (\x,4.3);}
  % bulk plasma region
  \fill[englime!22] (0.6,1.3) rectangle (7.4,4.3);
  \node[anchor=center, align=center] at (4.0,3.1)
    {quasineutral bulk plasma\\
      $n_e\sim 10^{16}\,\mathrm{m^{-3}}$, $T_e\sim 3\,\mathrm{eV}$};
  % upper sheath
  \fill[engorange!25] (0.6,4.3) rectangle (7.4,4.45);
  \node[font=\scriptsize, anchor=west] at (7.5,4.37) {ground sheath};
  % lower (driven) sheath, thicker
  \fill[engred!22] (0.6,1.0) rectangle (7.4,1.3);
  \node[font=\scriptsize, anchor=west] at (7.5,1.15) {RF sheath};
  % wafer on powered electrode
  \fill[engblue!30] (1.6,0.7) rectangle (6.4,1.0);
  \node[font=\scriptsize, anchor=north] at (4.0,0.68) {wafer};
  % powered electrode
  \fill[enggray!50] (0.6,0.35) rectangle (7.4,0.7);
  \node[font=\scriptsize, anchor=north] at (4.0,0.33)
    {powered electrode};
  % directional ion flux arrows through RF sheath onto wafer
  \foreach \x in {2.2,3.2,4.0,4.8,5.8}{
    \draw[engred, ->, thick] (\x,1.6) -- (\x,1.05);}
  \node[font=\scriptsize, engred, anchor=west] at (6.7,1.7)
    {ions $\Gamma_i$};
  % RF generator
  \draw[engblack, thick] (4.0,0.35) -- (4.0,-0.6);
  \draw[engblack, thick] (3.5,-0.6) rectangle (4.5,-1.4);
  \node at (4.0,-1.0) {$\sim$};
  \node[font=\scriptsize, anchor=north] at (4.0,-1.45)
    {RF $13.56\,\mathrm{MHz}$, $P_{\mathrm{rf}}$};
  % blocking capacitor / self-bias
  \draw[engblack, thick] (4.0,-0.05) -- (3.7,-0.05);
  \draw[engblack, thick] (3.7,0.10) -- (3.7,-0.20);
  \draw[engblack, thick] (3.55,0.10) -- (3.55,-0.20);
  \node[font=\scriptsize, anchor=east] at (3.5,-0.05) {$C_b$};
\end{tikzpicture}
\caption[Capacitively coupled plasma etch reactor]{A capacitively coupled
plasma (CCP) reactor for reactive-ion etching. Radio-frequency power at
$13.56\,\si{\mega\hertz}$ drives the lower electrode that carries the wafer;
the upper showerhead electrode is grounded and admits the process gas. The
bulk plasma stays quasineutral, and the entire applied potential drops
across two thin sheaths, one at each electrode. The blocking capacitor
$C_b$ in series with the powered electrode lets the electrode charge to a
negative direct-current self-bias, so the time-averaged sheath at the wafer
is much deeper than the sheath at the grounded wall. Ions accelerated across
that deep sheath strike the wafer along the surface normal, which is the
origin of anisotropic etching. The operating numbers shown anchor the
end-to-end CCP case study.}
\label{fig:vol04ch13:ccp-reactor}
\end{figure}


\subsection{The reactor and its operating point}

Take a parallel-plate CCP reactor with an electrode diameter of $300\,\si{
\milli\meter}$ (so the electrode area is $A = \pi(0.15)^{2} \approx
0.071\,\si{\meter\squared}$), an electrode gap of $30\,\si{\milli\meter}$,
fed with fluorine-bearing gas at a pressure of $5\,\si{\pascal}$ (about
$40\,\text{mtorr}$) and driven at $13.56\,\si{\mega\hertz}$ with a delivered
power of $P_{\mathrm{rf}} = 300\,\si{\watt}$. The discharge sustains a bulk
plasma with electron temperature $T_e = 3\,\si{\electronvolt}$ and bulk
density $n_e = 2\times 10^{16}\,\si{\per\meter\cubed}$, both in the range a
Langmuir probe would return for this class of tool. These few numbers,
power, pressure, area, gap, temperature, and density, fix the entire etch
balance.

A first check places the discharge in its operating regime. The Debye
length is $\lambda_D = \sqrt{\varepsilon_0 k_B T_e/(n_e e^{2})} =
\sqrt{(8.85\times 10^{-12})(4.8\times 10^{-19})/[(2\times 10^{16})(1.6
\times 10^{-19})^{2}]} \approx 9\times 10^{-5}\,\si{\meter}$, about
$90\,\si{\micro\meter}$, far smaller than the $30\,\si{\milli\meter}$ gap,
so the bulk is a genuine quasineutral plasma bounded by thin sheaths. The
ionisation fraction is low: at $5\,\si{\pascal}$ and room temperature the
neutral density is $n_n = p/(k_B T) \approx 5/[(1.38\times 10^{-23})(300)]
\approx 1.2\times 10^{21}\,\si{\per\meter\cubed}$, so the ionisation
fraction is $n_e/n_n \approx 2\times 10^{16}/1.2\times 10^{21} \approx
2\times 10^{-5}$, the weakly ionised regime in which most of the gas is
neutral and the reactive chemistry runs on radicals rather than on ions.

\subsection{Self-bias and the sheath voltage}

The blocking capacitor in series with the powered electrode is what makes a
CCP etch tool work. Because the electrode cannot pass direct current, it
charges to a negative self-bias $V_{\mathrm{dc}}$ that makes the
time-averaged electron and ion currents balance over an RF cycle. For an
asymmetric reactor, in which the powered electrode is smaller than the
grounded chamber, the self-bias approaches the peak RF voltage. We take the
RF voltage amplitude that delivers $300\,\si{\watt}$ to this discharge to be
$V_{\mathrm{rf}} \approx 400\,\si{\volt}$, so the self-bias settles near
$V_{\mathrm{dc}} \approx -350\,\si{\volt}$. The time-averaged voltage across
the wafer sheath is then of order $\overline{V}_s \approx 350\,\si{\volt}$,
two orders of magnitude above the $\sim 16\,\si{\volt}$ floating sheath of
an unbiased wall computed earlier in the chapter. The deep, driven sheath is
the entire point: it is the accelerator that turns thermal ions into the
directed beam that etches straight down.

\subsection{Ion energy and flux at the wafer}

Ions enter the sheath at the Bohm velocity, $u_B = \sqrt{k_B T_e/m_i}$. For
fluorine-process plasmas the dominant ion mass is set by the feed; take an
effective ion mass near argon-to-$\mathrm{CF}_3^+$ scale, $m_i \approx 40\,
\text{u} = 6.6\times 10^{-26}\,\si{\kilogram}$, so $u_B = \sqrt{(4.8\times
10^{-19})/(6.6\times 10^{-26})} \approx 2.7\times 10^{3}\,\si{\meter\per
\second}$. The ion flux to the wafer is the Bohm flux evaluated at the
sheath edge, where the density has fallen to about $0.6\,n_e$:
\[
\Gamma_i = 0.6\,n_e u_B = 0.6(2\times 10^{16})(2.7\times 10^{3})
\approx 3.2\times 10^{19}\,\si{\per\meter\squared\per\second}.
\]
Each ion falls through the sheath and arrives with an energy set by the
time-averaged sheath voltage, $E_i \approx e\overline{V}_s \approx
350\,\si{\electronvolt}$, in the range that drives anisotropic etching. The
ion-current density at the wafer is $J_i = e\Gamma_i = (1.6\times
10^{-19})(3.2\times 10^{19}) \approx 5.1\,\si{\ampere\per\meter\squared}$,
about $0.5\,\si{\milli\ampere\per\centi\meter\squared}$, and the ion power
delivered to the surface is $J_i\overline{V}_s \approx 5.1\times 350 \approx
1.8\times 10^{3}\,\si{\watt\per\meter\squared}$, a heat load the wafer chuck
must remove to hold temperature.

\subsection{Etch rate, selectivity, and anisotropy}

The vertical etch rate follows from the flux of removed atoms. In an
ion-assisted etch the removal rate is the ion flux times a yield $Y$, the
number of substrate atoms removed per incident ion, which for
silicon under a few-hundred-electronvolt fluorine-ion bombardment is of
order $Y \approx 1$. The areal removal rate is then $Y\Gamma_i \approx 3.2
\times 10^{19}\,\si{\per\meter\squared\per\second}$. Silicon has an atomic
number density of $n_{\mathrm{Si}} = 5.0\times 10^{28}\,\si{\per\meter
\cubed}$, so the vertical etch rate is
\[
\dot{z} = \frac{Y\Gamma_i}{n_{\mathrm{Si}}}
= \frac{3.2\times 10^{19}}{5.0\times 10^{28}}
\approx 6.4\times 10^{-10}\,\si{\meter\per\second},
\]
about $0.64\,\si{\nano\meter\per\second}$, or $38\,\si{\nano\meter}$ per
minute. A $300\,\si{\nano\meter}$ trench therefore takes about eight
minutes, a realistic figure for this class of recipe.

Selectivity is the ratio of the film etch rate to the mask etch rate. If
the photoresist mask etches at $0.1\,\si{\nano\meter\per\second}$ under the
same bombardment, the selectivity is $0.64/0.1 \approx 6{:}1$, enough to
clear a $300\,\si{\nano\meter}$ film through a $200\,\si{\nano\meter}$ mask
with margin. Raising the bias to speed the etch raises the mask erosion
faster, so the selectivity falls as the bias climbs, the trade drawn in
Figure~\ref{fig:vol04ch13:etch-bias}.

Anisotropy is the degree to which the etch runs vertically rather than
sideways. Define $A = 1 - \dot{z}_{\mathrm{lat}}/\dot{z}_{\mathrm{vert}}$,
so $A = 1$ is a perfectly vertical wall and $A = 0$ is a fully isotropic
etch. The lateral rate is the chemical (unbiased) component, since only the
ions are directional; if the spontaneous chemical etch contributes
$0.05\,\si{\nano\meter\per\second}$ of lateral removal against the
$0.64\,\si{\nano\meter\per\second}$ vertical rate, then $A = 1 - 0.05/0.64
\approx 0.92$, a steep but not perfectly vertical wall. Deepening the sheath
raises the directional ion rate without raising the isotropic chemical rate,
which pushes $A$ toward one, the trend the figure shows.

% Figure: etch behaviour versus RF bias power in a CCP reactor. The directed
% ion energy rises with applied bias (the self-bias voltage grows), which
% raises the vertical etch rate and the anisotropy but lowers the
% mask-to-film selectivity because the harder ion bombardment also attacks
% the mask. The curves are schematic; the numbers track the CCP case study.
\begin{figure}[htbp]
\centering
\begin{tikzpicture}
\begin{axis}[
  width=0.82\textwidth, height=0.52\textwidth,
  axis y line*=left,
  xlabel={RF bias power $P_{\mathrm{bias}}$ (W)},
  ylabel={etch rate (nm\,min$^{-1}$) and anisotropy ($\times100$)},
  xmin=0, xmax=400, ymin=0, ymax=320,
  grid=both, grid style={enggray!20},
  every axis plot/.append style={thick},
  legend pos=north west, legend cell align=left,
]
  % etch rate: rises and saturates as ion-energy-driven removal dominates
  \addplot[engblue, domain=0:400, samples=80]
    {300*(1-exp(-x/130))};
  \addlegendentry{vertical etch rate}
  % anisotropy (fraction times 100): rises toward 1 with directed ions
  \addplot[engorange, domain=0:400, samples=80]
    {100*(1-exp(-x/70))};
  \addlegendentry{anisotropy $A\times100$}
\end{axis}
\begin{axis}[
  width=0.82\textwidth, height=0.52\textwidth,
  axis y line*=right,
  axis x line=none,
  ylabel={selectivity (film\,:\,mask)},
  xmin=0, xmax=400, ymin=0, ymax=12,
  every axis plot/.append style={thick},
  legend pos=north east, legend cell align=left,
]
  % selectivity: falls as harder ion bombardment erodes the mask
  \addplot[engred, domain=0:400, samples=80, dashed]
    {2+9*exp(-x/120)};
  \addlegendentry{selectivity}
\end{axis}
\end{tikzpicture}
\caption[Etch behaviour versus RF bias power]{Etch behaviour in a
capacitively coupled reactor as the radio-frequency bias power on the wafer
electrode is raised. Higher bias deepens the self-bias voltage and the
sheath, so ions arrive with more directed energy. The vertical etch rate
(left axis, solid blue) climbs and saturates once ion-energy-driven removal
dominates the chemical floor, and the anisotropy (left axis, solid orange,
shown as a fraction times one hundred) approaches one as the directed ion
flux overwhelms the isotropic chemical component. The mask-to-film
selectivity (right axis, dashed red) falls in the same direction, because
the harder bombardment that straightens the side walls also erodes the
mask. The working point of an etch recipe sits where the anisotropy is
adequate and the selectivity has not yet collapsed. The curves are
schematic; the operating numbers are carried in the CCP case study.}
\label{fig:vol04ch13:etch-bias}
\end{figure}


\subsection{The loading effect and the verdict}

A real wafer does not etch at the single-feature rate, because the supply of
reactive radicals is finite and shared across all the open area. The
\engterm{loading effect} ties the etch rate to the total exposed film area:
when more of the wafer is open, each feature draws from the same radical
budget and every feature etches slower. A simple balance writes the etch
rate as
\[
\dot{z} = \frac{\dot{z}_0}{1 + k A_{\mathrm{open}}},
\]
where $\dot{z}_0$ is the rate at vanishing open area, $A_{\mathrm{open}}$ is
the open fraction, and $k$ measures how strongly the radical supply is
depleted. For a wafer with $30\%$ open area and a loading constant $k =
1.5$, the etch rate falls to $\dot{z}_0/(1 + 1.5\times 0.30) = \dot{z}_0/
1.45 \approx 0.69\,\dot{z}_0$, a $31\%$ slowdown relative to a nearly bare
wafer. A process tuned on a sparse test pattern will therefore over-etch a
production wafer's sparse regions while still clearing its dense regions, so
the etch time must be set for the loaded rate and the uniformity checked
across the real layout.

The verdict on the recipe is workable with the same caution the tokamak
case carried. The reactor delivers a $\sim 38\,\si{\nano\meter}$-per-minute
silicon etch at $6{:}1$ selectivity and $0.92$ anisotropy, adequate for a
$300\,\si{\nano\meter}$ feature, but the selectivity and the loading margin
both sit close enough to their limits that a $30\%$ open-area swing or a
bias increase to speed the etch would erode the mask or starve the dense
regions. A process engineer responds by holding the bias at the lowest
value that keeps the anisotropy above specification, by adding a
polymer-forming gas to passivate the side walls and protect the mask, and
by setting the etch time from the loaded rate measured on a production-
representative pattern rather than on a bare-wafer coupon. The case
reproduces, at room scale, the same discipline the tokamak case demanded:
the operating point must clear every requirement with enough margin to
survive the scatter in the real process, and the margin, not the peak etch
rate, sets the recipe.

\section{A third case study: a Hall-effect ion thruster}\pathtag{standard}

The tokamak case sat at the hot, dense corner of the plasma map and the etch
reactor at the cool, dense corner. Electric propulsion occupies a third
position: a cool, tenuous, partially magnetised plasma whose entire purpose
is to throw ions overboard at high speed. We design a Hall-effect thruster
end to end, from the discharge and ionisation, through the beam voltage and
current, to the exhaust velocity, the specific impulse, the thrust, and the
efficiency, and we close by spending the thruster against a real mission's
propellant budget. The crossed-field physics of Section 13.2 and the
ionisation and sheath physics of the earlier sections all reappear, now
serving propulsion rather than confinement or etching. The geometry is drawn
in Figure~\ref{fig:vol04ch13:hall-thruster} and the voltage trade in
Figure~\ref{fig:vol04ch13:hall-performance}.

% Figure: Hall-effect thruster cross-section. An annular discharge channel
% carries an axial electric field and a radial magnetic field; the crossed
% fields trap electrons in an azimuthal Hall current while xenon ions are
% accelerated out of the channel as the beam. The end-to-end case study
% carries the numbers.
\begin{figure}[htbp]
\centering
\begin{tikzpicture}[scale=1.0]
  % discharge channel walls (annular, shown in cross-section: two boxes)
  \fill[enggray!15] (0,0.7) rectangle (3.2,1.5);
  \fill[enggray!15] (0,-1.5) rectangle (3.2,-0.7);
  \draw[enggray, thick] (0,0.7) rectangle (3.2,1.5);
  \draw[enggray, thick] (0,-1.5) rectangle (3.2,-0.7);
  \node[enggray, font=\scriptsize, anchor=west] at (0.1,1.1) {dielectric wall};
  % anode at the back of the channel
  \draw[engred, line width=1.2pt] (0.15,-0.65) -- (0.15,0.65);
  \node[engred, font=\scriptsize, anchor=east, rotate=90] at (0.05,0) {anode (+)};
  % magnetic coils / poles producing radial B
  \node[engblue, font=\scriptsize, anchor=south] at (2.5,1.55) {N};
  \node[engblue, font=\scriptsize, anchor=north] at (2.5,-1.55) {S};
  % radial B field arrows (across the channel)
  \foreach \x in {1.4,2.0,2.6} {
    \draw[engblue, ->, thick] (\x,0.65) -- (\x,-0.6);
  }
  \node[engblue, font=\scriptsize, anchor=south west] at (2.6,0.0) {$\vect{B}$ (radial)};
  % axial E field arrow inside the channel
  \draw[engorange, ->, line width=1pt] (0.4,0) -- (3.0,0);
  \node[engorange, font=\scriptsize, anchor=south] at (1.6,0.05) {$\vect{E}$ (axial)};
  % azimuthal Hall electron current (into/out of page markers)
  \node[englime, font=\scriptsize] at (1.9,0.35) {$\otimes$};
  \node[englime, font=\scriptsize] at (1.9,-0.35) {$\odot$};
  \node[englime, font=\scriptsize, anchor=west] at (3.25,0.4)
    {Hall electron current};
  % accelerated ion beam leaving the channel
  \foreach \y in {-0.35,0,0.35} {
    \draw[engred, ->, thick] (3.0,\y) -- (4.3,\y);
  }
  \node[engred, font=\scriptsize, anchor=west] at (4.0,-0.7) {ion beam};
  % cathode neutraliser
  \fill[engblack] (4.2,1.1) circle (0.08);
  \node[engblack, font=\scriptsize, anchor=west] at (4.3,1.1)
    {cathode (neutraliser)};
  \draw[engblack, ->, dashed] (4.2,1.0) .. controls (4.0,0.3) .. (3.8,0.05);
\end{tikzpicture}
\caption[Hall-effect thruster]{Cross-section of one side of an annular
Hall-effect thruster. An anode at the back of the dielectric channel sets up
an axial electric field $\vect{E}$, while magnet coils drive a radial
magnetic field $\vect{B}$ across the channel. The crossed fields trap the
light electrons in an azimuthal $E\times B$ drift, the Hall current, which
greatly raises their residence time and lets them ionise the injected xenon
efficiently. The heavy ions have a Larmor radius larger than the channel and
are not magnetised, so they fall freely down the axial field and leave as a
high-velocity beam. An external cathode supplies electrons both to sustain
the discharge and to neutralise the departing beam, without which the
spacecraft would charge negative and pull the ions back. The end-to-end case
study carries the discharge, beam, thrust, and efficiency numbers.}
\label{fig:vol04ch13:hall-thruster}
\end{figure}


\subsection{The thruster and its operating point}

Take a Hall thruster with an annular discharge channel of mean diameter
$100\,\si{\milli\meter}$ and channel width $15\,\si{\milli\meter}$, fed with
xenon and drawing a discharge power of $P_{d} = 4.5\,\si{\kilo\watt}$ at a
discharge voltage $V_{d} = 300\,\si{\volt}$. The radial magnetic field in
the channel is $B \approx 20\,\si{\milli\tesla}$, chosen by the design
principle that orders the whole device: the field must magnetise the
electrons while leaving the ions free. An electron at the few-electronvolt
discharge temperature has a Larmor radius $r_{L,e} = m_{e}v_{\perp}/(eB)
\approx (9.11\times 10^{-31})(10^{6})/[(1.6\times 10^{-19})(0.02)] \approx
3\times 10^{-4}\,\si{\meter}$, about $0.3\,\si{\milli\meter}$, far smaller
than the $15\,\si{\milli\meter}$ channel, so the electrons are trapped. A
xenon ion accelerated to the beam energy has a Larmor radius of metres, far
larger than the channel, so it is unmagnetised and falls freely down the
axial field. The same field that holds the electrons lets the ions go, which
is the operating principle of the device.

The trapped electrons cannot cross the radial field to reach the anode
directly; they drift azimuthally in the crossed axial $\vect{E}$ and radial
$\vect{B}$ fields at the $E\times B$ velocity, $v_{E\times B} = E/B$. With
the $300\,\si{\volt}$ falling across an acceleration length of $\sim 20\,\si{
\milli\meter}$, the axial field is $E \approx 300/0.02 = 1.5\times
10^{4}\,\si{\volt\per\meter}$, so the azimuthal Hall drift is $v_{E\times B}
= 1.5\times 10^{4}/0.02 = 7.5\times 10^{5}\,\si{\meter\per\second}$. This
closed Hall current is the engine of ionisation: the long azimuthal path
gives each electron many chances to ionise an incoming xenon atom before it
finally crosses to the anode by collisions.

\subsection{Beam current, exhaust velocity, and specific impulse}

A xenon ion ($m_{i} = 131\,\text{u} = 2.18\times 10^{-25}\,\si{\kilogram}$)
falling through the discharge voltage gains kinetic energy $eV_{d}$, so its
exhaust velocity is
\[
v_{\mathrm{ex}} = \sqrt{\frac{2eV_{d}}{m_{i}}}
= \sqrt{\frac{2(1.6\times 10^{-19})(300)}{2.18\times 10^{-25}}}
\approx 2.1\times 10^{4}\,\si{\meter\per\second},
\]
about $21\,\si{\kilo\meter\per\second}$. The specific impulse referenced to
this ideal exhaust is $I_{sp} = v_{\mathrm{ex}}/g_{0} = 2.1\times 10^{4}/9.81
\approx 2150\,\si{\second}$. Real thrusters fall short of the ideal because
not every ion is born at the anode potential and some leave at an angle, so
a voltage utilisation $\eta_{v}\approx 0.95$ and a beam divergence
correction $\cos\theta$ with $\theta\approx 15^{\circ}$ ($\cos\theta\approx
0.97$) reduce the effective exhaust to $v_{\mathrm{eff}}\approx 0.92\,
v_{\mathrm{ex}}\approx 1.9\times 10^{4}\,\si{\meter\per\second}$, an
effective $I_{sp}\approx 1980\,\si{\second}$.

The beam current is set by how much of the discharge current is carried by
ions out the front rather than by electrons to the anode. Write the beam
current as $I_{b} = \eta_{b}I_{d}$ with a current utilisation $\eta_{b}
\approx 0.75$, a typical fraction. The discharge current is $I_{d} =
P_{d}/V_{d} = 4500/300 = 15\,\si{\ampere}$, so the ion beam carries $I_{b} =
0.75\times 15 = 11.25\,\si{\ampere}$. The ion mass flow follows from the
beam current and the ion charge,
\[
\dot{m}_{i} = \frac{I_{b}m_{i}}{e}
= \frac{(11.25)(2.18\times 10^{-25})}{1.6\times 10^{-19}}
\approx 1.5\times 10^{-5}\,\si{\kilogram\per\second},
\]
about $15\,\si{\milli\gram\per\second}$ of ionised xenon leaving as beam.

\subsection{Thrust, mass utilisation, and efficiency}

The thrust is the momentum flux of the beam, $T = \dot{m}_{i}v_{\mathrm{eff}}
+ \dot{m}_{n}v_{n}$, where the second term is the small contribution of
un-ionised neutrals leaking out at thermal speed. Neglecting the neutral
thrust, the ion thrust is
\[
T \approx \dot{m}_{i}v_{\mathrm{eff}}
= (1.5\times 10^{-5})(1.9\times 10^{4})
\approx 0.29\,\si{\newton},
\]
about $290\,\si{\milli\newton}$, the order quoted for a multi-kilowatt Hall
thruster. The mass utilisation efficiency measures how much of the injected
xenon is actually ionised and accelerated rather than escaping cold. If the
thruster is fed a total $\dot{m}_{\mathrm{tot}} = 18\,\si{\milli\gram\per
\second}$, the mass utilisation is $\eta_{m} = \dot{m}_{i}/\dot{m}_{\mathrm{
tot}} = 15/18 \approx 0.83$. The neutral leak that limits $\eta_{m}$ is the
same un-ionised fraction the etch case study saw at low ionisation; here a
strong Hall current is what pushes the ionised fraction high enough to make
the thruster efficient.

The anode efficiency collects the chain into one figure of merit, the ratio
of beam kinetic power to discharge power:
\[
\eta_{a} = \frac{\tfrac12\dot{m}_{i}v_{\mathrm{eff}}^{2}}{P_{d}}
= \frac{\tfrac12(1.5\times 10^{-5})(1.9\times 10^{4})^{2}}{4500}
\approx \frac{2.7\times 10^{3}}{4500} \approx 0.60.
\]
A $60\%$ anode efficiency is in the right range for this class of thruster.
The efficiency factors cleanly into the pieces computed above, $\eta_{a}
\approx \eta_{b}\,\eta_{v}\,\eta_{m}\cos^{2}\theta \approx (0.75)(0.95)(0.83)
(0.94) \approx 0.56$, close to the direct estimate; the small gap is the
neutral and double-ion bookkeeping a detailed model would carry. The voltage
trade in Figure~\ref{fig:vol04ch13:hall-performance} shows why $300\,\si{
\volt}$ is a sensible operating point: it sits near the knee where the
efficiency has nearly saturated while the thrust-to-power has not yet
collapsed.

% Figure: Hall-thruster performance trade. Thrust efficiency and thrust-to-
% power both vary with discharge (beam) voltage at fixed power: higher voltage
% raises exhaust velocity and specific impulse but lowers the thrust produced
% per unit power. The case study reads its operating point off this trade.
\begin{figure}[htbp]
\centering
\begin{tikzpicture}
\begin{axis}[
  width=0.82\textwidth, height=0.52\textwidth,
  xlabel={discharge voltage $V_d$ (V)},
  ylabel={anode efficiency $\eta_a$},
  xmin=100, xmax=600, ymin=0.3, ymax=0.75,
  axis y line*=left,
  grid=both, grid style={enggray!20},
  every axis plot/.append style={thick},
]
  % efficiency rises and saturates with voltage
  \addplot[engblue, domain=100:600, samples=80]
    {0.72 - 18/(x-40)};
  \label{plot:halleff}
\end{axis}
\begin{axis}[
  width=0.82\textwidth, height=0.52\textwidth,
  xmin=100, xmax=600,
  ymin=40, ymax=90,
  axis y line*=right,
  axis x line=none,
  ylabel={thrust-to-power (mN/kW)},
  every axis plot/.append style={thick},
]
  % thrust-to-power falls as voltage rises (T/P ~ 1/sqrt(Vd) at fixed eta)
  \addplot[engorange, domain=100:600, samples=80]
    {850/sqrt(x)};
  \addlegendimage{engblue}\addlegendentry{anode efficiency}
  \addlegendentry{thrust-to-power}
  \draw[dashed, enggray] (axis cs:300,40) -- (axis cs:300,90);
  \node[enggray, font=\scriptsize, anchor=south, rotate=90]
    at (axis cs:300,55) {case operating point};
\end{axis}
\end{tikzpicture}
\caption[Hall-thruster voltage trade]{The discharge-voltage trade in a
Hall-effect thruster at fixed input power. Raising the discharge voltage
raises the exhaust velocity and the specific impulse, and the anode
efficiency (blue, left axis) climbs and saturates as a larger fraction of
the supplied energy ends up as directed beam kinetic energy rather than as
ionisation cost and wall losses. The thrust produced per unit power (orange,
right axis) falls with voltage, because at fixed power a faster exhaust
carries the same momentum flux in fewer, faster ions. A mission that is
propellant-limited chooses high voltage for high specific impulse; a mission
that is power-limited or thrust-limited accepts lower voltage for more thrust
per kilowatt. The case study reads its $300\,\si{\volt}$ operating point off
the knee of this trade. The curves are schematic, drawn to show the shape of
the trade rather than the performance of a specific thruster.}
\label{fig:vol04ch13:hall-performance}
\end{figure}


\subsection{The mission budget and the verdict}

A thruster earns its place by the propellant it saves over a mission, so we
spend it against a delta-v. Take a $2000\,\si{\kilo\gram}$ spacecraft that
must raise its orbit with a $\Delta v = 2000\,\si{\meter\per\second}$
manoeuvre. The Tsiolkovsky equation gives the propellant mass
\[
m_{p} = m_{0}\left(1 - e^{-\Delta v/v_{\mathrm{eff}}}\right)
= 2000\left(1 - e^{-2000/19000}\right)
\approx 2000(1 - 0.900) \approx 200\,\si{\kilogram}
\]
of xenon. A chemical bipropellant at $v_{\mathrm{eff}} = I_{sp}g_{0} =
300\times 9.81 \approx 2940\,\si{\meter\per\second}$ would need $m_{p} =
2000(1 - e^{-2000/2940}) \approx 2000(1 - 0.506) \approx 990\,\si{\kilogram}$
for the same manoeuvre, nearly five times as much. The saving of $\sim
790\,\si{\kilogram}$ is the entire case for electric propulsion, and it pays
the spacecraft designer in payload mass or launch cost.

The cost is time. At $T = 0.29\,\si{\newton}$ the acceleration of the loaded
spacecraft is $a = T/m_{0} = 0.29/2000 \approx 1.5\times 10^{-4}\,\si{\meter
\per\second\squared}$, so accumulating $2000\,\si{\meter\per\second}$ takes
$t = \Delta v/a \approx 2000/1.5\times 10^{-4} \approx 1.3\times 10^{7}\,\si{
\second}$, about $150\,\si{\day}$ of nearly continuous thrust. The propellant
throughput over that time is $\dot{m}_{\mathrm{tot}}\times t = (1.8\times
10^{-5})(1.3\times 10^{7}) \approx 234\,\si{\kilogram}$, consistent with the
$200\,\si{\kilogram}$ of beam xenon plus the neutral leak, a useful internal
check that the thrust, flow, and mission numbers close on each other.

The verdict carries the same discipline as the other two cases. The thruster
delivers the manoeuvre on roughly one fifth the propellant of a chemical
stage, at the price of a months-long burn and a $4.5\,\si{\kilo\watt}$ power
demand the spacecraft's solar array must supply for the whole transfer. The
binding constraints are not the plasma physics, which is comfortable, but the
power available and the thruster's erosion lifetime, since the same energetic
ions that carry the beam also sputter the channel walls over thousands of
hours. The design choice that sets the mission is the discharge voltage: a
higher voltage raises $I_{sp}$ and cuts propellant further but lowers the
thrust per kilowatt and stretches the burn, so the operating point is read
off the trade in Figure~\ref{fig:vol04ch13:hall-performance} against the
mission's tolerance for transfer time, exactly as the tokamak's operating
point was read off its triple-product margin and the etch reactor's off its
selectivity and loading margin.

\begin{mastery}
The current-utilisation factor $\eta_{b}\approx 0.75$ the case took as given
hides one of the longest-running puzzles in plasma engineering. Classical
collision theory predicts that magnetised electrons cross the radial field
to the anode only by collisions, at a rate that would make the electron
current a small fraction of what is observed. The measured cross-field
electron transport in a Hall thruster is one to two orders of magnitude
larger than classical, the same \engterm{anomalous transport} first
quantified by Bohm in mid-twentieth-century magnetised discharges and written
as the \engterm{Bohm diffusion} coefficient $D_{\mathrm{Bohm}} = k_{B}T_{e}/
(16eB)$. Bohm diffusion scales as $1/B$ rather than the classical $1/B^{2}$,
so it dominates at high field exactly where classical theory predicts the
best confinement. The mechanism is turbulent: small fluctuations in the
plasma density and potential set up fluctuating $E\times B$ drifts that carry
electrons across the field far faster than collisions alone. In a Hall
thruster the anomalous transport sets the electron current to the anode, the
discharge current, and therefore the efficiency, so a predictive thruster
model cannot be built from classical transport and must either measure the
anomalous coefficient or resolve the turbulence directly. The same anomalous
cross-field transport limits the confinement time of every magnetic fusion
device, which is why the energy confinement time $\tau_{E}$ that drives the
Lawson balance is taken from empirical confinement scaling laws rather than
computed from first principles: the turbulent transport that sets it is not
yet calculable from the equilibrium alone.
\end{mastery}

\section{An applied vignette: the plasma arc cutter}\pathtag{standard}

The three case studies sat at the extremes of the density-temperature map.
A shorter example brings the same physics down to a tool a fabrication shop
keeps in the corner: the plasma arc cutter that slices steel plate. It is a
thermal-equilibrium arc, the Saha regime of the foundations section rather
than the non-equilibrium discharge of the etch reactor, and it shows the
chapter's formulas working at the scale of a hand tool.

A plasma cutter strikes a constricted electric arc between a tungsten
electrode and the workpiece, then forces a gas, usually compressed air or
nitrogen, through a narrow copper nozzle that squeezes the arc to a thin,
fast jet. The constriction raises the current density and the temperature in
the channel to $T \approx 1.5\times 10^{4}$ to $3\times 10^{4}\,\si{
\kelvin}$, hot enough that the Saha equation, not an electron-impact
discharge model, fixes the ionisation. For the air arc near $2\times
10^{4}\,\si{\kelvin}$, the ionisation fraction computed exactly as in the
arc-plasma worked example is a few percent to tens of percent, enough to
carry the hundreds of amperes the torch draws and to deliver the jet of hot
ionised gas that melts and blows away the metal.

Take a hand torch running $I = 60\,\si{\ampere}$ at an arc voltage $V =
120\,\si{\volt}$, so the electrical power into the arc is $P = IV = 7.2\,\si{
\kilo\watt}$. The nozzle orifice is $d = 1.2\,\si{\milli\meter}$, so the
current density in the channel is $J = I/(\pi d^{2}/4) = 60/[\pi(1.2\times
10^{-3})^{2}/4] \approx 5.3\times 10^{7}\,\si{\ampere\per\meter\squared}$, a
density that a solid copper conductor could not sustain without melting but
that the gas plasma carries because its ionised channel is continually
replenished. The arc plasma at this temperature has a Spitzer-like
resistivity of order $\eta \approx 10^{-4}\,\si{\ohm\meter}$ (much higher
than a fusion plasma because the arc is two orders of magnitude cooler), so
the ohmic dissipation per unit volume is $p_{\Omega} = \eta J^{2} \approx
(10^{-4})(5.3\times 10^{7})^{2} \approx 2.8\times 10^{11}\,\si{\watt\per
\meter\cubed}$, an enormous volumetric heating that holds the channel at its
operating temperature against the convective and radiative losses to the
cold gas sheath around it. The same $T_{e}^{-3/2}$ resistivity scaling that
made a fusion plasma a near-perfect conductor here keeps the cooler arc
dissipative enough to stay hot, the self-regulating balance that lets an arc
sit at a stable temperature.

The cut itself is a heat-transfer problem fed by the arc. The jet must
deliver enough power to melt a kerf through the plate and blow the melt
clear. To cut mild steel of thickness $h = 10\,\si{\milli\meter}$ at a
travel speed $u = 0.5\,\si{\meter\per\minute} = 8.3\times 10^{-3}\,\si{\meter
\per\second}$ with a kerf width $w = 1.5\,\si{\milli\meter}$, the volume
removal rate is $\dot{V} = h w u = (0.01)(1.5\times 10^{-3})(8.3\times
10^{-3}) \approx 1.2\times 10^{-7}\,\si{\meter\cubed\per\second}$. Steel of
density $7800\,\si{\kilogram\per\meter\cubed}$ removed at this rate is a mass
flow $\dot{m} = 9.7\times 10^{-4}\,\si{\kilogram\per\second}$. Raising that
steel from room temperature through its melt point and supplying the latent
heat of fusion costs roughly $c_{p}\Delta T + L_{f} \approx (500)(1500) +
2.7\times 10^{5} \approx 1.0\times 10^{6}\,\si{\joule\per\kilogram}$, so the
power that must reach the kerf is $\dot{m}\times 10^{6} \approx 9.7\times
10^{2}\,\si{\watt}$, about $1\,\si{\kilo\watt}$. The torch delivers
$7.2\,\si{\kilo\watt}$, so the coupling efficiency from electrical input to
useful melting is of order $15\%$, the rest carried off as radiation, as
heat conducted into the plate, and as enthalpy of the spent gas jet. A
shop's experience that a $60\,\si{\ampere}$ torch cuts roughly
$10\,\si{\milli\meter}$ plate at a modest speed is the same balance read in
the other direction.

The engineering constraints follow the same pattern as the case studies. The
nozzle erodes because the constricted arc concentrates the heat on the
copper orifice, so consumable life is set by the arc current and the
duty cycle, and pushing the current to cut thicker plate trades cut capacity
against consumable cost, the same energy-against-erosion trade the Hall
thruster faced at its channel walls. The cut quality, square edge versus
bevelled, is set by how tightly the nozzle constricts the arc, the same
directionality-versus-spread trade the etch sheath set for ion bombardment.
The plasma cutter is the chapter's physics at the scale of a tool a single
operator carries, and it obeys the same laws as the reactor that fills a
building.

\section{A second applied vignette: the magnetron sputtering source}\pathtag{standard}

The etch reactor removed material and the PECVD tool added it from gas-phase
chemistry. A third deposition tool adds material from a solid source, and it
turns the crossed-field trapping of Section 13.2 into a manufacturing advantage.
The magnetron sputtering source coats everything from architectural glass to
hard-disk platters to the reflective layer of a snack-food bag, and it works by
knocking atoms off a solid target with ion bombardment and letting them settle
on a substrate. Its physics is the same $E\times B$ electron trap the Hall
thruster used, redeployed to build a dense, low-pressure plasma against a target
rather than to expel a beam.

A planar magnetron places permanent magnets behind a flat target that is biased
several hundred volts negative, Figure~\ref{fig:vol04ch13:magnetron-sputter}. A
centre pole and a surrounding ring pole of opposite sign arch the magnetic field
over the target face, so that above the target the field runs nearly parallel to
the surface while the electric field of the sheath points into it. Electrons
sputtered or ionised near the target are trapped in this crossed-field region:
they gyrate on the field, are repelled from the target by the sheath, and drift
azimuthally at the $E\times B$ velocity in a closed loop, the race track. The
long confined path gives each electron many ionising collisions before it
finally escapes across the field to the anode, so the discharge sustains a dense
plasma at a pressure ten to a hundred times lower than an unmagnetised diode
would need. The low pressure is the entire point: sputtered atoms travel to the
substrate with few gas collisions, so they arrive with their sputter energy and
build a dense, adherent film, and the trapping is what makes that low-pressure
operation possible.

% Figure: planar magnetron sputtering source. Permanent magnets behind the
% target set up an arched magnetic field over the target face; the crossed
% radial E and the arched B trap electrons in a closed E-cross-B race track,
% intensifying ionisation just above the erosion groove. Ions are pulled to the
% target, sputter atoms off it, and the atoms deposit on the substrate above.
\begin{figure}[htbp]
\centering
\begin{tikzpicture}[>=Latex, scale=1.0]
  % target (cathode) slab
  \fill[enggray!25] (-3.2,0) rectangle (3.2,0.5);
  \node[enggray, font=\scriptsize, anchor=north east] at (3.2,0) {target (cathode)};
  % magnet bar cross-sections behind the target: outer poles S, centre pole N
  \fill[engblue!30] (-3.0,-0.7) rectangle (-2.2,0.0);
  \fill[engred!30]  (-0.4,-0.7) rectangle (0.4,0.0);
  \fill[engblue!30] (2.2,-0.7) rectangle (3.0,0.0);
  \node[font=\scriptsize] at (-2.6,-0.35) {S};
  \node[font=\scriptsize] at (0.0,-0.35) {N};
  \node[font=\scriptsize] at (2.6,-0.35) {S};
  % arched field lines from centre N to each outer S, above the target face
  \draw[engred, thick, ->] (-0.2,0.5) .. controls (-1.0,1.4) and (-1.8,1.4) .. (-2.6,0.5);
  \draw[engred, thick, ->] (0.2,0.5) .. controls (1.0,1.4) and (1.8,1.4) .. (2.6,0.5);
  % electron race-track marker: two crosses where field is most horizontal
  \node[engorange, font=\footnotesize] at (-1.4,1.15) {$\otimes$};
  \node[engorange, font=\footnotesize] at (1.4,1.15) {$\otimes$};
  \node[engorange, font=\scriptsize, anchor=south] at (0.0,1.5)
    {trapped-electron race track};
  % erosion grooves under the arches
  \draw[engblack, thick] (-1.7,0.5) .. controls (-1.4,0.28) and (-1.1,0.28) .. (-0.8,0.5);
  \draw[engblack, thick] (0.8,0.5) .. controls (1.1,0.28) and (1.4,0.28) .. (1.7,0.5);
  \node[engblack, font=\scriptsize, anchor=north] at (0.0,0.15) {erosion grooves};
  % sputtered atoms going up
  \draw[engblue, ->] (-1.25,0.6) -- (-1.0,2.4);
  \draw[engblue, ->] (1.25,0.6) -- (1.0,2.4);
  \draw[engblue, ->] (0.0,0.6) -- (0.0,2.4);
  % substrate at top
  \fill[englime!35] (-2.6,2.6) rectangle (2.6,3.0);
  \node[font=\scriptsize, anchor=south] at (0.0,3.0) {substrate (deposited film)};
  % incoming Ar+ ion pulled down to target
  \draw[engorange, ->] (2.0,2.2) -- (1.6,0.6);
  \node[engorange, font=\scriptsize, anchor=west] at (2.0,2.2) {Ar$^+$};
\end{tikzpicture}
\caption[Planar magnetron sputtering source]{A planar magnetron sputtering
source. Permanent magnets behind the target set up an arched magnetic field
(red) over the target face; the centre pole and the outer ring poles of
opposite sign make the field lines bow out and return, so above the erosion
groove the field runs nearly parallel to the target surface. Electrons emitted
from the negatively biased target are trapped there, gyrating on the field and
drifting azimuthally in the crossed radial electric and arched magnetic fields,
so they run in a closed race track (orange) that concentrates ionisation just
above the target. The dense plasma pulls argon ions into the target, where they
sputter atoms loose; the sputtered atoms travel up and deposit on the substrate
as a thin film. The magnetic trapping is what lets a magnetron run a dense
discharge at low pressure, which is why sputtered films are denser and cleaner
than those from an unmagnetised diode. The erosion grooves under the arches are
where the target wears fastest and set the target utilisation.}
\label{fig:vol04ch13:magnetron-sputter}
\end{figure}


Take a source with a $150\,\si{\milli\meter}$ circular target driven at
$P = 500\,\si{\watt}$ and a target voltage $V_t = 400\,\si{\volt}$, sputtering
aluminium in argon at $0.5\,\si{\pascal}$. The discharge current is $I = P/V_t =
500/400 = 1.25\,\si{\ampere}$, and if a fraction $\gamma \approx 0.7$ of that
current is ions arriving at the target, the ion flux is $\Gamma_i = \gamma I/(e
A_t)$ with $A_t = \pi(0.075)^{2} \approx 1.77\times 10^{-2}\,\si{\meter\squared}$,
so $\Gamma_i = (0.7)(1.25)/[(1.6\times 10^{-19})(1.77\times 10^{-2})] \approx
3.1\times 10^{20}\,\si{\per\meter\squared\per\second}$. Each $400\,\si{
\electronvolt}$ argon ion sputters aluminium with a yield $Y \approx 1.2$ atoms
per ion, so the sputtered-atom flux off the target is $Y\Gamma_i \approx
3.7\times 10^{20}\,\si{\per\meter\squared\per\second}$. Spread over a substrate
the same area at unit transport efficiency, and with aluminium's atomic density
$n_{\mathrm{Al}} = 6.0\times 10^{28}\,\si{\per\meter\cubed}$, the deposition rate
is $\dot{z} = Y\Gamma_i/n_{\mathrm{Al}} \approx 3.7\times 10^{20}/6.0\times
10^{28} \approx 6.2\times 10^{-9}\,\si{\meter\per\second}$, about $6\,\si{\nano
\meter}$ per second, or a few hundred nanometres a minute, the rate a coater is
built around. Raising the power raises the current and the rate in proportion,
until target heating or arcing sets the ceiling.

The engineering trades follow the same pattern the other tools set. The race
track wears a groove in the target because ionisation concentrates where the
field is most parallel, so target utilisation is only $\sim 30\%$ before the
groove cuts through and the target must be replaced, the same energy-against-
erosion trade the Hall thruster met at its channel walls. Reactive sputtering,
admitting oxygen or nitrogen to grow an oxide or nitride film, drives the source
into hysteresis: the reactive gas poisons the target surface, dropping the
sputter yield and the rate, so the process must be held on the steep part of the
hysteresis loop by feedback on the gas flow, the same margin-against-instability
discipline the tokamak's beta limit demanded. The magnetron is the chapter's
crossed-field physics earning its keep in a coating plant, and it reads by the
same operational formulas as the confinement and propulsion devices that look
nothing like it.

\section{A third applied vignette: the atmospheric-pressure plasma jet}\pathtag{standard}

Every plasma so far ran in a vacuum vessel. A growing family of tools runs in
open air at atmospheric pressure and near room temperature, and it turns the
non-equilibrium between hot electrons and cold heavy particles into a way to do
chemistry on surfaces and on living tissue without cooking either. The
atmospheric-pressure plasma jet is a thin tube through which a noble gas flows
past a powered electrode, producing a visible plume of cold plasma that a hand
can touch. Its physics is the same electron-heavy-particle energy separation the
fluorescent lamp uses, pushed to the extreme where the gas stays at body
temperature while the electrons reach a few electronvolts.

The separation is the entire point. In a low-pressure discharge the electrons
and the neutral gas equilibrate slowly because collisions are rare; at
atmospheric pressure collisions are frequent, yet the electron mass is so much
smaller than the neutral mass that an electron loses only a fraction $\sim
2m_e/M \approx 3\times 10^{-5}$ of its energy in each elastic collision with a
heavy particle. The electrons can therefore sit at $T_e \approx 1$ to $2\,\si{
\electronvolt}$ while the gas stays near $300\,\si{\kelvin}$, a factor of forty
thousand apart in temperature. The hot electrons drive the ionisation and the
radical chemistry; the cold gas carries those radicals to the surface without
thermal damage. A helium jet at a few kilovolts and tens of kilohertz produces
reactive oxygen and nitrogen species, ultraviolet photons, and a modest ion flux
in a plume a few millimetres across.

Take a jet driven at $V = 4\,\si{\kilo\volt}$ peak and $f = 20\,\si{\kilo\hertz}$
through a tube of $2\,\si{\milli\meter}$ bore with a helium flow of $2\,\si{\litre
\per\minute}$. The gas velocity in the tube is the volume flow over the area,
$u = \dot{V}/A = (2\times 10^{-3}/60)/[\pi(10^{-3})^{2}] \approx 11\,\si{\meter
\per\second}$, so a parcel of gas spends $\tau_{\mathrm{res}} = \ell/u \approx
0.02/11 \approx 2\,\si{\milli\second}$ in a two-centimetre active region, long
enough for many discharge cycles at $20\,\si{\kilo\hertz}$ to seed it with
radicals. The power coupled into the discharge is modest, of order a few watts,
so the gas heating is small: a few watts spread over the flowing gas raises its
temperature by $\Delta T = P/(\dot{m}c_p) \approx 3/[(6\times 10^{-6})(5200)]
\approx 100\,\si{\kelvin}$ at most for a helium flow of $\sim 6\,\si{\milli\gram
\per\second}$, and in practice much less because most of the power leaves as
radiation and radical enthalpy, so the plume reaches the target only tens of
degrees above ambient. The touchable plume at a few electronvolts of electron
temperature is the working demonstration of the non-equilibrium the lamp
introduced.

The applications follow the chemistry rather than the heat. Surface treatment
uses the radical flux to raise the surface energy of a polymer so that ink or
adhesive wets it, a step in printing and packaging that replaces wet chemical
priming. Sterilisation uses the reactive oxygen and nitrogen species and the
ultraviolet to kill bacteria on heat-sensitive medical instruments. Plasma
medicine, a field opened in the 2010s, applies the same jet to wounds, where the
reactive species accelerate healing and disinfect without thermal injury; the
dose is set by the radical flux and the exposure time, the same flux-times-time
accounting the etch and deposition tools used, now delivered to tissue rather
than to silicon. The engineering constraint is the reverse of the fusion
divertor's: the jet must deliver enough reactive flux to do the chemistry while
keeping the gas temperature and the ion energy low enough to leave the substrate
undamaged, so the design knob is the power that raises the radical production
without raising the gas heating past the tolerable limit. The atmospheric jet is
the chapter's physics run at the gentlest corner of the map, and it obeys the
same energy-balance and flux accounting as the reactor that fills a building.

\section{A fourth applied vignette: the negative-ion neutral-beam injector}\pathtag{standard}

The tokamak case took the auxiliary heating power as a given. The device that
supplies it is a plasma engine in its own right, and at reactor scale it forces
a change of technology that shows how a scaling law reshapes a design. A neutral
beam heats a magnetised plasma by firing fast neutral atoms across the confining
field, which the field cannot deflect because they carry no charge; the atoms
penetrate to the core, ionise by collision, and deposit their energy where the
plasma needs it. Building the beam means ionising a gas, accelerating the ions
through a large voltage, and then neutralising them again before they reach the
plasma, and the neutralisation step is where the scaling bites.

A beam that must penetrate a dense reactor plasma needs a high energy per
nucleon, of order $1\,\si{\mega\electronvolt}$ for the deuterium beams ITER will
use, far above the tens or hundreds of kilovolts of earlier machines. The
trouble is the neutraliser. A positive ion picks up an electron from a gas cell
with a cross-section that collapses at high energy, so the neutralisation
efficiency of a positive-ion beam falls toward a few percent above a few hundred
kilovolts per nucleon and is hopeless at a mega-electronvolt. A negative ion, an
atom carrying one extra electron bound by only $\sim 0.75\,\si{\electronvolt}$,
sheds that loose electron easily at any energy, so a negative-ion beam
neutralises at $\sim 60\%$ efficiency even at $1\,\si{\mega\electronvolt}$. The
scaling of the charge-exchange cross-section with energy is the whole reason a
reactor-scale injector must start from negative ions, a harder species to make
and to handle, rather than from the positive ions that served every earlier
machine.

Trace the sizes. ITER's injector aims to deliver $\sim 16.5\,\si{\mega\watt}$ of
neutral deuterium at $1\,\si{\mega\electronvolt}$. The neutral power is the beam
current times the voltage times the neutralisation efficiency, $P = \eta_n I V$,
so the negative-ion current needed is $I = P/(\eta_n V) = (1.65\times 10^{7})/
[(0.6)(10^{6})] \approx 27\,\si{\ampere}$ of negative ions, an enormous
current of a fragile species that a stray electron or a background gas atom can
strip before it is accelerated. The source is a caesium-seeded plasma that makes
negative ions on a metal surface; the accelerator pulls them through a stack of
grids at up to a megavolt while a magnetic filter and electron-suppression grids
remove the co-extracted electrons that would otherwise carry away the power and
damage the grids. Each requirement, the caesium chemistry, the megavolt
insulation in vacuum, the electron suppression, is a discipline of its own, and
the test facility that qualifies the ITER injector is itself a large physics
plant.

The vignette carries the same lesson the case studies did, read from the supply
side. The heating power the tokamak balance treated as a number is the output of
a plasma device whose design is set by a cross-section scaling law, exactly as
the divertor's design was set by the parallel-flux scaling and the Hall
thruster's by the electron-confinement scaling. When the physics of a component
changes character with size, as neutralisation does between the laboratory and
the reactor, the engineering must change technology with it, and reading which
scaling forces the change is the judgement the chapter's apparatus is for.

\section{A fifth applied vignette: the dense plasma focus neutron source}\pathtag{standard}

The confinement schemes so far held a plasma against its own pressure for a
long time or compressed it for a short one. A third route confines nothing at
all and lets the plasma's own current pinch it, briefly and violently, into a
hot dense filament that fuses. The dense plasma focus is a table-scale device
that turns a capacitor bank into a burst of neutrons, and it earns a place here
because it runs the chapter's pinch physics, the frozen-in field, and the
Rayleigh-Taylor instability all in a single microsecond event, and because it
is a working neutron and x-ray source rather than a reactor candidate.

A dense plasma focus is two coaxial electrodes, an inner anode inside a
cylindrical cathode cage, in a low-pressure deuterium fill. A capacitor bank
of tens of kilojoules discharges across the gap, breaking the gas down into a
sheath of current that the frozen-in field carries. The self-field of the
current, crossed with the current itself, drives a $\vect{J}\times\vect{B}$
force that sweeps the sheath down the coaxial gun like a snowplough, then runs
it off the end of the anode where it collapses radially onto the axis. The
radial collapse is a z-pinch: the azimuthal field of the axial current
squeezes the plasma into a filament a millimetre across and a centimetre long,
raising its density to $\sim 10^{25}\,\si{\per\meter\cubed}$ and its
temperature to a kilovolt for a few tens of nanoseconds. In that pinch the
deuterium fuses, and a $10\,\si{\kilo\joule}$ device emits $\sim 10^{8}$ to
$10^{9}$ neutrons per shot.

Trace the pinch with the chapter's tools. The pinch balances the plasma
pressure against the magnetic pressure of its own current, the Bennett
relation $\mu_0 I^{2} = 8\pi N k_B(T_e + T_i)$, where $N$ is the line density
of ions per unit length. Take a current $I = 300\,\si{\kilo\ampere}$ at the
pinch and a temperature $k_B(T_e+T_i) \approx 2\,\si{\kilo\electronvolt} =
3.2\times 10^{-16}\,\si{\joule}$. The line density the pinch confines is $N =
\mu_0 I^{2}/[8\pi k_B(T_e+T_i)] = (4\pi\times 10^{-7})(3\times 10^{5})^{2}/
[8\pi(3.2\times 10^{-16})] \approx 1.4\times 10^{19}\,\si{\per\meter}$, so a
one-centimetre pinch holds about $1.4\times 10^{17}$ ions, and at a pinch
radius of $1\,\si{\milli\meter}$ the density is $N/(\pi r^{2}) \approx
1.4\times 10^{19}/[\pi(10^{-3})^{2}] \approx 4.5\times 10^{24}\,\si{\per\meter
\cubed}$, consistent with the quoted pinch density. The confinement time is the
time the pinch holds before it flies apart, the radial Alfv\'en time $\tau
\sim r/v_A$; with the pinch field $B = \mu_0 I/(2\pi r) = (4\pi\times 10^{-7})
(3\times 10^{5})/[2\pi(10^{-3})] \approx 60\,\si{\tesla}$ and a mass density
$\rho = N m_D/(\pi r^{2}) \approx 1.5\times 10^{-2}\,\si{\kilogram\per\meter
\cubed}$, the Alfv\'en speed is $v_A = B/\sqrt{\mu_0\rho} \approx 60/\sqrt{
(4\pi\times 10^{-7})(1.5\times 10^{-2})} \approx 4.4\times 10^{5}\,\si{\meter
\per\second}$, so $\tau \sim 10^{-3}/4.4\times 10^{5} \approx 2\,\si{\nano
\second}$, tens of nanoseconds once the pinch dynamics are included. The pinch
lives and dies on the Alfv\'en time, exactly as the disruption worked example
found for the tokamak, because both are ideal-MHD events.

The device is not a reactor and does not pretend to be, for the same reason
the theta-pinch was not: the pinch is destroyed by the sausage and kink
instabilities the moment it forms, the $m=0$ and $m=1$ modes that a straight
current column cannot avoid, and it is the disruption of the pinch that heats
the ions and drives the neutron burst, a beam-target rather than a thermal
yield. That makes the dense plasma focus a poor path to net energy but an
excellent pulsed source: it is used for neutron radiography, for testing
fusion-relevant materials against neutron and x-ray pulses, and as a compact
source for activation analysis. The engineering trade is the one the whole
chapter has drawn, read at the scale of a bench device. The same pinch that
would have to be held steady for power must here only flash, so the
instabilities that defeat confinement become the mechanism that makes the
source work, and a device that fails as a reactor succeeds as an instrument.

\section{Worked examples}\pathtag{standard}

The operational formulas of the chapter reduce to a small number of
arithmetic moves. Three worked instances fix them in mind; each is the
template for one of the calculation exercises.

\paragraph{Debye length and plasma frequency of a fluorescent lamp.}
Take $n_{e} = 10^{17}\,\si{\per\cubic\meter}$ and $T_{e} = 1\,\si{
\electronvolt}$, so $k_{B}T_{e} = 1.6\times 10^{-19}\,\si{\joule}$. The
Debye length is
\[
\lambda_{D} = \sqrt{\frac{\varepsilon_{0}k_{B}T_{e}}{n_{e}e^{2}}}
= \sqrt{\frac{(8.85\times 10^{-12})(1.6\times 10^{-19})}
{(10^{17})(1.6\times 10^{-19})^{2}}}
\approx 2.4\times 10^{-5}\,\si{\meter},
\]
about $24\,\si{\micro\meter}$. The plasma frequency is
\[
\omega_{p} = \sqrt{\frac{n_{e}e^{2}}{m_{e}\varepsilon_{0}}}
= \sqrt{\frac{(10^{17})(1.6\times 10^{-19})^{2}}
{(9.11\times 10^{-31})(8.85\times 10^{-12})}}
\approx 1.8\times 10^{10}\,\si{\radian\per\second},
\]
or $f_{p} = \omega_{p}/2\pi \approx 2.8\,\si{\giga\hertz}$. Both numbers
agree with the table produced by \texttt{code/plasma\_parameters.py}. The
lamp is a few centimetres across, so $\lambda_{D}$ is smaller than the
device by a factor of order $10^{3}$, and the plasma definition holds.

\paragraph{Cyclotron frequency and Larmor radius in a fusion field.}
A thermal electron at $T = 10\,\si{\kilo\electronvolt}$ has perpendicular
speed $v_{\perp} = \sqrt{2k_{B}T/m_{e}} \approx \sqrt{2(1.6\times
10^{-15})/(9.11\times 10^{-31})} \approx 5.9\times 10^{7}\,\si{\meter\per
\second}$. In a $5\,\si{\tesla}$ field the cyclotron frequency is
$\omega_{c} = eB/m_{e} = (1.6\times 10^{-19})(5)/(9.11\times 10^{-31})
\approx 8.8\times 10^{11}\,\si{\radian\per\second}$, and the Larmor radius
is $r_{L} = v_{\perp}/\omega_{c} \approx 6.7\times 10^{-5}\,\si{\meter}$,
about $67\,\si{\micro\meter}$. A deuteron at the same temperature has
$v_{\perp} = \sqrt{2k_{B}T/m_{D}} \approx 9.8\times 10^{5}\,\si{\meter\per
\second}$ and $r_{L} = m_{D}v_{\perp}/(eB) \approx 4.1\times
10^{-3}\,\si{\meter}$, about $4\,\si{\milli\meter}$. The ion Larmor radius
is roughly $60$ times the electron's, the square root of the mass ratio,
and it sets the spatial scale of the turbulence that limits confinement.

\paragraph{Hall-thruster exit velocity and specific impulse.}
A singly ionised xenon atom ($m = 131\,\text{u} = 2.18\times
10^{-25}\,\si{\kilogram}$) accelerated through a $300\,\si{\volt}$ drop
gains kinetic energy $eV = (1.6\times 10^{-19})(300) = 4.8\times
10^{-17}\,\si{\joule}$, so its exit speed is $v = \sqrt{2eV/m} =
\sqrt{2(4.8\times 10^{-17})/(2.18\times 10^{-25})} \approx 2.1\times
10^{4}\,\si{\meter\per\second}$. The specific impulse is $I_{sp} =
v/g_{0} = 2.1\times 10^{4}/9.81 \approx 2200\,\si{\second}$, an order of
magnitude above a chemical bipropellant and the reason electric
propulsion wins for high-$\Delta v$ missions where propellant mass
dominates the budget.

\paragraph{Electron-cyclotron-resonance heating condition.}
A microwave source heats a magnetised plasma efficiently when its
frequency matches the electron cyclotron frequency, so the rotating
electric field stays in phase with the gyrating electrons and pumps energy
into them turn after turn. The resonance condition is $\omega = \omega_c =
eB/m_e$, which fixes the field at which a given source resonates. For the
$2.45\,\si{\giga\hertz}$ magnetrons used in commercial
electron-cyclotron-resonance (ECR) etch and ion sources, $\omega = 2\pi
(2.45\times 10^{9}) = 1.54\times 10^{10}\,\si{\radian\per\second}$, so the
resonant field is
\[
B_{\mathrm{ECR}} = \frac{m_e\omega}{e}
= \frac{(9.11\times 10^{-31})(1.54\times 10^{10})}{1.6\times 10^{-19}}
\approx 8.75\times 10^{-2}\,\si{\tesla},
\]
about $87.5\,\si{\milli\tesla}$ ($875\,\text{G}$), a field a
permanent-magnet assembly reaches
easily. The resonance is a surface in the chamber, the locus where $B$
passes through $87.5\,\si{\milli\tesla}$; electrons crossing that surface
gain energy and ionise the gas there, which is why ECR sources produce
dense plasma at low pressure without an immersed electrode to erode. The
same arithmetic for the $28\,\si{\giga\hertz}$ gyrotrons used to heat
tokamak plasmas gives a resonant field near $1\,\si{\tesla}$, well inside
the operating field of a fusion device, so the gyrotron beam can be aimed
to deposit its power on a chosen flux surface for current drive and
instability control.

\paragraph{Ambipolar diffusion in a discharge.}
Electrons diffuse faster than ions because they are lighter, but they
cannot run ahead indefinitely: any charge separation sets up an electric
field that retards the electrons and pulls the ions along, so the two
species diffuse together at a shared \engterm{ambipolar} rate. The
ambipolar diffusion coefficient is
\[
D_a = D_i\!\left(1 + \frac{T_e}{T_i}\right),
\]
which for a low-temperature discharge with $T_e \gg T_i$ reduces to $D_a
\approx D_i T_e/T_i$. Take an argon RIE plasma at $10\,\si{\pascal}$ with
$T_e = 3\,\si{\electronvolt}$ and a room-temperature ion population $T_i =
0.03\,\si{\electronvolt}$, so $T_e/T_i = 100$. The free ion diffusion
coefficient at this pressure is of order $D_i \approx 5\times
10^{-2}\,\si{\meter\squared\per\second}$, so $D_a \approx 5\times 10^{-2}
\times 100 = 5\,\si{\meter\squared\per\second}$. The diffusion time across a
chamber half-width $L = 0.1\,\si{\meter}$ is $\tau = L^{2}/D_a \approx
(0.1)^{2}/5 \approx 2\,\si{\milli\second}$, the time the plasma takes to
fill or empty the chamber by transport to the walls. The ambipolar field
that ties the two species together is what carries ions to the wafer at the
controlled flux an etch process depends on; raising $T_e$ raises $D_a$ and
the wall loss, which is one reason a discharge cannot be driven to
arbitrarily high electron temperature without losing its plasma to the
walls.

\paragraph{Dust-grain charging and the floating potential.}
A small grain of dust suspended in a plasma charges until it collects equal
electron and ion fluxes, the same flux-balance that fixes the potential of
any isolated surface. The floating potential of a grain relative to the
plasma is, for a hydrogenic plasma,
\[
\phi_f = -\frac{k_B T_e}{2e}\ln\!\left(\frac{m_i}{2\pi m_e}\right),
\]
the sheath drop derived earlier without the presheath term. For an argon
plasma at $T_e = 3\,\si{\electronvolt}$, the logarithm is $\ln(40\times
1836/2\pi) \approx 9.4$, so $\phi_f \approx -(3/2)(9.4) \approx
-14\,\si{\volt}$. A spherical grain of radius $a = 1\,\si{\micro\meter}$
behaves as a vacuum capacitor $C = 4\pi\varepsilon_0 a = 4\pi(8.85\times
10^{-12})(10^{-6}) \approx 1.1\times 10^{-16}\,\si{\farad}$, so it carries a
charge
\[
Q = C\phi_f \approx (1.1\times 10^{-16})(-14) \approx -1.6\times
10^{-15}\,\si{\coulomb},
\]
about $-10^{4}$ electron charges. The result that a micron grain holds ten
thousand surplus electrons explains the behaviour of dusty plasmas:
like-charged grains repel, and at high grain density the repulsion orders
them into Coulomb crystals; in a semiconductor tool the same charge levitates
particles in the sheath and then drops them onto the wafer when the plasma
extinguishes, a notorious source of defects.
Figure~\ref{fig:vol04ch13:dust-charging} shows the flux balance and the
charged grain inside its screening cloud.

% Figure: dust-grain charging in a plasma. A small isolated grain floats to
% a negative potential where the electron flux it collects equals the ion
% flux. The left panel shows the two fluxes crossing at the floating
% potential; the right panel sketches the negatively charged grain
% surrounded by its ion-rich screening cloud.
\begin{figure}[htbp]
\centering
\begin{tikzpicture}[font=\small]
% left panel: flux balance
\begin{axis}[
  name=flux,
  width=0.52\textwidth, height=0.42\textwidth,
  xlabel={grain potential $\phi_g$ (V)},
  ylabel={collected flux (arb.\ units)},
  xmin=-12, xmax=2, ymin=0, ymax=2.6,
  grid=both, grid style={enggray!20},
  every axis plot/.append style={thick},
  legend pos=north west, legend cell align=left,
]
  % electron flux: Boltzmann-suppressed as grain goes negative
  \addplot[engblue, domain=-12:1, samples=80]
    {2.2*exp(x/3)};
  \addlegendentry{electron flux}
  % ion flux: nearly flat (orbit-limited, weak rise toward negative phi)
  \addplot[engred, domain=-12:1, samples=80]
    {0.35*(1-x/12)};
  \addlegendentry{ion flux}
  % floating potential marker (where they cross, near phi=-5.4 here)
  \draw[dashed, enggray] (axis cs:-5.4,0) -- (axis cs:-5.4,0.55);
  \node[enggray, font=\scriptsize, anchor=south] at (axis cs:-5.4,0.6)
    {$\phi_f$};
\end{axis}
% right panel: charged grain with screening cloud
\begin{scope}[shift={(9.2,1.7)}]
  % screening cloud (ion-rich)
  \fill[engorange!18] (0,0) circle (1.7);
  \node[engorange!80!black, font=\scriptsize, anchor=south]
    at (0,1.75) {ion-rich sheath ($\sim\lambda_D$)};
  % grain
  \fill[enggray!60] (0,0) circle (0.5);
  \node[font=\scriptsize] at (0,0) {$-Q$};
  % a few + charges in the cloud
  \foreach \a in {30,110,200,300}{
    \node[engred, font=\scriptsize] at (\a:1.15) {$+$};}
\end{scope}
\end{tikzpicture}
\caption[Dust-grain charging in a plasma]{Charging of a small isolated grain
immersed in a plasma. Electrons are far more mobile than ions, so a grain
initially collects electrons faster than ions and charges negative. As the
grain potential falls, its electron collection is Boltzmann-suppressed while
its ion collection is nearly flat, and the two fluxes cross at the floating
potential $\phi_f$ (left panel). The grain settles there, carrying a
negative charge $-Q$ set by its capacitance and $\phi_f$, and sits inside
an ion-rich screening cloud a few Debye lengths across (right panel). The
same flux-balance argument fixes the charge on every wall, probe, and
particle in a plasma, and the mutual repulsion of like-charged grains is
what orders the Coulomb crystals seen in dusty laboratory plasmas.}
\label{fig:vol04ch13:dust-charging}
\end{figure}


\paragraph{Diamagnetic drift and the current a pressure gradient drives.}
A pressure gradient across a magnetic field drives a fluid drift even when no
single particle drifts, the \engterm{diamagnetic drift}, because more particles
gyrate through a fixed point from the high-pressure side than from the low. The
drift velocity of species $s$ is $\vect{v}_{Ds} = -\grad p_s\times\vect{B}/
(q_s n_s B^{2})$, opposite for ions and electrons because of the charge, so the
two species drift in opposite directions and carry a net \engterm{diamagnetic
current}. Take the case-study tokamak with an edge pressure gradient that falls
the full core pressure $p = 3.8\times 10^{5}\,\si{\pascal}$ over a scale length
$L_p = 0.1\,\si{\meter}$ at $B = 5.3\,\si{\tesla}$. The electron diamagnetic
speed is
\[
v_{De} = \frac{k_B T_e}{eB L_p}
= \frac{1.92\times 10^{-15}}{(1.6\times 10^{-19})(5.3)(0.1)}
\approx 2.3\times 10^{4}\,\si{\meter\per\second},
\]
using $p_e/n_e = k_B T_e$ and $|\grad p_e|/p_e = 1/L_p$. The diamagnetic current
density is $J_D = |\grad p|/B = (3.8\times 10^{5}/0.1)/5.3 \approx 7.2\times
10^{5}\,\si{\ampere\per\meter\squared}$, a current the plasma carries by virtue
of its confinement alone. This is the current that, crossed with the field,
supplies the $\vect{J}\times\vect{B}$ force that holds the pressure gradient up,
so it is not an accident of the edge but the diamagnetic signature of a confined
plasma: a magnetically confined plasma is always diamagnetic, expelling a little
of the field it sits in, and a magnetic pickup coil outside the plasma reads the
stored pressure through exactly this current. The diamagnetic-loop measurement
of plasma beta rests on this arithmetic.

\paragraph{Magnetic-mirror confinement time set by scattering.}
The mirror worked example earlier found the instantaneous loss fraction; the
confinement time follows from how fast collisions refill the empty loss cone.
A mirror-confined ion is lost when a collision scatters its velocity vector into
the loss cone, which takes a $90^{\circ}$ deflection, so the confinement time is
of order the ninety-degree scattering time times a logarithmic factor for the
mirror ratio, $\tau_{\mathrm{mirror}} \approx \tau_{90}\ln R_m$. Take a mirror
plasma of deuterons at $T = 10\,\si{\kilo\electronvolt}$, $n = 10^{20}\,\si{\per
\meter\cubed}$, mirror ratio $R_m = 8$. The ion-ion collision time at these
conditions is $\tau_{ii} \approx 6\times 10^{-3}\,\si{\second}$ from the
collision-frequency scaling, so $\tau_{\mathrm{mirror}} \approx (6\times
10^{-3})\ln 8 \approx (6\times 10^{-3})(2.1) \approx 1.3\times 10^{-2}\,\si{
\second}$, about thirteen milliseconds. Compare this to the triple-product
requirement: at $n = 10^{20}$ the Lawson criterion needs $\tau_E \approx
3\times 10^{21}/(10^{20}\times 10) \approx 3\,\si{\second}$, more than two
hundred times the mirror confinement time. The mirror falls short of ignition by
two orders of magnitude not because its physics is wrong but because collisional
scattering into the loss cone empties it far faster than fusion can heat it, the
quantitative reason the mirror line did not reach a reactor and closed toroidal
confinement, which has no loss cone, took over.

\paragraph{Child-Langmuir sheath thickness at a biased wall.}
A space-charge-limited sheath in front of a strongly biased electrode obeys
the Child-Langmuir law, the same one that sets the current in a vacuum
diode. The ion current density a sheath of voltage $V_{s}$ and thickness $s$
can pass is
\[
J_{i} = \frac{4}{9}\varepsilon_{0}
\sqrt{\frac{2e}{m_{i}}}\,\frac{V_{s}^{3/2}}{s^{2}}.
\]
Inverting for the thickness ties the sheath width to the bias and the Bohm
flux that feeds it. For the CCP etch case, with $J_{i}\approx 5\,\si{\ampere
\per\meter\squared}$, a sheath voltage $V_{s}\approx 350\,\si{\volt}$, and
the argon-scale ion mass $m_{i} = 6.6\times 10^{-26}\,\si{\kilogram}$,
\[
s = \left[\frac{4}{9}\varepsilon_{0}\sqrt{\frac{2e}{m_{i}}}
\frac{V_{s}^{3/2}}{J_{i}}\right]^{1/2}.
\]
The factor $\sqrt{2e/m_{i}} = \sqrt{2(1.6\times 10^{-19})/(6.6\times
10^{-26})} \approx 2.2\times 10^{3}$, and $V_{s}^{3/2} = 350^{1.5} \approx
6.55\times 10^{3}$, so the bracket is $\tfrac49(8.85\times 10^{-12})(2.2
\times 10^{3})(6.55\times 10^{3})/5 \approx 1.1\times 10^{-5}\,\si{\meter
\squared}$ and $s \approx 3.4\times 10^{-3}\,\si{\meter}$. The
$3.4\,\si{\milli\meter}$ result is approximate, since the
Child-Langmuir form ignores the finite ion entry velocity and the RF
modulation of the sheath, but it lands at the right scale: a few millimetres,
still many Debye lengths thick against the $90\,\si{\micro\meter}$
Debye length the case computed. The lesson the calculation carries is the
scaling: the sheath thickens as $V_{s}^{3/4}$ and thins as the ion current
rises, so deepening the bias to raise ion energy also widens the dark space,
which the chamber gap must accommodate.

\paragraph{Bremsstrahlung versus line radiation and the impurity limit.}
A fusion plasma loses energy to bremsstrahlung at a volumetric rate
$P_{\mathrm{br}} = 1.7\times 10^{-38}\,Z_{\mathrm{eff}}\,n_{e}^{2}\sqrt{T_{e}
/\text{eV}}\,\si{\watt\per\meter\cubed}$, where $Z_{\mathrm{eff}}$ is the
effective charge that weights each ion species by $Z^{2}$. For a pure
hydrogenic plasma $Z_{\mathrm{eff}} = 1$, but a small impurity fraction
raises it sharply because the radiation scales as $Z^{2}$. Take the
case-study plasma, $n_{e} = 10^{20}\,\si{\per\meter\cubed}$ at $T_{e} =
12\,\si{\kilo\electronvolt}$, and add a $1\%$ concentration of tungsten
($Z = 74$) sputtered from the divertor. The tungsten contribution to
$Z_{\mathrm{eff}}$ is $n_{W}Z_{W}^{2}/n_{e} = (0.01)(74)^{2} \approx 55$, so
$Z_{\mathrm{eff}}$ jumps from $1$ to about $56$. The bremsstrahlung loss
rises by the same factor, from the $\sim 9\,\si{\mega\watt}$ the case
computed to $\sim 500\,\si{\mega\watt}$, which exceeds the $170\,\si{\mega
\watt}$ of alpha heating and extinguishes the burn. The arithmetic is the
quantitative reason a fusion plasma tolerates only parts-per-ten-thousand of
high-$Z$ impurity: a contamination level that would be invisible in a
chemical process is fatal to the power balance, and impurity control is a
first-order engineering requirement, not a refinement.

\paragraph{Saha ionisation fraction of an arc plasma.}
A welding or cutting arc runs near atmospheric pressure at $T \approx
15{,}000\,\si{\kelvin}$, hot enough that the Saha equation, not a
non-equilibrium discharge model, sets the ionisation. For an argon arc
($E_{i} = 15.8\,\si{\electronvolt}$) the exponent is $E_{i}/k_{B}T =
(15.8)(1.6\times 10^{-19})/[(1.38\times 10^{-23})(1.5\times 10^{4})] \approx
12.2$, so $\exp(-12.2) \approx 5.0\times 10^{-6}$. The phase-space prefactor
$(m_{e}k_{B}T/2\pi\hbar^{2})^{3/2}$ evaluates to $\approx (9.11\times
10^{-31}\times 1.38\times 10^{-23}\times 1.5\times 10^{4}/[2\pi(1.05\times
10^{-34})^{2}])^{3/2} \approx (2.7\times 10^{18})^{3/2} \approx 4.5\times
10^{27}\,\si{\per\meter\cubed}$. The atmospheric neutral density is $n_{n}
\approx p/(k_{B}T) = 10^{5}/[(1.38\times 10^{-23})(1.5\times 10^{4})]
\approx 4.8\times 10^{23}\,\si{\per\meter\cubed}$. With $n_{i} = n_{e}$, the
Saha relation $n_{e}^{2}/n_{n} = (4.5\times 10^{27})(5.0\times 10^{-6})
\approx 2.3\times 10^{22}$ gives $n_{e} = \sqrt{(2.3\times 10^{22})(4.8\times
10^{23})} \approx 1.0\times 10^{23}\,\si{\per\meter\cubed}$, an ionisation
fraction of order $n_{e}/n_{n}\approx 20\%$. Ionising a fifth of the gas at
$15{,}000\,\si{\kelvin}$ makes the arc a strong conductor, able to
carry the hundreds of amperes a cutting torch draws, and the frequent
collisions at atmospheric pressure hold electrons and heavy particles at one
temperature, so a thermal arc is a genuine equilibrium plasma. A low-pressure
fluorescent lamp, ionised only about one part in a million and with its
electrons far hotter than its gas, sits at the opposite, non-equilibrium
extreme.

\paragraph{Alfv\'en transit time and the disruption timescale.}
The speed at which a magnetic disturbance propagates along a field line sets
the fastest timescale of an ideal-MHD instability. For the case-study
tokamak, $B = 5.3\,\si{\tesla}$ and a 50:50 D-T plasma at $n = 10^{20}\,\si{
\per\meter\cubed}$ has mass density $\rho = n m_{i} \approx (10^{20})(4.2
\times 10^{-27}) \approx 4.2\times 10^{-7}\,\si{\kilogram\per\meter\cubed}$,
so the Alfv\'en speed is
\[
v_{A} = \frac{B}{\sqrt{\mu_{0}\rho}}
= \frac{5.3}{\sqrt{(4\pi\times 10^{-7})(4.2\times 10^{-7})}}
\approx 7.3\times 10^{6}\,\si{\meter\per\second}.
\]
The Alfv\'en transit time across the $a = 2\,\si{\meter}$ minor radius is
$\tau_{A} = a/v_{A} \approx 2/(7.3\times 10^{6}) \approx 2.7\times 10^{-7}\,
\si{\second}$, a fraction of a microsecond. An ideal-MHD kink that goes
unstable grows on a few Alfv\'en times, so the loss-of-confinement phase of
a disruption develops in microseconds, far faster than any mechanical
protection can respond. The thermal quench that follows dumps the stored
energy in a millisecond and the current quench in tens of milliseconds; the
microsecond Alfv\'en time is why disruption mitigation must be triggered by
prediction before the instability grows, not by reaction after it does. The
same $\tau_{A}$ compared against the resistive time $\tau_{R}$ defines the
Lundquist number of the mastery box on reconnection.

\paragraph{Interchange instability growth at a curved plasma boundary.}
Where a magnetic field curves away from the plasma, the boundary behaves like
a heavy fluid sitting on a light one, and it goes unstable exactly as water
poured over air does, the \engterm{interchange} or flute instability that is
the plasma analogue of Rayleigh-Taylor. The effective gravity is the
centrifugal acceleration of the field-line curvature, $g_{\mathrm{eff}} =
v_{\mathrm{th}}^{2}/R_c$, and the growth rate of a boundary perturbation is
$\gamma = \sqrt{g_{\mathrm{eff}}/L_n}$ with $L_n$ the density gradient scale
length, the same $\sqrt{g/L}$ that sets the Rayleigh-Taylor rate. Take a
mirror-confined deuterium plasma at $T = 100\,\si{\electronvolt}$, so the
thermal speed is $v_{\mathrm{th}} = \sqrt{k_B T/m_D} = \sqrt{(1.6\times
10^{-17})/(3.34\times 10^{-27})} \approx 6.9\times 10^{4}\,\si{\meter\per
\second}$, with a field-curvature radius $R_c = 0.5\,\si{\meter}$ and a
density gradient $L_n = 0.05\,\si{\meter}$. The effective gravity is
$g_{\mathrm{eff}} = (6.9\times 10^{4})^{2}/0.5 \approx 9.5\times 10^{9}\,\si{
\meter\per\second\squared}$, so the growth rate is $\gamma = \sqrt{9.5\times
10^{9}/0.05} \approx 4.4\times 10^{5}\,\si{\per\second}$, an e-folding time of
a few microseconds. A boundary that goes unstable in microseconds cannot be
held by a field that curves the wrong way, which is the quantitative reason a
simple magnetic mirror leaks not only through the loss cone but through the
flutes that grow wherever the field is convex toward the plasma. The cure is a
minimum-B geometry, a field that increases in every direction away from the
plasma so the curvature everywhere points the stable way, the principle behind
the baseball-coil and Ioffe-bar mirrors that replaced the simple mirror.

\paragraph{Gyro-Bohm scaling and the confinement a bigger machine buys.}
The confinement time the Lawson balance needs is set by turbulent transport,
and the turbulence has a natural step size, the ion gyroradius, which is why
the confinement scaling carries the name gyro-Bohm. The gyro-Bohm diffusion
coefficient is the Bohm value reduced by the ratio of the gyroradius to the
machine size, $D_{\mathrm{gB}} = (\rho_i/a)\,D_{\mathrm{Bohm}}$ with
$D_{\mathrm{Bohm}} = k_B T_e/(16 e B)$. For the case-study tokamak at $T_e =
12\,\si{\kilo\electronvolt}$ and $B = 5.3\,\si{\tesla}$, the Bohm coefficient
is $D_{\mathrm{Bohm}} = (1.92\times 10^{-15})/[16(1.6\times 10^{-19})(5.3)]
\approx 1.4\times 10^{2}\,\si{\meter\squared\per\second}$. The ion gyroradius is $\rho_i =
m_D v_{\mathrm{th}}/(eB) \approx (3.34\times 10^{-27})(1.07\times 10^{6})/
[(1.6\times 10^{-19})(5.3)] \approx 4.2\times 10^{-3}\,\si{\meter}$, so with
$a = 2\,\si{\meter}$ the reduction factor is $\rho_i/a \approx 2.1\times
10^{-3}$ and $D_{\mathrm{gB}} \approx 0.30\,\si{\meter\squared\per
\second}$. The confinement time it implies is $\tau_E \sim a^{2}/D_{\mathrm{gB}}
\approx 4/0.30 \approx 13\,\si{\second}$, several times
longer than the $3\,\si{\second}$ the machine actually achieves, which is the
honest caution the gyro-Bohm estimate carries: the bare diffusion argument
overstates confinement by about fourfold because the turbulence is not a
simple random walk, and the working $\tau_E$ comes from empirical scaling laws
that fit whole databases of discharges. The one dependable content of gyro-Bohm is
the size scaling: because $D_{\mathrm{gB}}\propto\rho_i/a$ falls as the machine
grows, confinement improves faster than $a^{2}$, which is the theoretical
reason a burning plasma must be large.

\paragraph{Rogowski coil and the plasma current it reads.}
A tokamak measures its plasma current without touching the plasma, by a
\engterm{Rogowski coil}, a solenoid wound uniformly around a loop that encircles
the current. By Amp\`ere's law the line integral of $B$ around the loop is
$\mu_0 I_{\mathrm{enc}}$, and a coil of $N$ turns per unit length and
cross-section $A_c$ threaded by that field produces, when the current changes,
an output voltage $V = \mu_0 N A_c\,\dd I/\dd t$, so the integrated output is
proportional to the enclosed current regardless of where the current sits
inside the loop. Take a coil with $n = 2000$ turns spread around a $2\,\si{
\meter}$ circumference, so $N = 1000\,\si{\per\meter}$, each turn of area $A_c
= 1\,\si{\centi\meter\squared} = 10^{-4}\,\si{\meter\squared}$, reading a
plasma current that ramps to $15\,\si{\mega\ampere}$ in $t_r = 3\,\si{\second}$.
The mean $\dd I/\dd t = 5\times 10^{6}\,\si{\ampere\per\second}$, so the raw
coil voltage is $V = (4\pi\times 10^{-7})(1000)(10^{-4})(5\times 10^{6})
\approx 0.63\,\si{\volt}$, a signal an integrator turns into a current reading
of $\int V\,\dd t/(\mu_0 N A_c)$. The virtue the calculation shows is that the
reading is independent of the plasma's position and shape, since Amp\`ere's law
counts only the enclosed current, which is why the Rogowski coil is the primary
current diagnostic on every tokamak and why the same device, clamped around a
cable, is the clip-on ammeter an electrician uses. The coil measures $\dd I/\dd
t$, so a slow baseline drift in the integrator is its characteristic error, and
a well-built current measurement resets the integrator between shots.

\paragraph{Hoop force and the magnetic pressure on a toroidal field coil.}
The field that confines the plasma pushes just as hard on the coils that make
it, and the force sets the mechanical design of the magnet. A toroidal field
coil carrying the field $B$ feels an outward magnetic pressure $p_B = B^{2}/
(2\mu_0)$ on its inner face, the same magnetic pressure that balances the
plasma, now acting on the conductor. For the case-study field $B = 5.3\,\si{
\tesla}$, $p_B = (5.3)^{2}/(2\times 4\pi\times 10^{-7}) \approx 1.1\times
10^{7}\,\si{\pascal}$, about $110\,\text{bar}$ or $1600\,\text{psi}$, pressing
outward on every square metre of the coil bore. Integrated around the D-shaped
coil this becomes a net outward centring force, the hoop force, of order the
pressure times the bore area: for a coil enclosing a bore roughly $8\,\si{\meter}
\times 12\,\si{\meter}$, a projected area of order $50\,\si{\meter\squared}$,
the outward force approaches $p_B\times 50 \approx 5\times 10^{8}\,\si{\newton}$,
some fifty thousand tonnes, which is why the toroidal coils of a large tokamak
are wedged into a vault against a central column that carries the hoop load in
compression. The magnetic pressure that reads as a comfortable confinement
margin at the plasma is a punishing structural load at the coil, and the same
$B^{2}/2\mu_0$ that fixes the plasma beta fixes the steel a reactor needs, one
more place the fourth-power scaling of fusion power with field is paid for.

\begin{mastery}
The dust-charging estimate above used orbit-limited collection, the same
\engterm{orbital-motion-limited} (OML) theory that underlies the
small-probe Langmuir analysis. When the grain or probe radius is much
smaller than the Debye length, an attracted particle can be collected from
any impact parameter for which its angular momentum and energy still allow
the orbit to reach the surface, so the collection cross-section grows with
the attractive potential rather than saturating at the geometric area. For
a spherical collector the OML ion current scales as $I_i \propto a^{2}(1 -
e\phi/k_B T_i)$ in the attractive branch, a linear rise in collected
current with collector voltage rather than the flat saturation a planar
sheath would give. The practical consequence is that a Langmuir probe thin
enough to sit in the OML regime has no true ion saturation, and the
analyst must fit the linear OML form rather than read a flat plateau; using
the wrong collection model is one of the common ways a probe measurement
reports the wrong density. The same theory, applied to a spacecraft in the
ionosphere, predicts the floating potential a satellite reaches and the
charge it must shed to avoid arc damage to its surfaces.
\end{mastery}

\begin{mastery}
A high-intensity laser pulse pushes on a plasma through a force that survives
even after averaging over the fast optical oscillation, the \engterm{ponderomotive
force}. An electron in an oscillating field of non-uniform amplitude is driven
harder where the field is strong, so over a cycle it is pushed toward the weak
field, giving a time-averaged force $\vect{F}_p = -(e^{2}/4m_e\omega^{2})\grad
E^{2}$ proportional to the gradient of the field intensity. The force is
independent of the sign of the charge in the sense that it always expels
particles from the intense region, and it scales inversely with mass, so it acts
on the light electrons and leaves the heavy ions behind, setting up a charge
separation and a large electrostatic field. Two engineering consequences follow.
In laser-driven inertial fusion the ponderomotive force of the intense drive
light steepens density profiles and drives parametric instabilities that scatter
the drive and preheat the fuel, effects that the pulse shaping must suppress. In
laser-wakefield acceleration a short intense pulse expels electrons as it passes,
leaving a trailing plasma wave whose longitudinal field can reach $\sim
100\,\si{\giga\volt\per\meter}$, thousands of times the field a conventional
radio-frequency accelerator sustains, so a metre of plasma can do what a
kilometre of copper cavity does. The same plasma that must be held off the wall
in a fusion device becomes, in the wakefield accelerator, the accelerating
structure itself, because it can carry electric fields no solid material could
survive. Reading the ponderomotive force is how an engineer decides whether an
intense beam will push a plasma apart, drive an instability, or build an
accelerating wave.
\end{mastery}

\section{Failure: instabilities, disruption, sputtering}\pathtag{core}

Plasma instabilities are the working failure modes of every
magnetically confined plasma. They are the high-stakes content of
the chapter for fusion engineering and the working failure mode of
many industrial plasma devices.

The \engterm{kink instability} arises when the plasma current in a
tokamak exceeds a stability threshold (the Kruskal-Shafranov limit,
$q > 1$ at the edge of the plasma, where $q$ is the safety factor
that measures the helical pitch of the field lines). The plasma
deforms into a helical mode that grows on the Alfv\'en time scale;
in a fusion device, the consequence is loss of confinement and a
\engterm{disruption}, in which the plasma current rapidly extinguishes,
the magnetic energy is dumped, and large forces act on the vessel.
ITER disruptions are designed to be benign (mitigation through
massive-gas injection); a worst-case unmitigated disruption would
deposit tens of megajoules on vessel surfaces in milliseconds. The
engineering discipline of disruption mitigation is the working
content of much current fusion-engineering research.

The \engterm{edge-localised mode} (ELM) is a quasi-periodic burst at
the plasma edge in H-mode (high-confinement) tokamak operation.
Each ELM deposits a few percent of the stored plasma energy on the
divertor surfaces in $\sim 1\,\si{\milli\second}$. For ITER and
beyond, unmitigated ELMs would erode the divertor too rapidly for
practical fusion-power plants; ELM control by resonant magnetic
perturbations or by pellet pacing is a major engineering effort. The
mechanism ties the ELM to the same edge physics the H-mode buys. The
high-confinement mode builds a steep pressure pedestal at the plasma
edge, a narrow region where the pressure gradient and the associated
edge current rise sharply; that gradient drives a coupled
pressure-and-current instability, the peeling-ballooning mode, once the
pedestal climbs past a stability boundary. The mode relaxes the pedestal
in a burst, expels the few percent of stored energy, and the pedestal
then rebuilds, giving the quasi-periodic cycle. The engineering problem
is that the same steep pedestal that makes H-mode worth having is what
drives the ELM, so ELM control cannot simply flatten the edge without
losing the confinement; it must instead trigger many small ELMs before
the pedestal reaches the dangerous amplitude, by pellet pacing, or hold
the pedestal just below the boundary with a small resonant magnetic
perturbation that leaks enough edge transport to cap the gradient. The
target of a reactor-scale device is to keep each energy pulse below the
tungsten melt threshold, which for the $\sim 1\,\si{\milli\second}$
timescale is of order a fraction of a megajoule per square metre, and an
unmitigated ELM on a reactor would exceed that by more than an order of
magnitude.

\engterm{Sputtering} is the erosion of plasma-facing surfaces by
energetic-ion bombardment. In a fusion device the divertor (the
specialised inner surface that receives the heat flux from the
plasma exhaust) experiences ion fluxes of $\sim 10^{23}\,\si{\per
\meter\squared\per\second}$ at energies of tens to hundreds of
electronvolts; the resulting sputter-erosion rate of tungsten (the
chosen ITER divertor material for its high atomic number and low
sputter yield) is $\sim 1\,\si{\nano\meter\per\second}$ in normal
operation. The divertor lifetime under continuous high-power
operation is one of the open engineering questions of demonstration
fusion reactors.

In industrial plasmas, sputtering is sometimes deliberate (sputter
deposition of thin films onto a substrate by ion bombardment of a
target) and sometimes destructive (etching of unintended chamber
surfaces, contamination of products by sputtered material). The
chemistry of the working plasma determines which: high-Z ion
bombardment of a low-Z surface sputters efficiently; the reverse is
inefficient. Plasma-processing equipment is engineered to keep the
working surfaces ion-bombarded and the unwanted surfaces protected,
typically through magnetic field configuration and the geometry of
the chamber.

\begin{principle}[Confinement is always a stability problem]
Every successful plasma confinement scheme is the solution to a
stability problem: a magnetically confined plasma sits at a saddle
point of an effective free energy, and engineering interventions
must continually suppress the instabilities that would carry the
plasma off the saddle. Plasma engineering at industrial and fusion
scale is the engineering of active stabilisation against an
unforgiving palette of instabilities.
\end{principle}

\begin{estimation}
\textbf{Question.} A plasma globe with a $10\,\si{\centi\meter}$
radius spherical bulb operates at $25\,\si{\kilo\hertz}$ with a peak
voltage of about $5\,\si{\kilo\volt}$ between the central electrode
and the glass envelope. Estimate the electron temperature in the
filaments, the electron density, and the Debye length, from
operational observations.

\textbf{Estimate, before reading on.} State a factor-of-three answer
for $T_{e}$, $n_{e}$, and $\lambda_{D}$ before reading the working.

\textbf{Working.} The discharge filaments in a plasma globe are
filled with low-pressure xenon, argon, or neon at $\sim 100\,\si{
\pascal}$ ($\sim 1\,\text{torr}$). The applied voltage divided by
the bulb radius gives a field of $\sim 50\,\si{\kilo\volt\per\meter}$;
the electron drift in the field, balanced against ionisation
threshold ($\sim 15\,\si{\electronvolt}$ for argon), gives electron
temperatures of $1$ to a few electronvolts. The visible filament
width is $\sim 5\,\si{\milli\meter}$, comparable to the gas mean
free path at these conditions. The discharge current is small
($\sim 10\,\si{\milli\ampere}$ total), so $n_{e}\sim n_{e}/V \sim
I/(eAv_{e})$ with $A\sim 10^{-5}\,\si{\meter\squared}$ and
$v_{e}\sim 10^{6}\,\si{\meter\per\second}$ gives $n_{e}\sim
10^{16}\,\si{\per\meter\cubed}$. The Debye length is then
$\lambda_{D}\sim 75\,\si{\micro\meter}$, comparable to the resolved
filament structure under magnification.

\textbf{Postmortem.} The plasma globe is a working
laboratory-scale example of the operational definitions: $T_{e}$,
$n_{e}$, $\lambda_{D}$ are all observable order-of-magnitude
quantities in a $\$20$ household device. A reader who guessed within
a factor of ten on each was already in the working range; the
filament structure visible to the naked eye corresponds to spatial
scales much larger than $\lambda_{D}$, consistent with the
collective-behaviour definition of a plasma.
\end{estimation}

\section{Project}\pathtag{core}

\begin{project}[Characterise a plasma-globe discharge]
Observe and measure the discharge structure of a household plasma
globe and use the observations to estimate plasma parameters.

\textbf{Materials.}
A commercially available plasma globe (a $10$ to $15$-centimetre
diameter spherical bulb with central electrode, available for
$\$15$ to $\$50$); a smartphone or webcam capable of $30$- or
$60$-fps video and macro focusing; a metre rule; a multimeter; a
darkened room.

\textbf{Procedure.}
\begin{enumerate}
  \item Photograph the filaments in a darkened room with the camera
    on a tripod. Use the macro focus to image the discharge structure
    near the central electrode.
  \item Measure the diameter of the central electrode (typically
    $\sim 1\,\si{\centi\meter}$) and the bulb (typically $\sim
    10\,\si{\centi\meter}$).
  \item Estimate the visible filament width from a calibration on
    the photograph using the bulb diameter as the reference.
  \item Place your hand near the bulb surface. Observe how the
    filaments are drawn toward your hand. Estimate the displacement
    distance and time response.
  \item Measure the audible frequency of the discharge (if any) with
    a smartphone-based spectrum analyser.
  \item Estimate the operating power by measuring the supply current
    and voltage at the wall outlet (with appropriate caution).
\end{enumerate}

\textbf{Report.} A one-page report with the photographs, the
measured filament dimensions, the estimated electron temperature,
electron density, Debye length, and a discussion of why the visible
filament width is observable while the Debye sphere is not. Expected:
$T_{e}\sim 1$ to a few $\si{\electronvolt}$, $n_{e}\sim
10^{16}$ to $10^{17}\,\si{\per\meter\cubed}$, $\lambda_{D}\sim 30$ to
$100\,\si{\micro\meter}$. The hand-drawn-filament observation
illustrates the $E\times B$ drift and the response of the plasma to
external charged surfaces.
\end{project}

\section{Exercises}

\subsection*{Calculation}\pathtag{core}

\begin{exercise}
Compute the Debye length and plasma frequency of a fluorescent-lamp
plasma with $n_{e} = 10^{17}\,\si{\per\meter\cubed}$ and $T_{e} =
1\,\si{\electronvolt}$.
\paragraph{Solution sketch.} With $k_{B}T_{e} = 1.6\times 10^{-19}\,\si{
\joule}$, $\lambda_{D} = \sqrt{\varepsilon_{0}k_{B}T_{e}/(n_{e}e^{2})}
\approx 2.4\times 10^{-5}\,\si{\meter}$ ($24\,\si{\micro\meter}$), and
$\omega_{p} = \sqrt{n_{e}e^{2}/(m_{e}\varepsilon_{0})} \approx 1.8\times
10^{10}\,\si{\radian\per\second}$, so $f_{p}\approx 2.8\,\si{\giga\hertz}$.
This is the first worked example of Section~13.8; the lamp dimension is
$\sim 10^{3}\lambda_{D}$, confirming the plasma definition.
\end{exercise}

\begin{exercise}
Compute the cyclotron frequency and Larmor radius of a thermal
electron ($T = 10\,\si{\kilo\electronvolt}$) in a $5\,\si{\tesla}$
magnetic field. Repeat for a deuterium ion at the same temperature.
\paragraph{Solution sketch.} Electron: $v_{\perp} = \sqrt{2k_{B}T/m_{e}}
\approx 5.9\times 10^{7}\,\si{\meter\per\second}$, $\omega_{c} = eB/m_{e}
\approx 8.8\times 10^{11}\,\si{\radian\per\second}$, $r_{L} =
v_{\perp}/\omega_{c} \approx 67\,\si{\micro\meter}$. Deuteron: $v_{\perp}
\approx 9.8\times 10^{5}\,\si{\meter\per\second}$, $\omega_{c} = eB/m_{D}
\approx 2.4\times 10^{8}\,\si{\radian\per\second}$, $r_{L}\approx
4\,\si{\milli\meter}$. The ion radius exceeds the electron's by roughly
$\sqrt{m_{D}/m_{e}}\approx 60$. Worked in full in Section~13.8.
\end{exercise}

\begin{exercise}
Compute the $E\times B$ drift velocity in a $1\,\si{\tesla}$
magnetic field with a perpendicular electric field of $10\,\si{\kilo
\volt\per\meter}$. State the magnitude and direction relative to
$\vect{E}$ and $\vect{B}$.
\paragraph{Solution sketch.} The magnitude is $v_{E\times B} = E/B =
10^{4}/1 = 10^{4}\,\si{\meter\per\second}$, independent of charge and
mass. The direction is along $\vect{E}\times\vect{B}$, perpendicular to
both fields; with $\vect{E}$ and $\vect{B}$ given, take the cross product
in the stated coordinate frame. Both ions and electrons drift the same
way, so the $E\times B$ drift carries no net current.
\end{exercise}

\begin{exercise}
Compute the Alfv\'en speed and the magnetic pressure in an
ITER-relevant plasma: $B = 5\,\si{\tesla}$, $n_{e} = 10^{20}\,\si{\per
\meter\cubed}$, deuterium-tritium mixture.
\paragraph{Solution sketch.} A 50:50 D-T mixture has mean ion mass
$\approx 2.5\,\text{u} = 4.2\times 10^{-27}\,\si{\kilogram}$, so the mass
density is $\rho = n_{i}m_{i} \approx (10^{20})(4.2\times 10^{-27})
\approx 4.2\times 10^{-7}\,\si{\kilogram\per\cubic\meter}$. The Alfv\'en
speed is $v_{A} = B/\sqrt{\mu_{0}\rho} = 5/\sqrt{(4\pi\times 10^{-7})
(4.2\times 10^{-7})} \approx 6.9\times 10^{6}\,\si{\meter\per\second}$.
The magnetic pressure is $p_{B} = B^{2}/(2\mu_{0}) = 25/(2\times
4\pi\times 10^{-7}) \approx 1.0\times 10^{7}\,\si{\pascal}$, about
$100\,\text{atm}$, which the field must hold against the plasma pressure.
\end{exercise}

\begin{exercise}
Compute the plasma frequency of the F2 layer of the ionosphere at
peak density $10^{12}\,\si{\per\meter\cubed}$. Identify the highest
radio frequency that can be reflected off this layer.
\paragraph{Solution sketch.} $\omega_{p} = \sqrt{n_{e}e^{2}/(m_{e}
\varepsilon_{0})} = \sqrt{(10^{12})(1.6\times 10^{-19})^{2}/((9.11\times
10^{-31})(8.85\times 10^{-12}))} \approx 5.6\times 10^{7}\,\si{\radian
\per\second}$, so $f_{p}\approx 9\,\si{\mega\hertz}$. Waves below $f_{p}$
reflect; the highest reflected frequency at vertical incidence is the
critical frequency $\approx 9\,\si{\mega\hertz}$. Oblique incidence
raises the usable frequency by the secant of the incidence angle, the
basis of long-range HF skywave links.
\end{exercise}

\begin{exercise}
A xenon Hall thruster operates at $300\,\si{\volt}$ discharge voltage.
Compute the exit velocity of singly-ionised xenon and the
specific impulse.
\paragraph{Solution sketch.} Energy balance $eV = \tfrac12 m v^{2}$ with
$m = 131\,\text{u} = 2.18\times 10^{-25}\,\si{\kilogram}$ gives $v =
\sqrt{2eV/m} = \sqrt{2(1.6\times 10^{-19})(300)/(2.18\times 10^{-25})}
\approx 2.1\times 10^{4}\,\si{\meter\per\second}$. The specific impulse
is $I_{sp} = v/g_{0} \approx 2.1\times 10^{4}/9.81 \approx 2200\,\si{
\second}$. Worked in Section~13.8; real thrusters reach a somewhat lower
effective exit velocity because not all ions fall through the full
voltage.
\end{exercise}

\subsection*{Derivation}\pathtag{standard}

\begin{exercise}
Derive the Debye length formula $\lambda_{D} = \sqrt{\varepsilon_{0}
k_{B}T_{e}/(n_{e}e^{2})}$ from the linearised Poisson-Boltzmann
equation for a test charge inserted into a plasma.
\paragraph{Solution sketch.} Start from Poisson's equation $\nabla^{2}\phi
= -\rho/\varepsilon_{0}$ with the Boltzmann densities $n_{e} = n_{0}
\exp(e\phi/k_{B}T_{e})$ and $n_{i} = n_{0}\exp(-e\phi/k_{B}T_{i})$.
Linearise for $e\phi \ll k_{B}T$: $\rho \approx -n_{0}e^{2}\phi(1/k_{B}
T_{e} + 1/k_{B}T_{i})/1$, giving $\nabla^{2}\phi = \phi/\lambda_{D}^{2}$
with $1/\lambda_{D}^{2} = n_{0}e^{2}(T_{e}^{-1}+T_{i}^{-1})/(\varepsilon_{0}
k_{B})$. The spherically symmetric solution is $\phi = (q/4\pi\varepsilon_{0}
r)e^{-r/\lambda_{D}}$. Keeping only the mobile electrons recovers
$\lambda_{D} = \sqrt{\varepsilon_{0}k_{B}T_{e}/(n_{e}e^{2})}$.
\end{exercise}

\begin{exercise}
Derive the plasma frequency $\omega_{p} = \sqrt{n_{e}e^{2}/
(m_{e}\varepsilon_{0})}$ by considering a small charge separation in
a cold plasma and computing the restoring force on the displaced
electrons.
\paragraph{Solution sketch.} Displace a slab of electrons by $x$ relative
to the fixed ions. The exposed surface charge gives a uniform field $E =
n_{e}ex/\varepsilon_{0}$ (a parallel-plate result), so each electron feels
$m_{e}\ddot{x} = -eE = -(n_{e}e^{2}/\varepsilon_{0})x$. This is simple
harmonic motion with $\omega_{p}^{2} = n_{e}e^{2}/(m_{e}\varepsilon_{0})$,
hence $\omega_{p} = \sqrt{n_{e}e^{2}/(m_{e}\varepsilon_{0})}$. The ions
are taken as a fixed background because $\omega_{p}\propto m^{-1/2}$ makes
the ion oscillation $\sqrt{m_{i}/m_{e}}$ times slower.
\end{exercise}

\begin{exercise}
Derive the $E\times B$ drift formula by setting up the equation of
motion of a charged particle in crossed $\vect{E}$ and $\vect{B}$
fields, averaging over the gyration period.
\paragraph{Solution sketch.} Write $m\dot{\vect{v}} = q(\vect{E} +
\vect{v}\times\vect{B})$ and split $\vect{v} = \vect{v}_{d} +
\vect{v}_{g}$ into a constant drift plus gyration. Choosing $\vect{v}_{d}$
so that $q(\vect{E} + \vect{v}_{d}\times\vect{B}) = 0$ removes the electric
force from the gyration equation, leaving pure circular motion for
$\vect{v}_{g}$. Crossing the condition with $\vect{B}$ and using the BAC-CAB
identity with $\vect{v}_{d}\cdot\vect{B} = 0$ gives $\vect{v}_{d} =
\vect{E}\times\vect{B}/B^{2}$. The gyration averages to zero over a period,
so the guiding centre moves at exactly this drift.
\end{exercise}

\begin{exercise}
Derive the frozen-in-flux theorem from the ideal-MHD Ohm's law
$\vect{E} + \vect{v}\times \vect{B} = 0$ and Faraday's law.
\paragraph{Solution sketch.} Substitute $\vect{E} = -\vect{v}\times
\vect{B}$ into Faraday's law $\partial\vect{B}/\partial t = -\curl\vect{E}$
to get $\partial\vect{B}/\partial t = \curl(\vect{v}\times\vect{B})$.
Consider the flux $\Phi = \int_{S}\vect{B}\cdot\dd\vect{A}$ through a
surface moving with the fluid. Its total rate of change has two terms, the
local $\partial\vect{B}/\partial t$ and the boundary-motion term; combining
them with the identity above and Stokes' theorem shows $\dd\Phi/\dd t = 0$.
The flux through any comoving loop is constant, so field lines move with
the fluid.
\end{exercise}

\subsection*{Estimation}\pathtag{core}

\begin{exercise}
Estimate the fraction of the Earth's atmosphere that is ionised at
any given time, dominated by the ionosphere. State the assumptions
about ionosphere density, thickness, and total atmospheric mass.
\paragraph{Solution sketch.} Ionosphere: $n_{e}\sim 10^{11}\,\si{\per
\cubic\meter}$ over a $\sim 200\,\si{\kilo\meter}$ shell of Earth-radius
area $\sim 5\times 10^{14}\,\si{\meter\squared}$, giving $\sim
10^{31}$ free electrons, hence $\sim 10^{31}$ ionised atoms. Total
atmosphere: mass $\sim 5\times 10^{18}\,\si{\kilogram}$ of mean molar
mass $\sim 29\,\si{\gram\per\mole}$ is $\sim 10^{44}$ molecules. The
ionised fraction is therefore $\sim 10^{31}/10^{44} = 10^{-13}$. The
atmosphere is, to overwhelming precision, a neutral gas; the plasma lives
only in the thin upper shell.
\end{exercise}

\begin{exercise}
Estimate the propellant mass a Hall-effect thruster needs to deliver
$1000\,\si{\meter\per\second}$ of $\Delta v$ to a $500\,\si{\kilo
\gram}$ spacecraft, at specific impulse $3000\,\si{\second}$. Compare
to a chemical bipropellant (Isp $300\,\si{\second}$).
\paragraph{Solution sketch.} Tsiolkovsky: $m_{p} = m_{0}(1 -
e^{-\Delta v/(I_{sp}g_{0})})$. Electric: $I_{sp}g_{0} = 3000\times 9.81
\approx 2.94\times 10^{4}\,\si{\meter\per\second}$, so $\Delta v/(I_{sp}
g_{0}) = 1000/29400 \approx 0.034$ and $m_{p}\approx 500\times 0.0334
\approx 17\,\si{\kilogram}$. Chemical: $I_{sp}g_{0}\approx 2940\,\si{\meter
\per\second}$, ratio $0.34$, $m_{p}\approx 500\times 0.288 \approx
144\,\si{\kilogram}$. The electric thruster needs roughly one eighth the
propellant; the cost is long thrust duration at low thrust.
\end{exercise}

\begin{exercise}
Estimate the power required to ionise the working gas of a
fluorescent lamp, given $n_{e}\sim 10^{17}\,\si{\per\meter\cubed}$,
chamber volume $\sim 10^{-3}\,\si{\meter\cubed}$, and a
recombination rate $\sim 10^{-13}\,\si{\meter\cubed\per\second}$.
\paragraph{Solution sketch.} In steady state the ionisation rate balances
the recombination rate $R = \alpha n_{e}^{2}V = (10^{-13})(10^{17})^{2}
(10^{-3}) \approx 10^{18}\,\si{\per\second}$ ion-electron pairs recombining
per second. Each ionisation costs at least the ionisation energy, taken as
$\sim 15\,\si{\electronvolt}\approx 2.4\times 10^{-18}\,\si{\joule}$, so
the minimum power to sustain the plasma is $P\sim (10^{18})(2.4\times
10^{-18})\approx 2\,\si{\watt}$. The actual lamp draws more (tens of
watts) because most input power goes to excitation and radiation, not bare
ionisation.
\end{exercise}

\begin{exercise}
Estimate the disruption energy of an ITER full-current termination,
given a plasma current of $15\,\si{\mega\ampere}$ and a plasma
inductance of $\sim 10\,\si{\micro\henry}$. Compare to the kinetic
energy of a large lorry at $100\,\si{\kilo\meter\per\hour}$.
\paragraph{Solution sketch.} The stored magnetic energy is $W = \tfrac12
LI^{2} = \tfrac12(10^{-5})(15\times 10^{6})^{2} \approx 1.1\times
10^{9}\,\si{\joule}$, about $1\,\si{\giga\joule}$. A $40\,\si{\tonne}$
lorry at $100\,\si{\kilo\meter\per\hour}$ ($28\,\si{\meter\per\second}$)
carries $\tfrac12 m v^{2} = \tfrac12(4\times 10^{4})(28)^{2} \approx
1.6\times 10^{7}\,\si{\joule}$. The disruption energy is roughly $70$
lorries' worth, dumped in milliseconds, which is why disruption mitigation
is a first-order engineering problem for ITER.
\end{exercise}

\subsection*{Simulation}\pathtag{mastery}

\begin{exercise}
Implement a Boris-pusher single-particle simulator for a charged
particle in uniform $\vect{E}$ and $\vect{B}$ fields. Verify the
analytical $E\times B$ drift on a $10\,000$-step run and identify
the time step at which the integrator becomes inaccurate.
\paragraph{Solution sketch.} The reference implementation is
\texttt{code/boris\_pusher.py}: a Boris update splits each step into a
half electric kick, a magnetic rotation, and a second half kick. Averaging
the guiding-centre position over an integer gyro-period removes the
gyration and recovers $v_{E\times B} = E/B$ to better than $0.1\%$ at
$\sim 50$ steps per gyro-period. Accuracy degrades sharply once $\omega_{c}
\Delta t \gtrsim 0.3$ (fewer than $\sim 20$ steps per period): the rotation
angle per step grows, the orbit radius is mis-resolved, and the measured
drift drifts away from $E/B$. The energy in a pure magnetic field stays
bounded because the Boris rotation is exactly norm-preserving.
\end{exercise}

\begin{exercise}
Implement a one-dimensional particle-in-cell (PIC) plasma simulator
with $10^{4}$ electrons in a fixed-ion background. Excite the plasma
oscillation by an initial density perturbation and verify the
plasma frequency from the simulation output.
\paragraph{Solution sketch.} Place $10^{4}$ electrons on a periodic line
with a sinusoidal position perturbation against a uniform fixed-ion
background; deposit charge to a grid, solve Poisson's equation by FFT for
the field, and advance positions and velocities with a leapfrog step.
The total field energy oscillates; a discrete Fourier transform of its
time series shows a peak at $2\omega_{p}$ (the energy oscillates at twice
the field frequency), so the recovered $\omega_{p} = \sqrt{n_{e}e^{2}/
(m_{e}\varepsilon_{0})}$ matches the input density. Resolve at least
$\sim 10$ steps per plasma period and keep the grid spacing below a Debye
length, or the simulation heats numerically.
\end{exercise}

\subsection*{Design}\pathtag{standard}

\begin{exercise}
Design a Langmuir probe for an industrial RIE plasma operating at
$10\,\si{\pascal}$ argon. Specify the probe tip dimensions, the
sweep voltage range, the necessary sweep rate, and the expected
$T_{e}$ and $n_{e}$ ranges.
\paragraph{Solution sketch.} Tip: a tungsten wire $\sim 0.1\,\si{\milli
\meter}$ diameter, exposed length $\sim 5\,\si{\milli\meter}$, small
compared to the chamber but large compared to $\lambda_{D}\sim 10$ to
$40\,\si{\micro\meter}$, so thin-sheath theory applies. Sweep range:
$-60$ to $+30\,\si{\volt}$ around the expected plasma potential to capture
both saturation branches. Sweep rate: fast enough to beat drift yet slow
enough to resolve the exponential, $\sim 1\,\si{\milli\second}$ per sweep
($\sim 100\,\si{\kilo\volt\per\second}$). Expect $T_{e}\sim 2$ to
$5\,\si{\electronvolt}$ and $n_{e}\sim 10^{16}$ to $10^{17}\,\si{\per
\cubic\meter}$ for an RIE argon plasma. The analysis follows
\texttt{code/langmuir\_probe.py}: slope of $\ln(I-I_{\mathrm{ion}})$ for
$T_{e}$, saturation currents for $n_{e}$.
\end{exercise}

\begin{exercise}
Design a Hall thruster for a $1000\,\si{\kilo\gram}$ geostationary
satellite station-keeping over $15$ years. Specify the thruster
power, the propellant mass, the specific impulse, and the propellant
storage volume.
\paragraph{Solution sketch.} Station-keeping needs $\sim 50\,\si{\meter
\per\second}$ of $\Delta v$ per year, so $\sim 750\,\si{\meter\per\second}$
over $15$ years. At $I_{sp}\sim 1600\,\si{\second}$ ($I_{sp}g_{0}\approx
1.57\times 10^{4}\,\si{\meter\per\second}$), Tsiolkovsky gives $m_{p}
\approx 1000(1 - e^{-750/15700}) \approx 47\,\si{\kilogram}$ of xenon.
Liquid xenon stored at $\sim 1500\,\si{\kilogram\per\cubic\meter}$ needs
$\sim 0.03\,\si{\meter\cubed}$, a $\sim 35\,\si{\centi\meter}$ tank.
A $1.5\,\si{\kilo\watt}$ Hall thruster at $\sim 50\,\si{\milli\newton}$
thrust accumulates the impulse over many short burns; power and tank
volume are the design drivers.
\end{exercise}

\begin{exercise}
Design a magnetic confinement geometry for a small laboratory
plasma. Identify whether a magnetic-mirror, a stellarator, or a
small tokamak fits the laboratory's space and budget constraints,
and specify the expected $\beta$ and confinement time.
\paragraph{Solution sketch.} For a teaching laboratory with a $\sim 1\,
\si{\meter}$ footprint and a modest budget, a magnetic mirror or a simple
linear device is the realistic choice: a tokamak needs a transformer,
toroidal and poloidal coils, and disruption-safe construction, and a
stellarator needs precision-shaped coils, both beyond a teaching budget.
Expect low $\beta\sim 10^{-3}$ to $10^{-2}$ (magnetically dominated) and a
short confinement time $\tau_{E}\sim 10\,\si{\micro\second}$ to
$1\,\si{\milli\second}$; the device demonstrates confinement and
diagnostics, not net fusion. State the field strength, coil current, and
vacuum requirement as the binding constraints.
\end{exercise}

\subsection*{Diagnosis and reverse engineering}\pathtag{standard}

\begin{exercise}
A plasma-etch process is producing tapered side walls rather than
the specified vertical walls. Identify the candidate causes (excess
photoresist erosion, inadequate ion-energy directionality, wrong
chemistry choice) and the supplementary measurement that
distinguishes them.
\paragraph{Solution sketch.} Tapered walls mean the etch is gaining a
lateral (isotropic) component. Mask erosion shows as a receding, rounded
mask edge in cross-section SEM and correlates with mask-to-film selectivity;
measure the remaining mask thickness. Inadequate ion directionality
follows from too low a sheath voltage or too high a chamber pressure
(more ion-neutral collisions in the sheath); raise the RF bias or lower
the pressure and watch the taper, and confirm with a sheath-voltage or
ion-energy measurement. Wrong chemistry shows as a chemical (isotropic)
etch component independent of bias; switch to a more polymer-forming gas
mix and check the sidewall passivation. The single most diagnostic
supplementary measurement is the bias-voltage dependence of the taper
angle.
\end{exercise}

\begin{exercise}
A fluorescent lamp's lumen output drops below specification after
$10\,000$ hours. Identify the candidate failure modes (mercury
depletion, phosphor degradation, electrode erosion) and the
diagnostic test for each.
\paragraph{Solution sketch.} Mercury depletion (mercury consumed into the
glass and electrodes) shows as a drop in the $254\,\si{\nano\meter}$ UV
output and a colour shift; measure the spectral power distribution and the
UV line intensity. Phosphor degradation shows as reduced visible output at
unchanged UV excitation and as visible darkening of the coating; compare
visible-to-UV conversion efficiency. Electrode erosion shows as end
blackening, rising strike voltage, and flicker; measure the strike voltage
and inspect the electrode ends. The distinguishing test is whether the UV
line is still strong (then the fault is downstream in the phosphor) or
weak (then the fault is upstream in the discharge or mercury inventory).
\end{exercise}

\subsection*{Failure analysis}\pathtag{core}

\begin{exercise}
A tokamak experiences an unmitigated disruption causing $50\,\si{
\mega\joule}$ of plasma magnetic energy to be deposited on the
first wall in $\sim 10\,\si{\milli\second}$. Identify the failure
modes that would result (vessel deformation, vessel-wall melting,
runaway electron beam impact) and the engineering controls that
would mitigate each.
\paragraph{Solution sketch.} The energy density is enormous: $50\,\si{
\mega\joule}$ over $\sim 10\,\si{\meter\squared}$ in $10\,\si{\milli
\second}$ is $\sim 5\,\si{\mega\watt\per\centi\meter\squared}$ at the
surface. Vessel deformation comes from eddy currents and halo currents
crossing $\vect{B}$, giving large $\vect{J}\times\vect{B}$ forces; control
with mechanical bracing and conducting structures sized for the force.
Wall melting comes from the localised heat pulse exceeding the material's
melt threshold; control by spreading the load and by massive-gas or
shattered-pellet injection that radiates the energy isotropically before
it reaches the wall. Runaway electrons form when the induced electric
field accelerates electrons past the critical field; a multi-MeV beam can
drill through a wall, and the control is the same impurity injection plus
runaway-dissipation coils. Disruption prediction and mitigation systems
are the integrating control.
\end{exercise}

\begin{exercise}
A Hall thruster's life-test shows electrode erosion at twice the
predicted rate. Identify the candidate causes (excessive ion
acceleration onto walls, plasma-wall reconnection events, propellant
contamination) and the working diagnostic test for each.
\paragraph{Solution sketch.} Excess ion acceleration onto the channel
walls shows as a broadened ion-energy distribution and a wider plume
divergence; measure the angular ion-current distribution with a movable
Faraday probe and the energy spectrum with a retarding-potential analyser.
Anomalous near-wall conduction shows as wall-current signatures and as
erosion patterns concentrated near the channel exit; inspect the erosion
profile and correlate with magnetic-field shaping. Propellant contamination
(impurities sputtering or chemically attacking electrodes) shows as
foreign lines in the plume spectrum and as deposits on witness plates;
run a residual-gas analysis on the feed and an optical-emission scan of
the plume. The distinguishing measurement is the plume optical-emission
spectrum, which separates a clean-but-misdirected beam from a contaminated
one.
\end{exercise}

\begin{exercise}
A plasma-etch tool produces wafer-edge non-uniformity that exceeds
process specifications. Identify the candidate causes (gas-injection
non-uniformity, magnetic-field non-uniformity, edge ring
contamination) and the working tests that distinguish them.
\paragraph{Solution sketch.} Map the etch rate radially first to see the
shape of the non-uniformity. Gas-injection non-uniformity gives a pattern
that tracks the showerhead hole layout and shifts with flow rate; vary the
flow and watch the map. Magnetic-field non-uniformity (in a magnetically
enhanced tool) gives an azimuthal pattern that rotates or shifts with the
coil currents; perturb the field and re-map. Edge-ring contamination or
wear gives a sharp roll-off in the outer few millimetres that worsens with
ring age; inspect and replace the edge ring and re-measure. The working
discriminator is which control variable, gas flow, coil current, or ring
condition, moves the non-uniformity pattern.
\end{exercise}

\subsection*{Judgment}\pathtag{standard}

\begin{exercise}
A consulting client must decide between deploying capacitively
coupled or inductively coupled plasma sources for a new etch tool.
Identify the trade-offs in terms of ion energy, plasma density,
working pressure range, and process uniformity.
\paragraph{Solution sketch.} Capacitively coupled plasmas (CCP) couple
power through the sheath, so ion energy and plasma density are linked
through a single knob; they run at higher pressure, give lower density
($\sim 10^{15}$ to $10^{16}\,\si{\per\cubic\meter}$), and are simple and
cheap, suited to processes that need high ion energy and modest rate.
Inductively coupled plasmas (ICP) couple power through a coil into the
bulk, decoupling density from ion energy: a separate bias controls the
sheath voltage while the coil sets density ($\sim 10^{17}$ to
$10^{18}\,\si{\per\cubic\meter}$). ICP runs at lower pressure, gives
higher rate and better independent control, at higher cost and complexity.
Choose CCP for simple high-energy etch, ICP for high-rate work needing
independent control of ion energy and flux.
\end{exercise}

\begin{exercise}
A research group must decide between continuing the ITER tokamak
programme and shifting investment to inertial-confinement fusion.
Identify the technical, economic, and schedule trade-offs and the
working criterion for the decision.
\paragraph{Solution sketch.} Magnetic confinement (tokamak) has a longer
record of sustained burning-plasma physics and a clearer path to
steady-state operation, but carries disruption risk and large-device cost.
Inertial confinement has demonstrated net target gain (NIF, 2022) but
faces pulse-repetition-rate, target-supply, and energy-conversion problems
that stand between a physics demonstration and a power plant. The honest
criterion is not which physics is more elegant but which approach reaches
a net-electricity demonstration at acceptable cost and schedule; that
turns on engineering integration (tritium breeding, materials, balance of
plant) more than on the confinement scheme. A sound recommendation funds
the integrating engineering on the more mature line while keeping a
hedging investment in the other.
\end{exercise}

\begin{exercise}
A regulator must decide whether to permit a private fusion-energy
start-up to operate a sub-megawatt deuterium-tritium tokamak in a
residential area. Identify the technical and regulatory issues
(tritium handling, neutron flux, decommissioning) that govern the
decision.
\paragraph{Solution sketch.} The binding issues are radiological, not the
plasma itself. Tritium is a mobile radioactive gas requiring containment,
inventory accounting, and effluent monitoring; even a sub-megawatt D-T
device produces a neutron flux that activates structural materials and
demands shielding and exclusion zones. Decommissioning leaves
neutron-activated components classed as low-level or intermediate waste.
A residential siting decision turns on whether the operator can bound the
tritium inventory and neutron dose at the site boundary below regulatory
limits, and on a funded decommissioning and waste plan. The defensible
regulatory criterion is dose to the public at the boundary plus credible
accident analysis, judged against the same standards applied to other
small radiological facilities.
\end{exercise}
